% concatenated from revised-JHEP.tex, custom.tex
\documentclass[11pt,a4paper]{article}
\usepackage{jheppub}
\usepackage{mathrsfs}
\usepackage{longtable}
\usepackage{graphicx}%
\usepackage{multirow}%
\usepackage{amsmath,amssymb,amsfonts}%
\usepackage{amsthm}%
\usepackage{mathrsfs}%
\usepackage[title]{appendix}%
\usepackage[dvipsnames,svgnames,x11names]{xcolor}%
\usepackage{textcomp}%
\usepackage{manyfoot}%
\usepackage{booktabs}%

\usepackage{listings}
\usepackage{comment}
\usepackage[mathscr]{eucal}
\usepackage{braket}
\usepackage{longtable}
\usepackage{array}
%\usepackage{romannum}
\usepackage{bm}
\usepackage{subfigure}
\usepackage[normalem]{ulem}
\usepackage{mathtools}
\usepackage{stmaryrd}
\makeatletter
\newcommand{\fixed@sra}{$\vrule height 2\fontdimen22\textfont2 width 0pt\shortrightarrow$}
\newcommand{\shortarrow}[1]{%
  \mathrel{\text{\rotatebox[origin=c]{\numexpr#1*45}{\fixed@sra}}}
}
\makeatother
\usepackage[utf8]{inputenc}
\usepackage[T1]{fontenc}
\DeclareMathAlphabet{\mathbbold}{U}{bbold}{m}{n}
\usepackage{hyperref}
\setcounter{tocdepth}{2}
\hypersetup{
    colorlinks=true,                 % false: boxed links; true: colored links
    linkcolor=blue,%SeaGreen,                 % color of internal links (change box color with linkbordercolor)
    linktoc=all,                     % link over complete toc entries
    citecolor=blue,                 % color of links to bibliography
    filecolor=blue,                  % color of file links
    urlcolor=blue                 % color of external links
}
\usepackage{url}
\usepackage{amsthm}
\theoremstyle{plain}
\newtheorem{theorem}{Theorem}
\numberwithin{theorem}{section}
\newtheorem{proposition}[theorem]{Proposition}
\theoremstyle{definition}
\newtheorem{definition}[theorem]{Definition}
\newtheorem{hypothesis}[theorem]{Hypothesis}
\newtheorem{remark}[theorem]{Remark}
\newtheorem{lemma}[theorem]{Lemma}
\newtheorem{assumption}[theorem]{Assumption}
\newtheorem{corollary}[theorem]{Corollary}
\newtheorem*{theorem*}{Theorem}
\newtheorem{conjecture}[theorem]{Conjecture}
%%%%%%%%%%%%%%%%%%%%%%%%	Custom definitions 


%%%%%%%%%%%%%%%%%%%%%%%%%%%%%%%%    REFERENCES    ---     COLORS & FONT SIZES         %%%%%%%%%%%%%%%%%%%%%%%%%%%%%%%%%%%%%%%%%%

%%%%%%%%  Font Sizes

%\tiny
%\scriptsize
%\footnotesize
%\small
%\normalsize
%\large
%\Large
%\LARGE
%\huge
%\Huge


%%%%%%%%%%%%%%%%%%%%%%%%%%%%%%%%%%%%%%%%%%%%%%%%%%%     COLOUR   DEFINITIONS         %%%%%%%%%%%%%%%%%%%%%%%%%%%%%%%%%%%%%%%%%%%%%%%%%%%%%%%%%%%%%%%%%%%%%

%%%%%%%%%%%%%%%%%%%%  Default colors

%black, blue, brown, cyan, darkgray, gray, green, lightgray, lime, magenta, olive, orange, pink, purple, red, teal, violet, white, yellow
%

%%%%%%%%%%%%%%%%%%%%%%%$$$%%%%%%        SHORTCUTS FOR LANDAUER              %%%%%%%%%%%%%%%%%%%%%%%%%%%%%%%%%%%%%%%%%%%%%%%%%%%%%%%%%%%%


\newcommand{\textdef}[1]{{\color{purple}\textsf{#1}}}

%%%%%%%%%%%%%%%%%%%%%%%$$$%%%%%%        SHORTCUTS FOR AQFT              %%%%%%%%%%%%%%%%%%%%%%%%%%%%%%%%%%%%%%%%%%%%%%%%%%%%%%%%%%%%

\newcommand{\vNA}{von Neumann algebra }
\newcommand{\vNsA}{von Neumann subalgebra }



%%%%%%%%%%%%%%%%%%%%%%%$$$%%%%%%        SHORTCUTS FOR MATHS              %%%%%%%%%%%%%%%%%%%%%%%%%%%%%%%%%%%%%%%%%%%%%%%%%%%%%%%%%%%%


\newcommand{\rmO}{\text{O}}
\newcommand{\rmU}{\text{U}}
\newcommand{\GL}{\text{GL}}
\newcommand{\SL}{\text{SL}}
\newcommand{\SU}{\text{SU}}
\newcommand{\SO}{\text{SO}}
\newcommand{\Sp}{\text{Sp}}


\newcommand{\overbar}[1]{\mkern 1.5mu\overline{\mkern-1.5mu#1\mkern-1.5mu}\mkern 1.5mu}


%%%%%%%%%%%%%%%%%%%%%%%%%%%


\newcommand{\df}{{\rm d}}
\newcommand{\del}{\partial}
\newcommand{\wt}[1]{\widetilde{#1}}
\newcommand{\wh}[1]{\widehat{#1}}
\newcommand{\?}{\overset{?}{=}}


\newcommand{\defeq}{\mathrel{\mathop:}=}
\newcommand{\eqdef}{=\mathrel{\mathop:}}


\DeclareMathOperator{\Sch}{Sch}
\DeclareMathOperator{\sech}{sech}
\DeclareMathOperator{\csch}{csch}
\DeclareMathOperator{\arccosh}{arccosh}
\DeclareMathOperator{\arcsech}{arcsech}
\DeclareMathOperator{\arcsinh}{arcsinh}
\DeclareMathOperator{\arccsch}{arccsch}
\DeclareMathOperator{\arctanh}{arctanh}
\DeclareMathOperator{\arccoth}{arccoth}

\DeclareMathOperator{\tr}{tr}
\DeclareMathOperator{\Tr}{Tr}
\DeclareMathOperator{\rank}{rank}
\DeclareMathOperator{\diag}{diag}
\DeclareMathOperator{\Span}{span}
\DeclareMathOperator{\spec}{spec}
\DeclareMathOperator{\ad}{ad}
\DeclareMathOperator{\Ad}{Ad}
\DeclareMathOperator{\im}{im}
\DeclareMathOperator{\sign}{sign}
\DeclareMathOperator{\supp}{supp}
\renewcommand{\Re}{\operatorname{Re}}
\renewcommand{\Im}{\operatorname{Im}}

% \newcommand{\ad}[1]{\,\text{ad}_{#1}}
% \newcommand{\Ad}[1]{\,\text{Ad}_{#1}}


\newcommand{\abs}[1]{\left\vert#1\right\vert}

\newcommand{\ohgf}[1]{\tensor[_2]{F}{_1}\lB #1 \rB}   %% Ordinary Hypergeometric fn.

\newcommand{\rmnum}[1]{\romannumeral #1}
\newcommand{\Rmnum}[1]{\expandafter\@slowromancap\romannumeral #1@}



%%%%%%%%%%%%%%%%%%%%%%%%%%%%%%%%        EQUATION SHORTCUTS          %%%%%%%%%%%%%%%%%%%%%%%%%%%%%%%%%%%%%%%%%

\newcommand{\nn}{\nonumber}
\newcommand{\dsp}{\displaystyle}


\newcommand{\be}{\begin{equation}}
\newcommand{\ee}{\end{equation}}


\newcommand{\bea}[1]{\begin{align}#1\end{align}}
\newcommand{\bes}[1]{\begin{equation}\begin{split}#1 \end{split}\end{equation}}


%%%%%%%%%%%%%%%%%%%%%%%%%%%%%%%%        BRACKETS        %%%%%%%%%%%%%%%%%%%%%%%%%%%%%%%%%%%%%%%%%%%

\newcommand{\lB}{\left [}
\newcommand{\rB}{\right ]}
\newcommand{\lb}{\left (}
\newcommand{\rb}{\right )}
\newcommand{\lbr}{\left \{}						
\newcommand{\rbr}{\right \}}					
\newcommand{\lv}{\left\vert}
\newcommand{\rv}{\right\vert}
\newcommand{\lla}{\left\langle}
\newcommand{\rra}{\right\rangle}
\newcommand{\ld}{\left.}				        %% Left  continuation of bracketed expression.
\newcommand{\rd}{\right.}				        %% Right  continuation of bracketed expression.
\newcommand{\bfrac}[2]{\lb\frac{#1}{#2}\rb}  	%% Fractions in brackets.



%%%%%%%%%%%%%%%%%%%%%%%%%%%%%%%%%%%%%%%    SYMBOL SHORTCUTS     %%%%%%%%%%%%%%%%%%%%%%%%%%%%%%%%%%%%%%%%%%%%%%%%%%%%

\newcommand{\A}{\,\forall\,}
\newcommand{\C}{\mathbb{C}}
\newcommand{\I}{\infty}
\newcommand{\M}{\mathbb{M}}
\newcommand{\N}{\mathbb{N}}
\newcommand{\Qu}{\mathbb{H}}
\newcommand{\R}{\mathbb{R}}
\newcommand{\Z}{\mathbb{Z}}

%%%%%  Wide tilde

\newcommand{\wtA}{\widetilde{A}}
\newcommand{\wtB}{\widetilde{B}}
\newcommand{\wtC}{\widetilde{C}}
\newcommand{\wtD}{\widetilde{D}}
\newcommand{\wtE}{\widetilde{E}}
\newcommand{\wtF}{\widetilde{F}}
\newcommand{\wtG}{\widetilde{G}}
\newcommand{\wtH}{\widetilde{H}}
\newcommand{\wtI}{\widetilde{I}}
\newcommand{\wtJ}{\widetilde{J}}
\newcommand{\wtK}{\widetilde{K}}
\newcommand{\wtL}{\widetilde{L}}
\newcommand{\wtM}{\widetilde{M}}
\newcommand{\wtN}{\widetilde{N}}
\newcommand{\wtO}{\widetilde{O}}
\newcommand{\wtP}{\widetilde{P}}
\newcommand{\wtQ}{\widetilde{Q}}
\newcommand{\wtR}{\widetilde{R}}
\newcommand{\wtS}{\widetilde{S}}
\newcommand{\wtT}{\widetilde{T}}
\newcommand{\wtU}{\widetilde{U}}
\newcommand{\wtV}{\widetilde{V}}
\newcommand{\wtW}{\widetilde{W}}
\newcommand{\wtX}{\widetilde{X}}
\newcommand{\wtY}{\widetilde{Y}}
\newcommand{\wtZ}{\widetilde{Z}}


%%%%%  Greek Alphabet
\renewcommand{\a}{\alpha}	%%% Redefinition
\renewcommand{\b}{\beta}		%%% Redefinition
\renewcommand{\c}{\chi}		%%% Redefinition
\renewcommand{\d}{\delta}	%%% Redefinition
\newcommand{\e}{\epsilon}
\newcommand{\ve}{\varepsilon}
\newcommand{\f}{\phi}
\newcommand{\vf}{\varphi}
\newcommand{\g}{\gamma}
\newcommand{\h}{\eta}
\newcommand{\hb}{\hbar}
\newcommand{\hs}{\hslash}
\renewcommand{\i}{\iota}		%%% Redefinition
\renewcommand{\k}{\kappa}	%%% Redefinition
\renewcommand{\l}{\lambda}	%%% Redefinition
\newcommand{\mn}{\mu\nu}
\newcommand{\q}{\theta}	
\newcommand{\vq}{\vartheta}	
\renewcommand{\r}{\rho}		%%% Redefinition
\newcommand{\vr}{\varrho}	%%% Redefinition
\newcommand{\s}{\sigma}
\newcommand{\vs}{\varsigma}
\renewcommand{\t}{\tau}		%%% Redefinition
\renewcommand{\u}{\upsilon}	%%% Redefinition
\newcommand{\w}{\omega}
\newcommand{\vw}{\varomega}
\newcommand{\y}{\psi}
\newcommand{\z}{\zeta}


\newcommand{\D}{\Delta}
\newcommand{\F}{\Phi}
\newcommand{\G}{\Gamma}
\renewcommand{\L}{\Lambda}	%%% Redefinition
\newcommand{\Q}{\Theta}
\renewcommand{\S}{\Sigma}	%%% Redefinition
\newcommand{\W}{\Omega}
\newcommand{\Y}{\Psi}



%%%%%   CALIGRAPHIC
\newcommand{\cA}{\mathcal{A}}
\newcommand{\cB}{\mathcal{B}}
\newcommand{\cC}{\mathcal{C}}
\newcommand{\cD}{\mathcal{D}}
\newcommand{\cE}{\mathcal{E}}
\newcommand{\cF}{\mathcal{F}}
\newcommand{\cG}{\mathcal{G}}
\newcommand{\cH}{\mathcal{H}}
\newcommand{\cI}{\mathcal{I}}
\newcommand{\cJ}{\mathcal{J}}
\newcommand{\cK}{\mathcal{K}}
\newcommand{\cL}{\mathcal{L}}
\newcommand{\cM}{\mathcal{M}}
\newcommand{\cN}{\mathcal{N}}
\newcommand{\cO}{\mathcal{O}}
\newcommand{\cP}{\mathcal{P}}
\newcommand{\cQ}{\mathcal{Q}}
\newcommand{\cR}{\mathcal{R}}
\newcommand{\cS}{\mathcal{S}}
\newcommand{\cT}{\mathcal{T}}
\newcommand{\cU}{\mathcal{U}}
\newcommand{\cV}{\mathcal{V}}
\newcommand{\cW}{\mathcal{W}}
\newcommand{\cX}{\mathcal{X}}
\newcommand{\cY}{\mathcal{Y}}
\newcommand{\cZ}{\mathcal{Z}}

%%%%%   BLACKBOARD
\newcommand{\bA}{\mathbb{A}}
\newcommand{\bB}{\mathbb{B}}
\newcommand{\bC}{\mathbb{C}}
\newcommand{\bD}{\mathbb{D}}
\newcommand{\bE}{\mathbb{E}}
\newcommand{\bF}{\mathbb{F}}
\newcommand{\bG}{\mathbb{G}}
\newcommand{\bH}{\mathbb{H}}
\newcommand{\bI}{\mathbb{I}}
\newcommand{\bJ}{\mathbb{J}}
\newcommand{\bK}{\mathbb{K}}
\newcommand{\bL}{\mathbb{L}}
\newcommand{\bM}{\mathbb{M}}
\newcommand{\bN}{\mathbb{N}}
\newcommand{\bO}{\mathbb{O}}
\newcommand{\bP}{\mathbb{P}}
\newcommand{\bQ}{\mathbb{Q}}
\newcommand{\bR}{\mathbb{R}}
\newcommand{\bS}{\mathbb{S}}
\newcommand{\bT}{\mathbb{T}}
\newcommand{\bU}{\mathbb{U}}
\newcommand{\bV}{\mathbb{V}}
\newcommand{\bW}{\mathbb{W}}
\newcommand{\bX}{\mathbb{X}}
\newcommand{\bY}{\mathbb{Y}}
\newcommand{\bZ}{\mathbb{Z}}

%%%%%   SCRIPT
\newcommand{\sA}{\mathscr{A}}
\newcommand{\sB}{\mathscr{B}}
\newcommand{\sC}{\mathscr{C}}
\newcommand{\sD}{\mathscr{D}}
\newcommand{\sE}{\mathscr{E}}
\newcommand{\sF}{\mathscr{F}}
\newcommand{\sG}{\mathscr{G}}
\newcommand{\sH}{\mathscr{H}}
\newcommand{\sI}{\mathscr{I}}
\newcommand{\sJ}{\mathscr{J}}
\newcommand{\sK}{\mathscr{K}}
\newcommand{\sL}{\mathscr{L}}
\newcommand{\sM}{\mathscr{M}}
\newcommand{\sN}{\mathscr{N}}
\newcommand{\sO}{\mathscr{O}}
\newcommand{\sP}{\mathscr{P}}
\newcommand{\sQ}{\mathscr{Q}}
\newcommand{\sR}{\mathscr{R}}
\newcommand{\sS}{\mathscr{S}}
\newcommand{\sT}{\mathscr{T}}
\newcommand{\sU}{\mathscr{U}}
\newcommand{\sV}{\mathscr{V}}
\newcommand{\sW}{\mathscr{W}}
\newcommand{\sX}{\mathscr{X}}
\newcommand{\sY}{\mathscr{Y}}
\newcommand{\sZ}{\mathscr{Z}}

%%%%%%     Mathfrak
\newcommand{\fA}{\mathfrak{A}}
\newcommand{\fB}{\mathfrak{B}}
\newcommand{\fC}{\mathfrak{C}}
\newcommand{\fD}{\mathfrak{D}}
\newcommand{\fE}{\mathfrak{E}}
\newcommand{\fF}{\mathfrak{F}}
\newcommand{\fG}{\mathfrak{G}}
\newcommand{\fH}{\mathfrak{H}}
\newcommand{\fI}{\mathfrak{I}}
\newcommand{\fJ}{\mathfrak{J}}
\newcommand{\fK}{\mathfrak{K}}
\newcommand{\fL}{\mathfrak{L}}
\newcommand{\fM}{\mathfrak{M}}
\newcommand{\fN}{\mathfrak{N}}
\newcommand{\fO}{\mathfrak{O}}
\newcommand{\fP}{\mathfrak{P}}
\newcommand{\fQ}{\mathfrak{Q}}
\newcommand{\fR}{\mathfrak{R}}
\newcommand{\fS}{\mathfrak{S}}
\newcommand{\fT}{\mathfrak{T}}
\newcommand{\fU}{\mathfrak{U}}
\newcommand{\fV}{\mathfrak{V}}
\newcommand{\fW}{\mathfrak{W}}
\newcommand{\fX}{\mathfrak{X}}
\newcommand{\fY}{\mathfrak{Y}}
\newcommand{\fZ}{\mathfrak{Z}}
\newcommand{\fa}{\mathfrak{a}}
\newcommand{\fb}{\mathfrak{b}}
\newcommand{\fc}{\mathfrak{c}}
\newcommand{\fd}{\mathfrak{d}}
\newcommand{\fe}{\mathfrak{e}}
\newcommand{\ff}{\mathfrak{f}}
\newcommand{\fg}{\mathfrak{g}}
\newcommand{\fh}{\mathfrak{h}}
\newcommand{\fii}{\mathfrak{i}}		%%% Out of Order
\newcommand{\fj}{\mathfrak{j}}
\newcommand{\fk}{\mathfrak{k}}
\newcommand{\fl}{\mathfrak{l}}
\newcommand{\fm}{\mathfrak{m}}
\newcommand{\fn}{\mathfrak{n}}
\newcommand{\fo}{\mathfrak{o}}
\newcommand{\fp}{\mathfrak{p}}
\newcommand{\fq}{\mathfrak{q}}
\newcommand{\fr}{\mathfrak{r}}
\newcommand{\fs}{\mathfrak{s}}
\newcommand{\ft}{\mathfrak{t}}
\newcommand{\fu}{\mathfrak{u}}
\newcommand{\fv}{\mathfrak{v}}
\newcommand{\fw}{\mathfrak{w}}
\newcommand{\fx}{\mathfrak{x}}
\newcommand{\fy}{\mathfrak{y}}
\newcommand{\fz}{\mathfrak{z}}

%%%%%%%%%%%%%%%%%%%%%%%%%%%%%%%%%     END - SYMBOL DEFINITIONS       %%%%%%%%%%%%%%%%%%%%%%%%%%%%%%%%%%%%


%%%%%%%%%%%%%%%%%%%%%%%%


\newcommand{\rms}{{\mathrm{s}}}
\newcommand{\Ahat}{\widehat{A}}
\newcommand{\Ahateps}{\widehat{A}^{(\epsilon)}}
\newcommand{\Mhat}{\widehat{M}}
\newcommand{\hhat}[1]{\hat{h}_{#1}}
\newcommand{\az}{$\a$-$z$}
\newcommand{\Sgen}{S_{\text{gen}}}
\newcommand{\barSgen}{\bar{S}_{\text{gen}}}


\DeclareMathOperator{\Dom}{Dom}

\title{A general proof of integer R\'{e}nyi QNEC}


\author[a]{Tanay Kibe}
\emailAdd{tanay.kibe@ib.edu.ar}
\affiliation[a]{Instituto Balseiro, Centro 
At{\'o}mico Bariloche, S.C. de Bariloche, 8400, 
R{\'i}o Negro, Argentina}
\author[b]{and Pratik Roy}
\emailAdd{roy.pratik92@gmail.com }
\affiliation[b]{Institute of Mathematics, University of Warsaw, ul. Banacha 2, 02-097 Warsaw, Poland}



\begin{document}


\abstract{The R\'{e}nyi quantum null energy condition conjectures that the second null shape variation of the sandwiched R\'{e}nyi divergence (SRD) of an excited state relative to the vacuum is non-negative in local Poincar\'e-invariant quantum field theory, giving a one-parameter generalization of the quantum null energy condition (QNEC). We prove R\'{e}nyi QNEC for all integer R\'{e}nyi parameters $n\geq 2$ for von Neumann algebras carrying a half-sided modular inclusion structure. The only additional regularity assumption on the excited state is finiteness of the SRD relative to the vacuum. Concretely, for any $\sigma$-finite von Neumann algebra with such an inclusion, we prove log-convexity, under the associated null-translation semigroup, of the Kosaki $L^n$ norm of any normal positive functional with finite $L^n$ norm. 
}
\maketitle


\section{Introduction}


The quantum null energy condition (QNEC) was first conjectured as a non-gravitational limit of the quantum focusing conjecture in semi-classical gravity \cite{Bousso:2015mna}, and has since been proven in free quantum field theories (QFT) \cite{Bousso:2015wca,Malik:2019dpg}, holographic QFTs with anti de-Sitter duals \cite{Koeller:2015qmn}, and general local Poincar\'{e}-invariant QFTs \cite{Balakrishnan:2017bjg}. Moreover, mathematically rigorous proofs of the QNEC have been provided using von Neumann algebraic techniques \cite{Ceyhan:2018zfg, Hollands:2025glm}. 

Ceyhan and Faulkner \cite{Ceyhan:2018zfg}, and later Hollands and Longo \cite{Hollands:2025glm}, make Wall's information theoretic argument for QNEC \cite{Wall:2017blw} mathematically rigorous for arbitrary QFTs that admit the structure of half-sided modular inclusions (HSMI) \cite{Borchers:1991xk,Wiesbrock:1992mg,Borchers:1995zg,Borchers1996,Araki:2005we}. Relative entropy in QFT is well-defined and UV-finite, and \cite{Ceyhan:2018zfg,Hollands:2025glm} rigorously prove the QNEC as the non-negativity of the second null derivative of the relative entropy:
\begin{equation}
    \partial^2_v S(v) \geq0\,,\quad S(v)\coloneqq D(\Psi\|\Omega;M_v)\,,
\end{equation}
where $D(\Psi\|\Omega;M_v)$ is the relative entropy of the excited vector state $\Psi$ with respect to the QFT vacuum vector $\Omega$, for a one-parameter family of local algebras $M_v$, with $M_w \subset M_v$ if $w>v$. These proofs need additional regularity assumptions on the state $\Psi$ and do not apply to arbitrary finite-relative-entropy states.

Our interest in QNEC stems, among other things, from the fact that it is not merely a formal consistency condition on QFTs, and can lead to non-trivial quantum thermodynamic bounds on quantities such as entropy production in dynamical processes \cite{Kibe:2021qjy,Banerjee:2022dgv,Kibe:2024icu,Kibe:2025cqc} (see also \cite{Mezei:2019sla}). QNEC has also been applied to prove renormalization group irreversibility theorems in the presence of defects \cite{Casini:2023kyj}.
It is therefore interesting to explore generalisations of QNEC and investigate the bounds that they can reveal. 
Generalisations of relative entropy that obey the data processing inequality (DPI) are natural candidates to study in this context.
In this paper we shall focus on one such quantity, the sandwiched R\'enyi divergence (SRD) \cite{Muller-Lennert:2013liu,Wilde:2013bdg}, which is a one-parameter generalization of the relative entropy.

 In finite-dimensional systems, the SRD between two density matrices $\rho,\sigma$ is \cite{Muller-Lennert:2013liu}
\begin{equation}\label{Eq:SRD-finite}
    {D}_{\alpha}(\rho\|\sigma) \coloneqq \frac{1}{\alpha-1}\log \frac{\Tr\lB\lb\sigma^{\frac{1-\alpha}{2\alpha}}\rho\sigma^{\frac{1-\alpha}{2\alpha}}\rb^\alpha\rB}{\Tr(\rho)}\,.
\end{equation}
For normalized density matrices, the SRD is a measure of distinguishability between the two states. As we will review, the SRD can be defined for general QFTs  in the von Neumann algebraic setting \cite{Berta:2016vnw,Jencova:2016tqz,Jencova:2017txf}, and has been shown to satisfy the properties that are expected of a distinguishability measure for normalized states: DPI for $\a\in [1/2, 1)\cup(1,\I)$ and invariance under local unitaries. In the $\alpha \to 1$ limit, the SRD becomes the standard relative entropy $D(\rho\|\sigma)$,
\begin{equation}
    \lim_{\alpha \to 1} {D}_{\alpha}(\rho\|\sigma)= D(\rho\|\sigma)\,.
\end{equation}

Lashkari \cite{Lashkari:2018nsl} conjectured the R\'{e}nyi QNEC (RQNEC): second null shape variations of SRD in a Poincar\'{e}-invariant QFT  are non-negative. Schematically:
\begin{equation}
  \partial^2S_\a(v) \geq0, \quad S_\a(v)= D_\a(\Psi\|\Omega; M_v)\,,
\end{equation}
where $D_\a(\Psi\|\Omega, M_v)$ is the SRD defined for the one-parameter family of von Neumann algebras $M_v$, $\Psi$ is any excited state, and $\Omega$ is the QFT vacuum. For $\a>1$, RQNEC has been proven in free field theories, and counter-examples have been constructed for $\a<1$ \cite{Moosa:2020jwt,Roy:2022yzm}. In this paper we prove the RQNEC for integer $\alpha\geq 2$.

\subsection{Main results}
The principal result of this paper is an operator-algebraic log-convexity
theorem for noncommutative \(L^n\)-norms associated with half-sided modular
inclusions.  Its main physical consequence is the R\'enyi quantum null energy
condition for integer R\'enyi index.  We state the operator-algebraic result
first, before explaining its relation to quantum field theory.

\subsubsection{Abstract theorem}
Let \(M_B\subset M_A=:M\) be a half-sided modular inclusion with respect to a
faithful normal state \(\omega\) on a \(\sigma\)-finite von Neumann algebra
\(M\).  The half-sided modular inclusion determines a one-parameter unitary group \(U_c=e^{-icP}\), and hence a semigroup of normal
unital endomorphisms of \(M\),
\begin{equation}\label{eq:intro-Tc}
        T_c(a):=U_{-c}aU_c,
        \qquad c\geq0,
\end{equation}
with \(T_c(M)\subset M\).  Since \(\omega\circ T_c=\omega\), the maps \(T_c\)
extend to positive contractive strongly continuous semigroups
\[
        V_c^{(p)}:L^p(M)\longrightarrow L^p(M),
        \qquad 1<p<\infty,
\]
on Haagerup's noncommutative \(L^p\)-spaces.

Our main theorem is the following. Its precise formulation appears as
Theorem~\ref{thm:general-log-convexity}. Below $\tr$ denotes Haagerup's canonical linear functional on the $L^1$ space. 

\paragraph*{Theorem A (Log-convexity under HSMI).}
Let \(n\geq2\) be an integer and let \(x_0\in L^n(M)_+\), $x_0\neq0$, and
\begin{equation}\label{eq:intro-Qn}
        Q_n(c)
        :=
        \tr\left(\bigl(V_c^{(n)}x_0\bigr)^n\right),
        \qquad c\geq0,
\end{equation}
Then $Q_n(c)$ is log-convex: for \(c_0,c_1\geq0\) and
\(\lambda\in[0,1]\),
\begin{equation}\label{eq:intro-log-convexity}
\begin{aligned}
        Q_n\bigl((1-\lambda)c_0+\lambda c_1\bigr)
        &\leq
        Q_n(c_0)^{1-\lambda}Q_n(c_1)^\lambda .
\end{aligned}
\end{equation}
In particular, wherever \(Q_n\) is twice differentiable and non-zero,
\begin{equation}
        Q_n(c)Q_n''(c)-Q_n'(c)^2\geq0.
\end{equation}
A key point is that Theorem~A is stated for an arbitrary
\(x_0\in L^n(M)_+\).  No differentiability assumption and no domain membership
condition involving the generator of \(V_c^{(n)}\) is imposed on \(x_0\). 


%%%%%%%%%%%%%%%%%%%%%%%%%%%%%%%%%%%%%%%%%%%%%%%%%%%%%%%%%%

\subsubsection{QFT setup and R\'enyi QNEC}\label{sec:setup}
The physical consequence of the abstract Theorem~A is the integer R\'enyi QNEC described below.  The precise statement appears as Corollary~\ref{cor:integer-renyi-qnec-fixed}. We emphasize that the stress tensor formulas that appear in this section motivate the HSMI abstraction, but are not used in the proof.

Consider a local Poincar{\'e}-invariant QFT living on $d$-dimensional Minkowski space, with metric
\begin{equation}
	ds^2 = -dx^+\,dx^- + d\vec y^2\,,
\end{equation}
where $x^\pm$ are null coordinates and $\vec y$ denotes the $d-2$ coordinates transverse to the lightcone. The Rindler wedge $R$ is the region $\{x^+>0,x^-<0\}$, and we denote the algebra of bounded operators associated with $R$ by $M_R$. The regions $N_C=\{x^-=0,x^+>C(\vec y)\}$, where $C(\vec y)$ is a continuous function of the transverse coordinates along the entangling surface, are called null cuts. We consider the von Neumann algebra, $M_C=M_{N_C''}$, of bounded operators localised in the causal completion $N_C''$.
The Rindler wedge corresponds to $M_0=M_R$. The algebras are illustrated in Fig.~\ref{fig:nullcuts}.
\begin{figure}
    \centering
    \includegraphics[width=0.5\linewidth]{nullcuts.png}
    \caption{Schematic of the null cut algebras. The transverse directions are perpendicular to the plane of the figure. The null planes $x^\pm=0$ are labelled. Light blue regions correspond to the Rindler wedge $R: \{x^+>0, x^-<0\}$ and its commutant $R'$. Dark blue regions correspond to the null cut algebra $M_b$ and its commutant $M_b'$.}
    \label{fig:nullcuts}
\end{figure}


Denote by $\W \in \mathcal H$ the vacuum vector in the QFT Hilbert space $\mathcal H$, which is cyclic and separating for all the local algebras. For the Rindler cut, the modular unitary group $\D_{\W;M_R}^{is}$ implements Lorentz boosts that fix the entangling surface $x^-=0$. Define the modular Hamiltonian associated with a cut $A$ by $K_A=-\log\D_{\W;M_A}$. Then, we have \cite{Casini:2017roe}
\begin{equation}
	K_A = 2\pi \int_{x^-=0} (x^+-A(\vec y))T_{++}\,.
\end{equation}
The modular Hamiltonians associated with two null cuts $A(\vec y),B(\vec y)$ have the commutation relation
\begin{equation}
	[K_A,K_B] = 2\pi i(K_A-K_B) = (2\pi)^2 i P\,,
\end{equation}
where $P$ is the self-adjoint and positive semi-definite null translation generator acting on Rindler horizons as translations by $B(\vec y)-A(\vec y)$. For null cuts, $P$ is the averaged null energy (ANE) operator
\begin{equation}
	P= \int_{x^-=0} (B(\vec y)-A(\vec y)) T_{++}\,.
\end{equation}

Given two null cuts $A,B$, with $B(\vec y)\geq A(\vec y)$, define the one-parameter family of cuts  
\[
    A_b=A+ b(B-A).
\]
We introduce the notation
\[
    M_{b}\defeq M_{A_b} = U_{-b}M_A U_b, \quad U_b= e^{-ibP},
\]
and the algebras associated with the null cuts $A$ and $B$ are denoted by $M_A,M_B$, respectively. In cases of interest, these cuts are expected to satisfy the properties of half-sided modular inclusions (HSMI) \cite{Borchers:1991xk,Wiesbrock:1992mg,Borchers:1995zg,Borchers1996,Araki:2005we}.

With this motivation, the physical consequence of Theorem~A is the following integer R\'enyi QNEC.  
\paragraph*{Corollary B (Integer R\'enyi QNEC).}
Suppose that the half-sided modular inclusion \(M_B\subset M_A\) is realized by null-cut
algebras in a Poincar\'e-invariant quantum field theory, with common faithful normal vacuum functional
\(\omega\).  Let
\[
        M_c:=T_c(M_A)=U_{-c}M_AU_c,
        \qquad c\geq0.
\]
For every integer \(n\geq2\) and every non-zero normal positive functional
\(\psi\) satisfying
\begin{equation}\label{assum:Assum1}
        D_n(\psi\|\omega;M_A)<\I,
\end{equation}
the function
\begin{equation}\label{eq:intro-RQNEC}
        c\longmapsto D_n(\psi\|\omega;M_c)
\end{equation}
is convex on \([0,\I)\).  Equivalently, at points at which the second
derivative exists,
\begin{equation}
        \partial_c^2D_n(\psi\|\omega;M_c)\geq0.
\end{equation}
We emphasize that the only assumption on $\psi$ is the finiteness of $D_n(\psi\|\omega,M_A)$ for some null cut $M_A$.%, which is the minimal condition under which the RQNEC is well defined.

\subsection{Proof sketch and outline of the paper}
Haagerup's noncommutative \(L^p\)-spaces play a central role in our proof.  They are
defined through the modular crossed product of the von Neumann algebra and
provide a representation-independent replacement for Schatten spaces.
The sandwiched R\'enyi divergence admits a natural formulation in these
spaces \cite{Jencova:2016tqz,Jencova:2017txf}. In Sec.~\ref{Sec:prelim} we discuss some preliminaries on Haagerup's non-commutative $L^p$ spaces and define the SRD. 

If
\(h_\psi,h_\omega\in L^1(M)_+\) are the Haagerup densities of two normal
positive functionals \(\psi\) and \(\omega\), then, for \(n>1\), finiteness
of \(D_n(\psi\|\omega)\) is equivalent to the existence of a unique
\(x_\psi\in L^n(M)_+\) such that
\begin{equation}\label{eq:intro-kosaki-representative}
        h_\omega^{\eta_n}
        x_\psi
        h_\omega^{\eta_n}
        =
        h_\psi,
        \qquad
        \eta_n:=\frac{n-1}{2n}.
\end{equation}
In this case,
\begin{equation}\label{eq:intro-SRD-Haagerup}
        D_n(\psi\|\omega;M)
        =
        \frac{1}{n-1}
        \log
        \frac{\tr(x_\psi^n)}{\tr(h_\psi)}.
\end{equation}

We now describe the main analytic steps of the proof.

\paragraph*{1. Fixed-algebra reformulation. (Sec.~\ref{Sec:HSMI})}
Following \cite{Hollands:2025glm}, the SRD of a fixed state on the moving algebra
\(M_c\) is first rewritten as the SRD of a moving state on the fixed
algebra \(M_A\):
\begin{equation}
        D_n(\psi\|\omega;M_c)
        =
        D_n(\psi_c\|\omega;M_A),
        \qquad
        \psi_c:=\psi\circ T_c.
\end{equation}
For vector states, this corresponds to replacing \(\Psi\) by
\(\Psi_c=U_c\Psi\).

\paragraph*{2. The \(L^n\)-semigroup and cyclic generator identity. (Sec.~\ref{sec:cyclic-ward})}
The maps \(T_c\) lift to a positive contractive strongly continuous
semigroup \(V_c^{(n)}\) on \(L^n(M_A)\), and
\begin{equation}
        x_{\psi_c}=V_c^{(n)}x_\psi.
\end{equation}
Let \(G_n\) denote the generator of \(V_c^{(n)}\).  Modular and null
smoothings are used to construct a graph core for \(G_n\) on which we establish the cyclic
generator identity
\begin{equation}\label{eq:intro-cyclic-Ward}
        \sum_{k=0}^{n-1}
        \zeta_n^k
        \tr\left(
        y_1\cdots y_k
        (G_ny_{k+1})
        y_{k+2}\cdots y_n
        \right)
        =
        0,
        \qquad
        \zeta_n=e^{-2\pi i/n}.
\end{equation}
We refer to this as the cyclic Ward identity.
The identity is then extended
by graph-norm density to arbitrary
\(y_1,\ldots,y_n\in\operatorname{Dom}(G_n)\).

\paragraph*{3. Smooth-orbit log-convexity. (Sec.~\ref{sec:integer-renyi-qnec})}
Assume temporarily that \(x_\psi\in\operatorname{Dom}(G_n^2)\), and set
\[
        x_c:=V_c^{(n)}x_\psi,
        \qquad
        Q_n(c):=\tr(x_c^n).
\]
Then \(x_c\) is twice differentiable in \(L^n(M)\), and the continuity of
the multilinear product
\[
        L^n(M)\times\cdots\times L^n(M)
        \longrightarrow
        L^1(M)
\]
gives
\begin{align}\label{eq:derivatives}
\begin{aligned}
        &Q_n'(c)=
        n\tr\left((G_nx_c)x_c^{n-1}\right),\\
        &Q_n''(c)=
        n\tr\left((G_n^2x_c)x_c^{n-1}\right)
        +
        n\sum_{j=0}^{n-2}
        \tr\left(
        (G_nx_c)x_c^j
        (G_nx_c)x_c^{n-2-j}
        \right).
\end{aligned}
\end{align}
Applying the cyclic identity \eqref{eq:intro-cyclic-Ward} to the first term in the double derivative $Q''_n(c)$ above, and using the Cauchy--Schwarz inequality in the Hilbert space \(L^2(M)\)
yields
\begin{equation}
        Q_n(c)Q_n''(c)-Q_n'(c)^2\geq0.
\end{equation}
Thus \(Q_n\) is log-convex for generator-smooth initial data.

\paragraph*{4. Removal of the regularity assumption. (Sec.~\ref{sec:integer-renyi-qnec})}
For a general \(x_\psi\in L^n(M)_+\), let
\begin{equation}
        R_m:=m(m-G_n)^{-1},
        \qquad
        x_\psi^{(m)}:=R_m^2x_\psi,
\end{equation}
where $R_m$ is the Yosida resolvent of the generator $G$. Using the fact that $x_\psi^{(m)}\in \Dom(G_n^2)$, the generator-smooth argument applies to 
\[
        Q_n^{(m)}(c)
        :=
        \tr\left(
        \bigl(V_c^{(n)}x_\psi^{(m)}\bigr)^n
        \right).
\]
We establish that \(Q_n^{(m)}(c)\to Q_n(c)\) pointwise, and passing to the limit in
the multiplicative convexity inequality proves Theorem~A without any
generator-domain assumption.


The logical dependence of the proof is as follows. Proposition~\ref{prop:null-translation-Ln-semigroup} constructs the $L^n$-semigroup, and Proposition~\ref{prop:equal-dual-semigroups} constructs its dual. 
Propositions~\ref{lem:smooth-core-delta} and~\ref{prop:ward-graph-core-fixed} provide a modular-entire graph core for the generator of the $L^n$-semigroup.
Proposition~\ref{prop:core-cyclic-ward-fixed} proves the cyclic identity on this core, and Theorem~\ref{thm:full-domain-ward-fixed} extends it to the full generator domain. 
Corollary~\ref{cor:cyclic-identity-logconvexity-form} is the only consequence of Sec.~\ref{sec:cyclic-ward} used in Proposition~\ref{prop:smooth-orbit-log-convexity-fixed}, which establishes smooth-orbit log-convexity. 
Finally, Theorem~\ref{thm:general-log-convexity} removes the generator-domain assumption by Yosida regularization, and Corollary~\ref{cor:integer-renyi-qnec-fixed} translates the resulting log-convexity into integer RQNEC.

Sec.~\ref{Sec:conclusion} discusses related open problems and future directions.  Appendix~\ref{appsec:technical} compiles standard definitions and results on integrals and one-parameter semigroups on Banach spaces. Appendix~\ref{Sec:app-haagerupLp} includes useful analytical results in Haagerup $L^p$ spaces. In Appendix~\ref{Sec:misc-results-Sec45}, we collect some technical results that are used in Sec.~\ref{sec:cyclic-ward} and \ref{sec:integer-renyi-qnec}. For the special case \(n=2\), \(L^2(M)\) is a Hilbert space, and we give a shorter proof of RQNEC using the spectral calculus of the positive generator
\(P\) in Appendix~\ref{subsec:alpha-two-rigorous}.


%%%%%%%%%%%%%%%%%%%%%%%%%%%%%%%%%%%%%%%%%%%%%%%%%%%%%%%%%%%%%%%

\section{Noncommutative \texorpdfstring{$L^p$}{Lp}-framework}
\label{Sec:prelim}

This section reviews notation and key facts about Haagerup's and Kosaki's non-commutative $L^p$ spaces, following~\cite{Terp:1981lp,Hiai:2021qfv,Hiai:2021cks}, and defines the sandwiched R\'enyi divergence (SRD) following~\cite{Jencova:2017txf, Jencova:2016tqz}. 
Some standard results on Haagerup $L^p$ spaces are collected in Appendix~\ref{Sec:app-haagerupLp}.

 

%%%%%%%%%%%%%%%%%%%%%%%%%%%%%%%%%%%%%%%%%%%%%%%%%%%%%%%%%%
\subsection{Haagerup \texorpdfstring{$L^p$}{Lp}-spaces}
\label{subsec:Haagerup-Lp}

Let $M$ be a $\sigma$-finite von Neumann algebra and $M_*$ be its predual  defined as the set of all $\sigma$-weakly continuous linear functionals on $M$. 
Let $\omega\in M_*^+$ be faithful and normal, and let the crossed-product with respect to the modular automorphism group $\sigma^{\omega}$ be $\mathcal R= M\rtimes_{\sigma^\omega} \mathbb{R} $. The algebra $\mathcal R$ is semi-finite and has a (semi-finite) trace $\tau$. Let $L^0(\mathcal R,\t)$ be the set of $\tau$-measurable operators affiliated with $\mathcal R$. 
\begin{definition}[Haagerup's $L^p$ spaces]\label{def:Haagerup-space}
For $0<p\leq \I$, Haagerup's $L^p$ space $L^p(M)$ is defined by
\begin{equation}
    L^p(M)\defeq\{x\in  L^0(\mathcal R,\t): \hat{\sigma}_s(x)=e^{-s/p}x,\, \forall s\in \mathbb{R}\}\,,
\end{equation}
where $\hat{\sigma}_s(\cdot)$ is the dual action of the modular automorphism.
\end{definition}
The construction is independent, up to canonical isometry, of the choice
of the faithful normal state \(\omega\).
$L^p(M)$ are closed linear subspaces of $ L^0(\mathcal R,\t)$ and are linearly spanned by positive operators: $L^p(M)_+ = L^p(M)\cap L^0(\mathcal R,\t)_+$. 

\begin{definition}[tr]
Define a linear functional ${\tr}$ on $L^1(M)$ by
\begin{equation}
    {\tr}(h_\psi)\defeq\psi(1),\qquad \psi\in M_*, \; h_\psi \in L^1(M)\,.
\end{equation}
\end{definition}
Note that
\begin{equation}
   {\tr}(|h_\psi|)=\tr(h_{|\psi|}) = |\psi|(1)=\|\psi\|_1 \quad\A\psi\in M_*\,.
\end{equation}
\begin{definition}[Haagerup norm: $\|a\|_p$]\label{def:haagerup-norm}
For every $a\in L^p(M)$, define $\|a\|_p \in [0,\I)$ as
\begin{align}
    \|a\|_p \defeq&\ {\tr}(|a|^p)^{1/p}, \quad p\in(0,\I)\,,\\
    \|a\|_\I \defeq&\ \|a\|\,.
\end{align}
For $p\geq1$, $\|\cdot\|_p$ is a norm; for $p\in(0,1)$ it is a quasi-norm
\end{definition}


\begin{proposition}[Standard Haagerup \(L^p\)-facts \cite{Terp:1981lp}]
\label{prop:standard-Haagerup-facts}
Let $p\geq1$ and $q\leq\I$ with $1/p+1/q=1$.
    \begin{enumerate}
        \item If $a\in L^p(M)$ and $b\in L^q(M)$, then $ab,ba\in L^1(M)$ and
            \[
                \tr(ab) = \tr(ba)\,.
            \]
        \item  The map
            \begin{equation}
                M_*\ni\psi\mapsto h_\psi\in L^1(M)
            \end{equation}
                is an isometric linear bijection of $M_*$ onto $L^1(M)$.
        \item Let $a\in L^q(M)$ such that $a\geq 0$, then
                \[
                    \tr(ab) \geq0 ,\quad \forall\, b\in L^p(M)_+\,.
                \]
    \end{enumerate}
\end{proposition}
The following generalized H{\"o}lder inequality is from \cite[Theorem 9.17]{Hiai:2021cks}.
\begin{lemma}[Generalized H{\"o}lder inequality]\label{lem:general-holder}
Let $p,q,r\in(0,\I]$ with $1/r=1/p+1/q$. If $a\in L^p(M)$ and $b\in L^q(M)$, then $ab\in L^r(M)$ and
\[
    \|ab\|_r \leq \|a\|_p\|b\|_q\,.
\]
\end{lemma}
The following properties of Haagerup's $L^p$ spaces from \cite{Terp:1981lp} are useful for our purposes.
\begin{theorem}
    \begin{enumerate}
        \item $(L^p(M), \|\cdot\|_p)$ is a Banach space for all $p\in [1,\I]$.
        \item $L^2(M)$ is a Hilbert space with respect to the inner product: $\braket{a,b}_{L^2(M)} \defeq \tr(a^*b)= \tr(ba^*)$ for $a,b\in L^2(M)$.
        \item Let $1\leq p<\I$ and $1/p+1/q=1$. Then the dual Banach space of $L^p(M)$ is $L^q(M)$ with the duality pairing: $L^p(M) \times L^q(M)\ni(a,b) \mapsto \tr(ab) \in \mathbb C$.
    \end{enumerate}
\end{theorem}

We will also need the following result.
\begin{lemma}[\cite{KOSAKI198429}]\label{lem:modularflow-haagerup}
    For $\omega \in M_*^+$ faithful, and the modular automorphism group $\sigma^\omega_t$ with respect to $\omega$, we have
    \begin{equation}
        \sigma^\omega_t(x) = h_\omega^{it}x h_\omega^{-it}, \quad x\in M, \quad t\in \mathbb{R}\,.
    \end{equation}
\end{lemma}

The algebra $M$ can be densely embedded into Haagerup $L^p$ spaces as follows.
\begin{definition}[Dense embedding in $L^p(M)$ \cite{KOSAKI198429}]\label{def:symmetricembedding}
    Let \(1\leq p<\infty\), and $\omega$ be a distinguished faithful normal state and $h_\omega\in L^1(M)$ be its Haagerup density. Define the map
    \[
        j_p:M\longrightarrow L^p(M),
        \qquad
        j_p(a)
        :=
        h_\omega^{1/(2p)}a h_\omega^{1/(2p)}.
    \]
    Then $j_p(M)$ is dense in $L^p(M)$.
\end{definition}

Let $M^{\rm an} \subset M$ denote the elements of $M$ that are entire analytic with respect to the modular flow of the faithful normal state $\omega\in M_*^+$. For $a\in M^{\rm an}$ the map $t\to \sigma_t^{\omega}(a)$ can be extended to an entire analytic function $z\to\sigma_z^{\omega}(a)$ on the complex plane.  $M^{\rm an}$ is strong-$*$ dense in $M$ (see the construction of $b_\ve$ in \eqref{eq:b-eps-def}). Let $n$ be an integer. The following KMS placement formula for a product of $n$ modular entire elements will be extremely useful in the subsequent section.
\begin{lemma}[KMS placement formula]
\label{lem:KMS-placement}
Let $\omega\in M_*^+$ be faithful with Haagerup density $h_\omega\in L^1(M)_+$, and let $n$ be an integer \(n\geq2\).  Suppose that
\(a_1,\ldots,a_n\in M^{\rm an}\). Then
\begin{equation}
\label{eq:KMS-placement}
\tr\!\left(j_n(a_1)\cdots j_n(a_n)\right)
=
\omega\!\left(
a_n
\sigma_{-i/n}^\omega(a_1)
\sigma_{-2i/n}^\omega(a_2)
\cdots
\sigma_{-(n-1)i/n}^\omega(a_{n-1})
\right).
\end{equation}
\end{lemma}

\begin{proof}
Since
\[
        j_n(a)=h_\omega^{1/(2n)}ah_\omega^{1/(2n)},
\]
we have
\[
\begin{aligned}
\tr\!\left(j_n(a_1)\cdots j_n(a_n)\right)
&=
\tr\!\left(
h_\omega^{1/(2n)}
a_1h_\omega^{1/n}a_2
\cdots
h_\omega^{1/n}a_n
h_\omega^{1/(2n)}
\right)                                                   \\
&=
\tr\!\left(
a_nh_\omega^{1/n}a_1h_\omega^{1/n}a_2
\cdots
h_\omega^{1/n}a_{n-1}h_\omega^{1/n}
\right),
\end{aligned}
\]
where the second equality follows from the H\"older-compatible cyclicity of
the Haagerup functional $\tr$, Lemma~\ref{lem:holder-products-trace-cyclicity}.

For every modular-entire \(a\in M^{\rm an}\) and \(s\geq0\), the modular-flow
from Lemma~\ref{lem:modularflow-haagerup} gives
\begin{equation}
\label{eq:move-h-past-a}
        h_\omega^s a
        =
        \sigma_{-is}^\omega(a)h_\omega^s .
\end{equation}
Applying \eqref{eq:move-h-past-a} successively yields
\[
\begin{aligned}
&a_nh_\omega^{1/n}a_1h_\omega^{1/n}a_2\cdots
h_\omega^{1/n}a_{n-1}h_\omega^{1/n}                                  \\
&\qquad =
a_n
\sigma_{-i/n}^\omega(a_1)
\sigma_{-2i/n}^\omega(a_2)
\cdots
\sigma_{-(n-1)i/n}^\omega(a_{n-1})
h_\omega .
\end{aligned}
\]
All products above are H\"older-compatible and belong to \(L^1(M)\).
A final application of $\tr$ cyclicity, together with
\(\omega(b)=\tr(h_\omega b)\) for \(b\in M\), gives
\[
\begin{aligned}
\tr\!\left(j_n(a_1)\cdots j_n(a_n)\right)
&=
\tr\!\left(
h_\omega\,a_n
\sigma_{-i/n}^\omega(a_1)
\cdots
\sigma_{-(n-1)i/n}^\omega(a_{n-1})
\right)                                                   \\
&=
\omega\!\left(
a_n
\sigma_{-i/n}^\omega(a_1)
\cdots
\sigma_{-(n-1)i/n}^\omega(a_{n-1})
\right).
\end{aligned}   
\]
\end{proof}


\subsection*{Kosaki \texorpdfstring{$L^p$}{Lp} spaces}
From Definition~\ref{def:symmetricembedding}, $M$ is densely embedded into $L^1(M)$ as
\begin{equation}
    x\in M \mapsto h_\omega^{1/2} xh^{1/2}_\omega \in L^1(M)\,.
\end{equation}
Define the norm $\|h_\omega^{1/2} xh^{1/2}_\omega\|\defeq\|x\| $ on $h_\omega^{1/2} Mh^{1/2}_\omega\subset L^1(M)$, so that $h_\omega^{1/2} Mh^{1/2}_\omega \simeq M$. Then $L^1(M)$ and $h_\omega^{1/2} Mh^{1/2}_\omega$ are a pair of compatible Banach spaces.
\begin{definition}[Kosaki's $L^p$ spaces \cite{KOSAKI198429}]
    Let $1<p<\I$. Kosaki's symmetric $L^p$ space $L^p(M,\omega)$ is defined to be the complex interpolation space
    \begin{equation}
        C_{1/p}(h_\omega^{1/2} M h_\omega^{1/2},L^1(M))\,,
    \end{equation}
    with the complex interpolation norm $\|\cdot\|_{p,\omega}$. $L^1(M,\omega)=L^1(M)$ and $L^\I(M,\omega)=h_\omega^{1/2} M h_\omega^{1/2}$, with $\|h_\omega^{1/2} x h_\omega^{1/2}\|_{\I,\omega}=\|x\|$.
\end{definition}

\begin{theorem}[\cite{KOSAKI198429}]\label{thm:haagerup-kosaki}
    Let $p\geq 1$ and $q \leq\I$ with $\frac1p+\frac1q=1$, then
    \begin{align}
        &L^p(M,\omega)=h_\omega^{1/2q} L^p(M) h_\omega^{1/2q} \subset L^1(M)\,,\\
        &\|h_\omega^{1/2q} a h_\omega^{1/2q}\|_{p,\omega} = \|a\|_p, \quad a\in L^p(M)\,.
    \end{align}
    That is
    \begin{equation}
        L^p(M,\omega) \simeq L^p(M)
    \end{equation}
    via the isometry
    \begin{equation}
        a\mapsto h_\omega^{1/2q} a h_\omega^{1/2q}\,.
    \end{equation}
\end{theorem}

\subsection{Sandwiched R\'{e}nyi divergence}\label{ssec:srd}
Here, we define the SRD for the range of interest $\a\in (1,\I)$ for $\psi,\omega\in M_*^+$, with $\psi\neq0$, and $\omega$ faithful. We note, however, that the SRD can be defined without assuming faithfulness, and for the range $\alpha \in [1/2,1)\cup(1,\I)$.
\begin{definition}[SRD]\label{def:SRD}
Let $\alpha \in (1,\I)$. The SRD is defined as \cite{Jencova:2016tqz,Jencova:2017txf}
\begin{equation}
    D_\alpha(\psi\|\omega;M)\defeq\frac{1}{\alpha-1}\log\frac{Q_\alpha(\psi\|\omega;M)}{\psi(1)}\,,
\end{equation}
with
\begin{equation}\label{eq:QSRD}
 Q_\alpha(\psi\|\omega;M)\defeq
    \begin{cases}
        \|h_\psi\|_{\alpha,\omega}^\alpha, & h_\psi \in L^\alpha(M,\omega)\subset L^1(M),\\[1em]
        \I, \qquad \qquad &\text{otherwise}\,.
    \end{cases}
\end{equation}
\end{definition}

The following form of the SRD (from \cite[Lemma 8]{Kato:2023aro}) will be useful.
\begin{lemma}           \label{lem:y-op-defn-SRD}
Let $\alpha\in (1,\I)$, with $\psi,\omega \in M_*^+$, and $\omega$ faithful. The following are equivalent
\begin{enumerate}
    \item $Q_\a(\psi\|\omega;M)<\I$
    \item $h_\psi\in L^\alpha(M,\omega)\subset L^1(M)$
    \item There exists a unique $x\in L^\a(M)_+$ satisfying
    \begin{equation}			\label{eq:x-op-defn-SRD}
	L^\alpha(M,\omega)\ni h_\psi = h_\omega^{\frac{\a-1}{2\a}} x h_\omega^{\frac{\a-1}{2\a}}\,,
\end{equation} 
\end{enumerate}
In this case
 \begin{equation}
     Q_\a(\psi\|\omega;M)=\|h_\psi\|_{\alpha,\omega}^\alpha=\left\|x\right\|_\a^\a\,.
 \end{equation}
\end{lemma}

The SRD obeys the data processing inequality.
\begin{theorem}[Data Processing Inequality (DPI) \cite{Berta:2016vnw,Jencova:2016tqz,Jencova:2017txf}]\label{thm:DPI-range-SRD}
Let $\g:N\to M$ be a normal positive unital map between von Neumann algebras. Then we have
\begin{equation}
    D_{\a}(\psi\circ\g\|\omega\circ\g;N) \leq D_{\a}(\psi\|\omega;M)\,.
\end{equation}
\end{theorem}


\subsection{Extension of maps to \texorpdfstring{$L^p$}{Lp} spaces}\label{ssec:extension}

Maps between two von Neumann algebras $M,N$ can be lifted to bounded maps on their crossed products, and also to their Haagerup $L^p$ spaces for $p\geq1$. For the lift to crossed product algebras, one can use \cite[Theorem 4.1]{Haagerup:2008arm}. We will only require the extension theorem for $L^p$ spaces for maps $T: M \to M$, which we state below.

\begin{theorem}[Extension of maps to $L^p$ spaces]\label{thm:extension-Lp}
Let \(M\) be a \(\sigma\)-finite von Neumann algebra, let
\(\omega\in M_*^+\) be faithful and normal, and let
\(T:M\to M\) be a normal positive unital map satisfying
\[
        \omega\circ T=\omega.
\]
For \(1\leq p<\infty\), define the map on the dense subspace
\[
    T^{(p)}: j_p(M)\longmapsto j_p(T(M)).
\]
Then $T^{(p)}$ extends uniquely to a positive contraction
\[
        T^{(p)}:L^p(M)\longrightarrow L^p(M)\,.
\]
\end{theorem}
%%%%%%%%%%%%%%%%%%%%%%%%%%%%%%%%%%%%%%%%%%%%%%%%%%%%%%%%%%%%%%%%%%%%

\section{Half-sided modular inclusions and R\'{e}nyi QNEC}\label{Sec:HSMI}

The null variations appearing in the statement of QNEC and RQNEC can be rigorously defined for algebras carrying a half-sided modular inclusion (HSMI) structure.  In this section we recall the HSMI structure, collect the properties needed in later sections, and reformulate RQNEC in terms of a fixed algebra and a one-parameter family of null-deformed states, following~\cite{Hollands:2025glm}.  Our conventions for modular operators and relative modular operators are the same as \cite{Hiai:2021qfv}.
 
\subsection{Half-sided modular inclusions and null translations}

From this point onward, we work entirely in the following abstract operator-algebraic setting: a pair of von Neumann algebras $M_B\subset M_A$ with a common cyclic separating vector $\W$, or equivalently a common faithful normal state $\omega$, satisfying the half-sided modular inclusion condition in Definition \ref{def:HSMI}. All subsequent results depend only on this HSMI structure and its consequences collected in Theorem \ref{thm:modulartranslations}.


\begin{definition}[Half-sided modular inclusion]
\label{def:HSMI}
Let $M_B \subset M_A$. This inclusion is a half-sided modular inclusion if there is a faithful normal state $\omega$, common to both algebras, such that:
\begin{equation}
\Delta_{\omega;A}^{-is} M_B \Delta_{\omega;A}^{is} \subset M_B \qquad s \geq 0\,.
\end{equation}
\end{definition}

\begin{theorem}[Null translations]\label{thm:modulartranslations}
For an HSMI $M_B\subset M_A$ with common faithful normal state $\omega$, there is a family of unitary operators, which we refer to as null translations,
\[
    U_b\defeq e^{-i b P},\quad b\in \mathbb R\,,
\]
such that 
\[
P\defeq \frac{1}{2\pi} \overline{\left( \ln \Delta_{\omega;B} - \ln \Delta_{\omega;A} \right)}
\]
is self-adjoint and positive semi-definite \cite{Borchers:1991xk}, where the bar denotes the closure.
Furthermore, the following are true \cite{Borchers:1991xk,Wiesbrock:1992mg,Borchers1996,Araki:2005we,Borchers:1995zg}.
\begin{enumerate}
\item $\omega \circ {\rm Ad}(U_{-b})=\omega$ for all $b\in \mathbb R$.
\item The null translations $U_b$ generate the one-parameter family of sub-algebras:
\[
U_{-b} M_{A} U_{b} = M_{b}\,,
\]
with $M_{c}\subset M_{b}$ if $c\geq b$. Moreover $M_{{1}} =M_{B}$. This family of algebras has $\omega$ as a faithful normal state, and $M_b \subset M_A$ is a HSMI relative to $\omega$. The translations apply for all real $b$ and also apply for the commutants, which satisfy:
\[
 M_{b}' \subset M_{c}', \quad b \leq c\,.
\]
\item The modular automorphism of the algebra $M_A$ acts as
\[
\Delta_{\omega;A}^{is} M_{b} \Delta_{\omega;A}^{-is} = M_{b e^{-2\pi s}}\,,
\]
and similarly for the commutant.
\item The modular operators satisfy the Borchers covariance relation with the null translations
\[
    \Delta_{\omega;M_A}^{is} U_b \Delta_{\omega;M_A}^{-is} =   U_{e^{-2\pi s}b}  \,.
\]
Furthermore, the following properties hold:
\begin{align*}
 &J_{\omega;M_A} U_b J_{\omega;M_A}  =U_{-b} \\
  &\Delta^{is} _{\omega;M_{b_1}} \Delta_{\omega;M_{b_2}}^{-is} = U_{ (b_1-b_2) (e^{-2\pi s} -1)} \nonumber,
\end{align*}
where $J_{\omega;M_A}$ is the modular conjugation with respect to $\omega$.
\item $U_b$ is strongly continuous and can be analytically continued to complex $b$, where it is bounded by $1$ and strongly continuous for ${\rm Im}\, b \leq 0$.
\end{enumerate}
\end{theorem}

\subsection{Fixed algebra formulation of R\'{e}nyi QNEC}

For convenience, we introduce the notation
\begin{equation}			\label{eq:entropy-shorthand-1}
	S_\a(c)\defeq D_\a(\psi\|\omega;M_c)\,, \quad \psi, \omega\in (M_A)_*^+,
\end{equation}
with $\omega$ faithful.
Note that $S_\a(c)$ in \eqref{eq:entropy-shorthand-1} is the SRD for the fixed functionals $\psi,\omega \in  (M_A)_*^+$ on a one-parameter family of sub-algebras $M_c$. Following \cite{Hollands:2025glm}, we instead write $S_{\a}(c)$ as the SRD of a one-parameter family of functionals $\psi_c$ on the fixed algebra $M_A$. 
\begin{remark}
In this subsection, we view $\Ad(U_{-c}):M_A\to M_c$ as a unital $*$-isomorphism onto $M_c$. Beginning in Section \ref{sec:cyclic-ward} we will think of $T_c:M_A\to M_A$ defined in \eqref{eq:def-Tc} as a $*$-endomorphism of $M_A$ with range $M_c\subset M_A$.
\end{remark}

\begin{lemma}[SRD in terms of varying functionals on a fixed algebra]\label{lem:fixedalgebraSRD}
Under the hypotheses of Theorem~\ref{thm:modulartranslations}, for $\psi \in (M_A)_*^+$, $\psi\neq0$, and $\alpha\in (1,\I)$:
\begin{equation}
    S_\a(c)\defeq D_\a(\psi\|\omega;M_c)= D_\a(\psi_c\|\omega;M_A)\,,\qquad \psi_c\defeq \psi \circ {\rm Ad}(U_{-c})\in (M_A)_*^+.
\end{equation}
\end{lemma}
\begin{proof}
The map ${\rm Ad}(U_{-c}): M_A \mapsto M_c$ is a normal unital $*$-isomorphism and its inverse is also normal and unital. Applying the DPI from Theorem~\ref{thm:DPI-range-SRD} to this $*$-isomorphism and its inverse, combined with the invariance of $\omega$ under ${\rm Ad}(U_{-c})$ gives the result.
\end{proof}


We can now state the R\'{e}nyi QNEC conjecture precisely as follows.
\begin{conjecture}[R\'{e}nyi QNEC conjecture]\label{conj:RQNEC}
Under the hypotheses of Theorem~\ref{thm:modulartranslations}, for $\alpha\in (1,\I)$, with $\omega\in(M_A)_*^+$ faithful, and any non-zero $\psi\in (M_A)_*^+$ such that $D_\a(\psi\|\omega;M_A)<\I$, the function
\[
c\mapsto S_\a(c)= D_\a(\psi_c\|\omega;M_A)
\]
is convex.
\end{conjecture}
Counter-examples to the RQNEC have been found for $\a\in [1/2,1)$ in the context of null quantization in free field theories \cite{Moosa:2020jwt}, so the conjecture as stated is restricted to $\alpha>1$. The remainder of this paper proves the conjecture for integer $\a\geq2$.

%%%%%%%%%%%%%%%%%%%%%%%%%%%%%%%%%%%%%%%%%%%%%%%%%%%%%%%%%%%%%%%%%%%%%%%%%%%%%%%%%%%%%%%%%%%

\section{\texorpdfstring{$L^n$}{Ln} semigroup and the cyclic Ward identity}
\label{sec:cyclic-ward}

Fix an integer \(n\geq2\), and set
\[
        q_n:=\frac{n}{n-1},
        \qquad \eta_n\defeq \frac{n-1}{2n}=\frac{1}{2q_n}, \qquad
        \zeta_n:=e^{-2\pi i/n}.
\]
The purpose of this section is to prove the following theorem, which is Theorem~\ref{thm:full-domain-ward-fixed} below.
\begin{theorem*}
Let $(V_c^{(n)})_{c\geq0}$ be the Haagerup
\(L^n(M)\)-semigroup associated with the half-sided modular inclusion,
and let \(G_n\) be its generator.
For all $y_1,\ldots,y_n\in\Dom(G_n)$,
\begin{equation}\label{eq:full-domain-ward}
    \sum_{k=0}^{n-1}\zeta_n^k
  \tr\bigl(
    y_1\ldots y_k (G_ny_{k+1}) y_{k+2}\ldots y_n
  \bigr)=0\,.
\end{equation}
\end{theorem*}
The proof consists of four steps:
\begin{enumerate}
\item construct \(V_c^{(n)}\) and the dual semigroup \(V_c^{(n)}\) (Sec.~\ref{subsec:test-and-dual-semigroups});
\item construct semigroup and modular regularizations (Sec.~\ref{subsec:modular-semigroup-reg});
\item use the regularizations to construct a graph core for the generator of \(V_c^{(n)}\) (Sec.~\ref{subsec:ward-graph-core});
\item establish \eqref{eq:full-domain-ward} on the core and extend
it by graph-norm continuity (Sec~\ref{subsec:cyclic-ward-identity-fixed}).
\end{enumerate}

\subsection*{Assumptions and notation}
Throughout this section, we assume the HSMI $M_B\subset M_A$, with common faithful normal state $\omega$. We use the following shorthand:
\begin{equation}
  M\defeq M_A,\qquad h\defeq h_\omega,\qquad \sigma_t\defeq \sigma_t^\omega\,,
\end{equation}
where $h_\omega\in L^1(M)_+$ is the Haagerup density of $\omega$ and $\sigma_t^\omega$ is the modular automorphism.

Let
\begin{equation}\label{eq:def-Tc}
  T_c: M_A\mapsto M_c \subset M_A, \quad T_c(a)\defeq U_{-c}aU_c,\qquad c\geq 0\,, 
\end{equation}
be the null-translation endomorphism on $M_A$ obtained from the HSMI structure in Theorem~\ref{thm:modulartranslations}. 
We use the following consequences of Theorem~\ref{thm:modulartranslations}.:
\begin{enumerate}
    \item $T_c$ is a normal unital $*$-endomorphism of $M$, with
            \(
                T_{c+d}=T_cT_d,
            \)
    \item \(\omega\circ T_c=\omega\)
    \item  \(T_c(a)\to a\) in the strong-\(*\) topology as \(c\downarrow0\)
    \item The Borchers covariance relation (Theorem~\ref{thm:modulartranslations}(4)) implies for real $s$ and $c\geq 0$:
            \[
                 T_c(\sigma_s(a))=\sigma_s(T_{ce^{2\pi s}}(a)),\qquad a\in M\,.
            \]
\end{enumerate}

Using Definition~\ref{def:symmetricembedding} we define the two dense embeddings for the Banach space $L^n(M)$ and its dual $L^{q_n}(M)$ as:
\begin{equation}\label{eq:jn-symmetric-embedding}
  j_n(a)\defeq h^{1/(2n)}ah^{1/(2n)}\in L^n(M)\,,
\end{equation}
with dense range in $L^n(M)$, and
\begin{equation}  j_{q_n}(a)\defeq h^{\eta_n}ah^{\eta_n}\in L^{q_n}(M)\,,
\end{equation}
with dense range in $L^{q_n}(M)$.

We will also use several technical results collected in Appendices~\ref{appsec:technical}, \ref{Sec:app-haagerupLp}, and \ref{Sec:misc-results-Sec45}.
%%%%%%%%%%%%%%%%%%%%%%%%%%%%%%%%%%%%%%%%%%%%%%%%%%%%%%%%%%%%%%%%%%%%%%%%%%%%%

\subsection{Null translations as \texorpdfstring{$L^n(M)$}{Ln(M)} semigroups}
\label{subsec:test-and-dual-semigroups}

The goal of this section is to define the null translation semi-group and its dual semigroup. We start by defining it first on the dense subspace $j_n(M)\subset L^{n}(M)$. 
\begin{definition}[Null translations on the dense core]
\label{def:Vc-def-dense}
Fix \(c\geq0\).  Since \(T_c(M)\subset M\), define
\[
        \widetilde V_c^{(n)}:j_n(M)\longrightarrow j_n(M)
\]
by
\begin{equation}
\label{eq:Vc-defn}
        \widetilde V_c^{(n)}j_n(a)
        :=
        j_n(T_c(a)),
        \qquad a\in M.
\end{equation}
\end{definition}

\begin{proposition}[The null-translation semigroup on \(L^n(M)\)]
\label{prop:null-translation-Ln-semigroup}
For every \(c\geq0\), the map
\(\widetilde V_c^{(n)}\) extends uniquely to a positive contraction
\[
        V_c^{(n)}:L^n(M)\longrightarrow L^n(M).
\]
The family \(V^{(n)}=(V_c^{(n)})_{c\geq0}\) is a strongly continuous
contraction semigroup.  More precisely,
\begin{align}
        &V_{c+d}^{(n)}
        =
        V_c^{(n)}V_d^{(n)},
        &&c,d\geq0,                                      \label{eq:V-semigroup}\\
        &V_0^{(n)}
        =
        I,                                                \label{eq:V-identity}\\
        &\lim_{c\to c_0}
        \|V_c^{(n)}x-V_{c_0}^{(n)}x\|_n
        =
        0,
        &&x\in L^n(M),\quad c_0\geq0.                    \label{eq:V-strong-cont}
\end{align}
\end{proposition}

\begin{proof}
Fix \(c\geq0\).  The map \(T_c:M\to M\) is normal, unital, positive,
and \(\omega\)-preserving.  Therefore
Theorem~\ref{thm:extension-Lp} applies and shows that $\widetilde V_c^{(n)}$, from Definition~\ref{def:Vc-def-dense},
extends uniquely to a positive contraction
\[
        V_c^{(n)}:L^n(M)\longrightarrow L^n(M).
\]
In particular,
\begin{equation}
\label{eq:V-contraction}
        \|V_c^{(n)}\|_{L^n(M)\to L^n(M)}
        \leq1,
        \qquad c\geq0.
\end{equation}

We next verify the semigroup property.  For \(a\in M\) and \(c,d\geq0\),
\begin{align*}
        V_c^{(n)}V_d^{(n)}j_n(a)=
        j_n(T_c(T_d(a)))=
        j_n(T_{c+d}(a))=
        V_{c+d}^{(n)}j_n(a).
\end{align*}
Thus
\[
        V_c^{(n)}V_d^{(n)}
        =
        V_{c+d}^{(n)}
\]
on the norm-dense subspace \(j_n(M)\subset L^n(M)\).  Since both sides are
bounded operators on \(L^n(M)\), they agree on all of \(L^n(M)\).  The same
argument, using the fact that \(T_0\) is the identity operator in $M$, gives
\(V_0^{(n)}=I\).

It remains to prove strong continuity.  We first establish right continuity
at \(c=0\) on the dense subspace \(j_n(M)\).  Let \(a\in M\).  Since $T_c(a)=U_{-c}aU_c$,
and \(c\mapsto U_c\) is strongly continuous,  it is straightforward that $T_c(a) \to a$ strongly as $c\downarrow0$.
It is similarly straightforward to show strong-\(*\) convergence using
\(T_c(a)^*=T_c(a^*)\).  Moreover,
\[
        \|T_c(a)\|=\|a\|,
        \qquad c\geq0,
\]
because \(T_c\) is implemented by unitary conjugation.

The family \((T_c(a))_{c\geq0}\) is therefore uniformly bounded and
converges strongly to \(a\).  By the sandwich-continuity
Lemma~\ref{lem:sandwich-continuity},
\begin{align}
        \|V_c^{(n)}j_n(a)-j_n(a)\|_n
        &=
        \|j_n(T_c(a))-j_n(a)\|_n
        \longrightarrow0
        \qquad(c\downarrow0).                         \label{eq:V-core-cont}
\end{align}

We have shown that \(\|V_c^{(n)}\|\leq1\) for every \(c\geq0\), and \(V_c^{(n)}x\to x\) as \(c\downarrow0\) for every \(x\in j_n(M)\). Proposition~\ref{prop:semigroup}, applied with the dense set \(D=j_n(M)\) and uniform bound \(1\), therefore implies that
\((V_c^{(n)})_{c\geq0}\) is a strongly continuous semigroup on
\(L^n(M)\).
\end{proof}

Having established that $V_c^{(n)}$ is a strongly continuous contraction semigroup, we define its generator as follows; cf.~Definition \ref{def:generator}.
\begin{definition}[Generator of the null translation semigroup on $L^n(M)$]\label{def:Gn-def}
    We denote the generator of $V_c^{(n)}$ by $G_n$, which is defined as:
    \[  
    G_n x\defeq \lim_{c\downarrow 0}\frac{V_c^{(n)}x-x}{c},\qquad x\in\Dom(G_n)\,,
    \]
    where the limit is in $L^n(M)$.
\end{definition}


Since $1<n<\infty$, the space $L^{n}(M)$ is reflexive \cite{Pisier:2003dsk} and its dual Banach space is
\begin{equation}
  L^{q_n}(M)=L^n(M)^*
\end{equation}
under the pairing
\begin{equation}
  \langle x,y\rangle\defeq \tr(xy),\qquad x\in L^{q_n}(M),\quad y\in L^n(M)\,.
\end{equation}
Denote the dual semigroup on $L^{q_n}(M)$ by $\bigl(V_c^{(n)}\bigr)^*$.
Because $L^{n}(M)$ is reflexive, the dual semigroup is strongly continuous on $L^{q_n}(M)$ \cite{phillips1955adjoint,engel.nagel:00:one-parameter}.  It is also positive and contractive. We now establish that the dual semigroup also has a simple action \eqref{eq:dual-core-action} on the dual dense core.


\begin{lemma}[Action on the dual dense core]\label{lem:dual-core-action-fixed}
For every \(c\geq 0\) and every \(a\in M\), let $V_c^{(n)}$ be the semigroup defined by Proposition~\ref{prop:null-translation-Ln-semigroup}. Let \((V_c^{(n)})^*\) denote the Banach adjoint semigroup on
\(L^{q_n}(M)\), with respect to the pairing
\[
        \langle x,y\rangle_{q_n,n}\defeq \tr (xy),
        \qquad x\in L^{q_n}(M),\ y\in L^n(M)\,.
\]
Then
\begin{equation}\label{eq:dual-core-action}
    (V_c^{(n)})^*j_{q_n}(a)=j_{q_n}(T_c(a)).
\end{equation}
\end{lemma}

\begin{proof}
First note that for $a,b\in M$, the duality pairing between the dense cores $j_n(M)$ and $j_{q_n}(M)$ is given by
\begin{equation}\label{eq:core-pairing}
    \tr \bigl(j_n(a)j_{q_n}(b)\bigr) = \langle a^*\Omega,J_\omega b^*\Omega\rangle_\cH ,
        \qquad a,b\in M .
\end{equation}
For \(b\) entire analytic with respect to \(\sigma^\omega\), this follows from the Tomita relation
\[
        \s_{-i/2}(b)\W = \D_\w^{1/2}b\W = J_\w b^*\W .       
\] 
combined with the KMS placement formula from Lemma~\ref{lem:KMS-placement}:
\begin{equation*}
    \tr \bigl(j_n(a)j_{q_n}(b)\bigr)=\omega\bigl(a\,\sigma_{-i/2}(b)\bigr).
\end{equation*}
For general $b\in M$, this follows from the fact that the analytic elements form a strong-$*$ dense subalgebra $M^{\text{an}}$ of $M$, as is clear from the construction of $b_\ve$ in \eqref{eq:b-eps-def}.

Fix \(a,b\in M\) and \(c\geq 0\).  By definition of the Banach adjoint,
cyclicity of the Haagerup trace for H\"older-compatible products, and the
definition of the \(L^n\)-core action,
\[
\begin{aligned}
 \left\langle (V_c^{(n)})^*j_{q_n}(b),j_n(a)\right\rangle_{q_n,n}
 &= \left\langle j_{q_n}(b),V_c^{(n)}j_n(a)\right\rangle_{q_n,n}  \\
 &= \tr \bigl(j_{q_n}(b)j_n(T_c(a))\bigr)             \\
 &= \tr \bigl(j_n(T_c(a))j_{q_n}(b)\bigr).
\end{aligned}
\]
Using the core pairing formula \eqref{eq:core-pairing},
\[
        \tr \bigl(j_n(T_c(a))j_{q_n}(b)\bigr)
        =
        \langle T_c(a)^*\Omega,J_\omega b^*\Omega\rangle_\cH 
        =
        \langle U_{-c}a^*\Omega,J_\omega b^*\Omega\rangle_\cH
        =
        \langle a^*\Omega,U_cJ_\omega b^*\Omega\rangle_\cH .
\]
Now use \(J_\omega U_{-c}J_\omega=U_c\),
\[
        J_\omega T_c(b)^*\Omega
        =
        J_\omega U_{-c}b^*\Omega
        =
        U_cJ_\omega b^*\Omega ,
\]
to write
\[
        \langle a^*\Omega,U_cJ_\omega b^*\Omega\rangle_\cH
        =
        \langle a^*\Omega,J_\omega T_c(b)^*\Omega\rangle_\cH .
\]
Applying \eqref{eq:core-pairing} again gives
\[
        \langle a^*\Omega,J_\omega T_c(b)^*\Omega\rangle_\cH
        =
        \tr \bigl(j_n(a)j_{q_n}(T_c(b))\bigr)
        =
        \left\langle j_{q_n}(T_c(b)),j_n(a)\right\rangle_{q_n,n}.
\]
Hence
\[
        \left\langle (V_c^{(n)})^*j_{q_n}(b),j_n(a)\right\rangle_{q_n,n}
        =
        \left\langle j_{q_n}(T_c(b)),j_n(a)\right\rangle_{q_n,n}
\]
for every \(a\in M\).  Since \(j_n(M)\) is dense in \(L^n(M)\), the two
elements of \(L^{q_n}(M)=(L^n(M))^*\) agree:
\[
        (V_c^{(n)})^*j_{q_n}(b)=j_{q_n}(T_c(b)).
\]
This proves the claim.
\end{proof}

Define the null translation on the dual core as follows.
\begin{definition}[Null translations on the dual dense core]
\label{def:Vc-dual-def-dense}
Fix \(c\geq0\).  Since \(T_c(M)\subset M\), define
\[
        \widetilde V_c^{(q_n)}:j_{q_n}(M)\longrightarrow j_{q_n}(M)
\]
by
\begin{equation}
\label{eq:Vc-dual-defn}
        \widetilde V_c^{(q_n)}j_{q_n}(a)
        :=
        j_{q_n}(T_c(a)),
        \qquad a\in M.
\end{equation}
\end{definition}
The map $\widetilde V_c^{(q_n)}$ extends uniquely to a positive contraction on all of $L^{q_n}(M)$.
\begin{proposition}[The null-translation dual semigroup on \(L^{q_n}(M)\)]
\label{prop:null-translation-Lqn-semigroup}
For every \(c\geq0\), the map
\(\widetilde V_c^{(q_n)}\) extends uniquely to a positive contraction
\[
        V_c^{(q_n)}:L^{q_n}(M)\longrightarrow L^{q_n}(M).
\]
The family \(V^{(q_n)}=(V_c^{(q_n)})_{c\geq0}\) is a strongly continuous
contraction semigroup.  More precisely,
\begin{align}
        &V_{c+d}^{(q_n)}
        =
        V_c^{(q_n)}V_d^{(q_n)},
        &&c,d\geq0,                                      \label{eq:V-semigroup-dual}\\
        &V_0^{(q_n)}
        =
        I,                                                \label{eq:V-identity-dual}\\
        &\lim_{c\to c_0}
        \|V_c^{(q_n)}x-V_{c_0}^{(q_n)}x\|_{q_n}
        =
        0,
        &&x\in L^{q_n}(M),\quad c_0\geq0.                    \label{eq:V-strong-cont-dual}
\end{align}
\end{proposition}
\begin{proof}
    Straightforward using the same argument as Proposition~\ref{prop:null-translation-Ln-semigroup}
\end{proof}

We now identify the map $(V_c^{(n)})^*$ with $V_c^{q_n}$.
\begin{proposition}[Equivalence of Banach adjoint semigroups]\label{prop:equal-dual-semigroups}
    Let $(V_c^{(n)})^*$ denote the Banach adjoint semigroup of $V_c^{(n)}$ as in Lemma~\ref{lem:dual-core-action-fixed}. Let $V_c^{(q_n)}$ be the unique extension of the null translations to the Haagerup $L^{q_n}(M)$ space, as in Proposition~\ref{prop:null-translation-Lqn-semigroup}. Then
    \[
        (V_c^{(n)})^*=V_c^{(q_n)}
    \]
\end{proposition}
\begin{proof}
    Lemma~\ref{lem:dual-core-action-fixed} establishes that the dual semigroup $(V_c^{(n)})^*$ has the same action on the dense core $j_{q_n}(M)\subset L^{q_n}(M)$ as $\widetilde V^{(q_n)}_c$. The uniqueness of the extension to all of $L^{q_n}(M)$ implies equality of the two maps.
\end{proof}


%%%%%%%%%%%%%%%%%%%%%%%%%%%%%%%%%%%%%%%%%%%%%%%%%%%%%
\subsection{Semigroup and modular regularization}\label{subsec:modular-semigroup-reg}
In this section we construct a semigroup and modular regularization that is used to construct the graph core for $G_n$.

\subsubsection{Semigroup smoothing for \texorpdfstring{$L^n(M)$}{LnM}}
\label{subsec:null-smoothing}

Let $C_{c,0}^\infty((0,\infty))$ denote the smooth compactly supported functions on the half-line that vanish in a neighbourhood of $0$.  
\begin{definition}[Semigroup smoothing]
    For $f\in C_{c,0}^\infty((0,\infty))$ and $a\in M$, define the ultraweak integral\footnote{Ultraweak integrals are defined in Definition~\ref{def:ultraweak-int}. We use the same ultraweak-integration convention as in Definition~\ref{def:ultraweak-int} for any bounded ultraweakly measurable $M$-valued orbit, in particular for $t\mapsto T_t(a)$.}
\begin{equation}\label{eq:def-null-smearing}
  a_f\defeq \int_0^\infty f(t)T_t(a)\,dt\in M\,.
\end{equation}
Also define
\begin{equation}\label{eq:def-null-smearing-operator}
  N_fj_n(a)\defeq j_n(a_f)\,.
\end{equation}
\end{definition}

We can then express the action of $N_f$ on the dense core in terms of a Bochner integral in $L^n(M)$.
\begin{lemma}\label{lem:Nf-action}
$N_fj_n(a)$ can be expressed as
\begin{equation}\label{eq:nul-smearing-explicit}
  j_n(a_f) = N_fj_n(a)=\int_0^\infty f(t)V_t^{(n)}j_n(a)\,dt\,,
\end{equation}
where the integral is a Bochner integral in the Banach space $L^n(M)$.
\end{lemma}
\begin{proof}
$V_t^{(n)}(j_n(a))$ is strongly continuous in $t$. Multiplication by the smooth compactly supported $f(t)$ preserves continuity. Hence, $t\mapsto f(t)V_t^{(n)}j_n(a)$ is continuous and therefore (strongly) measurable \cite[Corollary 1.1.2]{arendt.batty.ea:11:vector-valued}.  Finally, by contractivity of $V_t^{(n)}$, $$\|f(t)V_t^{(n)} j_n(a)\|_n \leq |f(t)|\|j_n(a)\|_n,$$
and since $|f(t)|\|j_n(a)\|_n$ is integrable, $f(t)V_t^{(n)}j_n(a)$ is Bochner integrable.  
The result \eqref{eq:nul-smearing-explicit} then follows from Lemma~\ref{lem:jn-compatible-weak}, with $\mathcal S= (0,\I)$, $d\nu(t)=f(t)\,dt$, and $a_t=T_t(a)$.
\end{proof}

\begin{definition}[Generator of the null translation endomorphism $T_r$]
Define the generator\footnote{This is inspired by Definition \ref{def:generator} for the generator of a semigroup. However, we note that $r\mapsto T_r$ is not expected to be a strongly continuous semigroup on $M$ with operator norm.} of $T_r$ to be $\d$,
\begin{equation}\label{eq:delta-dom-def-main}
    \d a\defeq \lim_{r\downarrow0}\frac{T_r(a)-a}{r}, \quad \Dom(\d)\defeq\lbr a\in M:\lim_{r\downarrow0}\frac{T_r(a)-a}{r} \ \text{exists in operator norm} \rbr.
\end{equation}
\end{definition}
The domain $\Dom(\delta)$ is a unital $*$-subalgebra of M, and $\delta$ acts as a derivation on its domain:
\begin{equation}
    \delta(ab)= \delta(a)b + a\delta(b), \quad a,b\in \Dom \delta.
\end{equation}
See Lemma~\ref{lem:derivation-on-dom} for a proof of this claim.


The next lemma identifies the action of $G_n$ on the semigroup-smoothed core $j_n(a_f) \in \Dom(G_n)$ as $G_n j_n(a_f) = j_n(\delta a_f)$. This provides a correspondence between the null generators $\delta$ on $M$ and $G_n$ on $L^n(M)$ that we will use throughout the rest of this Section.
\begin{lemma}[Action of $G_n$ on semigroup smooth elements]
\label{lem:null-smoothing-fixed}
For $f\in C_{c,0}^\infty((0,\infty))$ and $a\in M$, we have $a_f\in\Dom(\d)$ with
\begin{equation}\label{eq:delta-af-defn}
    \delta a_f\defeq -\int_0^\infty f'(t)T_t(a)\,dt\in M\,.
\end{equation}
Equivalently, $N_fj_n(a)\in\Dom(G_n)$, with
\begin{equation}\label{eq:Gn-Nf-jn-a}
  G_nN_fj_n(a) = -\int_0^\infty f'(t)V_t^{(n)}j_n(a)\,dt\,,
\end{equation}
which can be expressed as
\begin{equation}\label{eq:delta-defn}
  G_nj_n(a_f)=j_n(\delta a_f)\,.
\end{equation}
\end{lemma}
\begin{proof}
We first establish that $a_f\in \Dom(\delta)$. Extend $f$ by zero to all of $\mathbb R$. Then
\begin{equation*}
    \frac{T_r(a_f)-a_f}{r}=\int_0^\I \frac{f(t-r)-f(t)}{r} T_t(a)\, dt.
\end{equation*}
Since $\frac{f(\cdot-r)-f(\cdot)}{r} \to -f'$ in $L^1((0,\I))$, and $\|T_r(a)\|\leq \|a\| $, we get
\begin{equation*}
    \Big\|\frac{T_r(a_f)-a_f}{r}+ \int_0^\I f'(t) T_t(a)\,dt\Big \| \leq \|a\| \Big\| \frac{f(\cdot-r)-f(\cdot)}{r}+f' \Big\|_{L^1} \to 0.
\end{equation*}
Thus,
\begin{equation*}
    a_f\in \Dom (\delta), \quad \delta a_f\defeq -\int_0^\infty f'(t)T_t(a)\,dt\in M\,.
\end{equation*}
For $r>0$, since $a_f\in \Dom(\delta)$,
\begin{equation*}
    \frac{V_r^{(n)}j_n(a_f)-j_n(a_f)}{r}=j_n \lb \frac{T_r(a_f)-a_f}{r}\rb \longrightarrow j_n(\delta a_f) \qquad\text{in }L^n(M),
\end{equation*}
by sandwich continuity Lemma~\ref{lem:sandwich-continuity}. Hence,
\begin{equation*}
    j_n(a_f) \in \Dom(G_n), \quad G_n j_n(a_f)= j_n(\delta a_f) \,.
\end{equation*}
Finally, Lemma~\ref{lem:jn-compatible-weak} gives
\begin{equation*}
    j_n(\delta a_f)= -\int_0^\I f'(t) V_t^{(n)} j_n(a)\, dt. 
\end{equation*}
\end{proof}

\subsubsection{Modular smoothing for \texorpdfstring{$L^n(M)$}{LnM} }
\label{subsec:modular-smoothing}

To define a modular regularization in $L^n(M)$, we first define the $L^n(M)$-extension of the modular automorphism group.
\begin{definition}[Modular flow extended to $L^n(M)$]
    Let \((\Sigma_s^{(n)})_{s\in\mathbb R}\) be the \(L^n(M)\) extension of the modular automorphism group $\s_s$, acting on the dense core by
\begin{equation}
  \Sigma_s^{(n)} j_n(a)\defeq j_n(\sigma_s(a)), \qquad a\in M .
\end{equation}
\end{definition}
Lemma~\ref{lem:modular-Ln-group} shows that \(\Sigma_s^{(n)}\) is a strongly continuous group of isometries on \(L^n(M)\).

\begin{lemma}
    For $y\in\Dom(G_n)$,
\begin{equation}\label{eq:Sigma-y-domain}
  \Sigma_s^{(n)}y\in\Dom(G_n), \qquad G_n\Sigma_s^{(n)}y=e^{2\pi s}\Sigma_s^{(n)}G_ny\,.
\end{equation}
\end{lemma}
\begin{proof}
    By Borchers covariance and \eqref{eq:Vc-defn},
\begin{equation}
  \Sigma_s^{(n)}V_t^{(n)}=V_{te^{-2\pi s}}^{(n)}\Sigma_s^{(n)},
  \qquad s\in\mathbb R,
  \quad t\geq 0\,.
\end{equation}
Equivalently,
\begin{equation}\label{eq:semigroup-covariance}
  V_t^{(n)}\Sigma_s^{(n)}=\Sigma_s^{(n)}V_{te^{2\pi s}}^{(n)}\,.
\end{equation}
Eq \eqref{eq:Sigma-y-domain} then follows from Definition~\ref{def:Gn-def}.
\end{proof}

Using the modular flow extended to $L^n(M)$ we now define the modular smoothing. 
\begin{definition}[Modular smoothing]
    We define a modular Gaussian smoothing operator for $\varepsilon>0$
\begin{equation}\label{eq:modular-gaussian-smoothing}
  \mathcal M_\varepsilon^{(n)}y\defeq \int_{\mathbb R}\gamma_\varepsilon(s)\Sigma_s^{(n)}y\,ds\,,
\end{equation}
where $\gamma_\varepsilon(s)\defeq \frac{1}{\sqrt{2\pi}\,\varepsilon}
  \exp\left(-\frac{s^2}{2\varepsilon^2}\right)$, as in \eqref{eq:gamma-appendix}, and the integral is a Bochner integral (Definition \ref{def:Bochner-int}). 
\end{definition}
The following lemma shows that the smoothed operators $\mathcal M_\ve^{(n)} y$ converge to $y$ as $\ve\to0$ in the graph norm for $G_n$ (Definition \ref{def:graph}).
\begin{lemma}[Graph convergence of modular smoothing]
\label{lem:modular-smoothing-fixed}
If $y\in\Dom(G_n)$, then $\mathcal M_\ve^{(n)}y$ converges in $G_n$-graph norm to $y$, i.e., as $\ve\to0$,
\begin{equation}
    \mathcal M_\varepsilon^{(n)}y\longrightarrow y,
  \qquad
  G_n\mathcal M_\varepsilon^{(n)}y\longrightarrow G_ny
  \qquad
  \text{in }L^n(M).
\end{equation}
\end{lemma}
\begin{proof}
We first show that \(\mathcal M_\varepsilon^{(n)}y\in\Dom (G_n)\) and compute its
\(G_n\)-image.  Since \(s\mapsto \Sigma_s^{(n)}y\) is norm-continuous and
\(\Sigma_s^{(n)}\) is an isometry on $L^n(M)$,
\[
        \int_{\mathbb R}\gamma_\varepsilon(s)
        \|\Sigma_s^{(n)}y\|_n\,ds
        =
        \|y\|_n
        <\infty .
\]
Thus $\cM_\ve^{(n)}y$ is a Bochner integral in $L^n(M)$.

Moreover, by \eqref{eq:Sigma-y-domain},
\[
        \left\|
        \gamma_\varepsilon(s)G_n\Sigma_s^{(n)}y
        \right\|_n
        =
        \gamma_\varepsilon(s)e^{2\pi s}\|G_ny\|_n .
\]
which is integrable because
\[
        \int_{\mathbb R}\gamma_\varepsilon(s)e^{2\pi s}\,ds
        =
        \exp(2\pi^2\varepsilon^2)
        <\infty .
\]
Let \(\mathcal G(G_n)\subset L^n(M)\oplus L^n(M)\) denote the graph of \(G_n\). 
Since \(G_n\) is closed (Theorem \ref{thm:generator-is-closed}), \(\mathcal G(G_n)\) is a closed subspace of
\(L^n(M)\oplus L^n(M)\), hence a Banach space.  The map
\[
        s\longmapsto
        \left(
            \gamma_\varepsilon(s)\Sigma_s^{(n)}y,\,
            \gamma_\varepsilon(s)e^{2\pi s}\Sigma_s^{(n)}G_ny
        \right)
\]
is Bochner integrable as a \(\mathcal G(G_n)\)-valued function. Therefore its
Bochner integral again lies in \(\mathcal G(G_n)\).  Equivalently,
\begin{equation}\label{eq:M-eps-domain}
    \mathcal M_\varepsilon^{(n)}y\in\Dom (G_n), 
    \qquad 
    G_n\mathcal M_\varepsilon^{(n)}y
    =
    \int_{\mathbb R}
    \gamma_\varepsilon(s)e^{2\pi s}
    \Sigma_s^{(n)}G_ny\,ds .
\end{equation}

It remains to prove the two convergence statements.  
The first convergence follows from Lemma~\ref{lem:approx-id} applied to $F(s)=\Sigma_s^{(n)}y$ and the approximate identity $\gamma_\varepsilon(s)\,ds$ (Definition \ref{def:approx-id}). 
For the generator part, define the finite measure
\[
        d\mu_\varepsilon(s)
        \defeq 
        \gamma_\varepsilon(s)e^{2\pi s}\,ds .
\]
Completing the square gives
\[
        \gamma_\varepsilon(s)e^{2\pi s}
        =
        e^{2\pi^2\varepsilon^2}
        \frac{1}{\sqrt{2\pi}\varepsilon}
        \exp\!\left(
            -\frac{(s-2\pi\varepsilon^2)^2}{2\varepsilon^2}
        \right).
\]
Hence
\(
        \mu_\varepsilon(\mathbb R)
        =
        e^{2\pi^2\varepsilon^2}
        \to 1 ,
\) 
and $\mu_\ve(\{|s|\geq\delta\})\to0$ for every $\d>0$. Thus $\mu_\ve$ is an approximate identity.  Applying Lemma~\ref{lem:approx-id} to $F(s)= \Sigma_s^{(n)} G_n y$ gives
\[
        G_n\mathcal M_\varepsilon^{(n)}y
        =
        \int_{\mathbb R}
        \gamma_\varepsilon(s)e^{2\pi s}
        \Sigma_s^{(n)}G_ny\,ds
        \longrightarrow
        G_ny
        \qquad\text{in }L^n(M).
\]
\end{proof}

The modular-Gaussian smoothing is compatible with the dense embedding $j_n$ as follows.
\begin{lemma}\label{lem:modular-gaussian-dense-core}
Let $\mathcal M_\varepsilon^{(n)}$ be the modular-Gaussian smoothing defined in \eqref{eq:modular-gaussian-smoothing}. 
Then
\begin{equation}\label{eq:modular-gaussian-dense-core-act}
     \mathcal M_\varepsilon^{(n)}j_n(a)=j_n\!\left(\int_{\mathbb R}\gamma_\varepsilon(s)\,\sigma_s(a)\,ds\right),\qquad a\in M,
\end{equation}
as an ultraweak integral in $M$. 
\end{lemma}

\begin{proof}
Follows from Lemma~\ref{lem:jn-compatible-weak}, with $\mathcal S=\mathbb R$, $d\nu(s)= \gamma_{\varepsilon}(s)\, ds$, and $a_s = \sigma_s(a)$.
\end{proof}

\subsection{A Ward graph core of the semigroup generator}
\label{subsec:ward-graph-core}

In this section we combine the semigroup and modular regularizations to define a graph core of the semigroup generator which is modular entire. 

For $a\in M$, $f\in C_{c,0}^\infty((0,\infty))$, and $\varepsilon>0$, define the ultraweak integral
\begin{equation}\label{eq:defn-double-smearing}
  a_{f,\varepsilon}\defeq \int_{\mathbb R}\gamma_\varepsilon(s)\sigma_s(a_f)\,ds\in M\,.
\end{equation} 
The element $a_{f,\varepsilon}$ is entire analytic for the modular group (Lemma~\ref{lem:gaussian-modular-entire}).  More explicitly, for $z\in\mathbb C$,
\begin{equation}
  \sigma_z(a_{f,\varepsilon})
  =\int_{\mathbb R}\gamma_\varepsilon(s-z)\sigma_s(a_f)\,ds\,,
\end{equation}
where the integral is understood in the usual weakly entire sense. The smoothened operators $a_{f,\varepsilon}$ allow us to define null derivatives in an appropriate sense on a smooth core.
\begin{definition}[Ward core]\label{def:smooth-core-null}
    Let $M^W\subset M$ %(W for Ward) 
    be the unital $*$-subalgebra generated by all elements of the form $a_{f,\varepsilon}$ with $a\in M$, $f\in C_{c,0}^\infty((0,\infty))$, and $\varepsilon>0$. Define
    \[
        \cD_n^{W}\defeq j_n(M^W)\subset L^n(M).
    \]
\end{definition}
\begin{proposition}[Smooth core for the null derivation]
\label{lem:smooth-core-delta}
Let $M^W\subset M$ and  $\cD_n^{W}$ be as in Definition~\ref{def:smooth-core-null}.  Then:
\begin{enumerate}
  \item[(i)] $M^W\subset\Dom(\delta)$, and $\delta$ acts as a derivation on $M^W$:
  \begin{equation}
  \label{eq:delta-derivation}
    \delta(ab)=\delta(a)b+a\,\delta(b),\qquad a,b\in M^W\,.
  \end{equation}
  \item[(ii)] Every $a\in M^W$ is modular-entire, and so is $\delta (a)$. Moreover, $\sigma_z(a)\in \Dom(\delta)$ for all $a\in M^W$, $z\in \mathbb C$, and Borchers covariance relation holds in complex form:
  \begin{equation}
  \label{eq:complex-borchers}
    \delta(\sigma_z(a))=e^{2\pi z}\sigma_z(\delta a),\qquad z\in\mathbb C\,.
  \end{equation}
  \item[(iii)] $\cD_n^{W}\subset\Dom(G_n)$, and
  \begin{equation}
  \label{eq:Gn-on-core}
    G_nj_n(a)=j_n(\delta a),\qquad a\in M^W\,.
  \end{equation}
\end{enumerate}
\end{proposition}

\begin{proof}
\emph{of (i).} 

For each generator $a_{f,\varepsilon}$, consider the difference quotient
\begin{equation}\label{eq:difference-quot-Tr}
  \frac{T_r(a_{f,\varepsilon})-a_{f,\varepsilon}}{r}
  =\int_{\mathbb R}\gamma_\varepsilon(s)\, e^{2\pi s}\sigma_s\!\left(\frac{T_{re^{2\pi s}}(a_f)-a_f}{r e^{2\pi s}}\right)ds\,.
\end{equation}
Since \eqref{eq:difference-quot-Tr} is an ultraweak integral, using
\begin{equation*}
    \left\|\int F_r(s)ds\right\| = \sup_{\|\vf\|_{M_*}\leq 1}\lv \int\vf\lb F_r(s)\rb ds\rv \leq \int \|F_r(s)\|ds,
\end{equation*}
we get
\begin{equation*}
    \Bigg\| \frac{T_r(a_{f,\varepsilon})-a_{f,\varepsilon}}{r}-\int_\mathbb{R}\gamma_\varepsilon(s) e^{2\pi s} \sigma_s(\delta a_f)\Bigg\| \leq \int_\mathbb{R}\gamma_\varepsilon(s)e^{2\pi s} \Bigg\|\frac{T_{re^{2\pi s}}(a_f)-a_f}{r e^{2\pi s}}-\delta a_f\Bigg\|\, ds \,.
\end{equation*}
The integrand in the upper bound above tends point-wise to zero as $r\downarrow0$ since $a_f \in \Dom(\delta)$.
Specializing the Lipschitz bound from Lemma~\ref{lem:lipschitz-af} to $u=re^{2\pi s}>0$ and using the operator-norm isometry of $\sigma_s$ implies that the integrand in the difference quotient for $a_{f,\varepsilon}$ \eqref{eq:difference-quot-Tr} is dominated by 
$2\gamma_\varepsilon(s)\,e^{2\pi s}\,\|\delta a_f\|$ uniformly in $r>0$. The dominating bound is integrable since
\begin{equation*}
  \int_{\mathbb R}\gamma_\varepsilon(s)\,e^{2\pi s}\,ds=e^{2\pi^2\varepsilon^2}<\I.
\end{equation*}
By dominated convergence we obtain $a_{f,\varepsilon}\in\Dom(\delta)$ with
\begin{equation}
\label{eq:delta-on-generators}
  \delta(a_{f,\varepsilon})=\int_{\mathbb R}\gamma_\varepsilon(s)\,e^{2\pi s}\sigma_s(\delta a_f)\,ds\,.
\end{equation}
Membership in $\Dom(\delta)$ is preserved under the $*$-operation: since $T_t$ is a $*$-endomorphism, $a_{f,\varepsilon}^*=(a^*)_{\bar f,\varepsilon}\in\Dom(\delta)$. 
Thus, $M^W$, the unital $*$-subalgebra generated by these elements  is contained in $\Dom(\delta)$, and the derivation property \eqref{eq:delta-derivation} follows from Lemma~\ref{lem:derivation-on-dom}.

\emph{Proof of (ii)}. 

Each generator $a_{f,\varepsilon}$ is modular-entire by construction (see Lemma \ref{lem:gaussian-modular-entire}), with explicit complex extension
\begin{equation*}
  \sigma_z(a_{f,\varepsilon})
  =\int_{\mathbb R}\gamma_\varepsilon(s-z)\sigma_s(a_f)\,ds,\qquad z\in\mathbb C\,.
\end{equation*}
Modular-entirety is preserved under sums, products, and adjoints, so all of $M^W$ is modular-entire.  
Since $\gamma_\varepsilon(s)e^{2\pi s}$ is a shifted Gaussian, the derivative in  \eqref{eq:delta-on-generators} is also a modular analytic element with complex extension
\begin{equation}\label{eq:sigmaz-afe}
    \sigma_z(\delta(a_{f,\varepsilon}))=\int_{\mathbb R}\gamma_\varepsilon(s-z)\,e^{2\pi( s-z)}\sigma_s(\delta a_f)\,ds\,.
\end{equation}
This extends to the full algebra $M^W$, i.e., $\delta(M^W)$ is modular entire,  by the derivation rule from (i).

By the real Borchers
covariance relation, $T_t\sigma_s=\sigma_sT_{te^{2\pi s}}$,
we have 
\begin{align}
 \frac{T_r(\sigma_z(a_{f,\varepsilon}))
             -\sigma_z(a_{f,\varepsilon})}{r}
 &=
 \int_{\mathbb R}
   \gamma_\varepsilon(s-z)
   \frac{T_r(\sigma_s(a_f))-\sigma_s(a_f)}{r}\,ds
 \notag\\
 &=
 \int_{\mathbb R}
   \gamma_\varepsilon(s-z)e^{2\pi s}
   \sigma_s\left(
      \frac{T_{re^{2\pi s}}(a_f)-a_f}
           {re^{2\pi s}}
   \right)\,ds.
 \label{eq:complex-modular-difference-quotient}
\end{align}

For each fixed \(s\in\mathbb R\),
\[
 \frac{T_{re^{2\pi s}}(a_f)-a_f}{re^{2\pi s}}
 \longrightarrow \delta a_f
 \qquad\text{in operator norm as }r\downarrow0.
\]
Moreover, the Lipschitz estimate of
Lemma~\ref{lem:lipschitz-af}, and the operator norm isometry of $\sigma_s$ implies that the integrand in the difference quotient \eqref{eq:complex-modular-difference-quotient} is dominated by $2|\gamma_{\varepsilon}(s-z)| e^{2\pi s}\|\delta a_f\|$.
For fixed \(z\in\mathbb C\),
\[
 |\gamma_\varepsilon(s-z)|
 =
 e^{(\Im z)^2/(2\varepsilon^2)}
 \gamma_\varepsilon(s-\Re z),
\]
and therefore
\(
 s\longmapsto
 |\gamma_\varepsilon(s-z)|e^{2\pi s}
\)
is integrable.


Dominated convergence in operator norm now yields
\begin{equation}\label{eq:delta-complex-modular-generator}
 \delta\bigl(\sigma_z(a_{f,\varepsilon})\bigr)
   =
   \int_{\mathbb R}
      \gamma_\varepsilon(s-z)e^{2\pi s}
      \sigma_s(\delta a_f)\,ds.
\end{equation}
In particular,
\(\sigma_z(a_{f,\varepsilon})\in\Dom(\delta)\).
Combining \eqref{eq:sigmaz-afe} and \eqref{eq:delta-complex-modular-generator}, we get that \eqref{eq:complex-borchers} holds for every generator
\(a_{f,\varepsilon}\). Since \(\sigma_z\) is multiplicative on
modular-entire elements, \(\Dom(\delta)\) is a unital
\(*\)-algebra, and \(\delta\) is a derivation on its domain, the
identity is preserved under finite sums and products. Moreover,
\((a_{f,\varepsilon})^*=(a^*)_{\bar f,\varepsilon}\), so it is also
preserved under adjoints. Hence
\[
 \delta(\sigma_z(a))=e^{2\pi z}\sigma_z(\delta a),
 \qquad a\in M^W.
\]

\emph{Proof of (iii).}

Let $a\in M^W\subset\Dom(\delta)$.  Since $a\in \Dom(\delta)$, \eqref{eq:Vc-defn} gives as $r\downarrow0$
\begin{equation*}
  \frac{V_r^{(n)}j_n(a)-j_n(a)}{r}
  =j_n\!\left(\frac{T_r(a)-a}{r}\right)
  \longrightarrow j_n(\delta a) \qquad\text{in }L^n(M),
\end{equation*}
where the convergence uses operator-norm convergence of the difference quotient on $a$ together with sandwich continuity (Lemma~\ref{lem:sandwich-continuity}).  Hence $j_n(a)\in\Dom(G_n)$ with $G_nj_n(a)=j_n(\delta a)$, and $\cD_n^{W}\subset\Dom(G_n)$.
\end{proof}

Before proving that $\mathcal D_n^W$ is a graph core for $G_n$, it is useful to recall the definition of the resolvent of a semigroup generator. See Definition~\ref{def:resolvent-appendix} and Appendix \ref{appsubsec:semigrp} for more details.
\begin{remark}[Resolvent of $G_n$]\label{rem:resolvent}
Denote by $\rho(G_n)$ the resolvent set of $G_n$, defined as:
\[
    \rho(G_n)\defeq \{\lambda\in \mathbb C: \lambda-G_n \text { has a bounded inverse}\}\,.
\]
For $\l\in\r(G_n)$, denote the resolvent by
\[
    R(\lambda,{G_n}) \defeq (\lambda-G_n)^{-1}\,.
\]
By the Hille–Yosida theorem (\cite{engel.nagel:00:one-parameter} and Theorem \ref{thm:hille-yosida}), $G_n$ is a closed densely defined operator with $(0,\I)\subset \rho(G_n)$, and $\|R(\lambda,{G_n})\|\leq \frac1\lambda$. Moreover, the resolvent admits the Laplace representation:
\[
    R(\lambda,{G_n}):=\int_0^\I e^{-\lambda t} V_t^{(n)} \,dt\,:=\lim_{s\to\I}\int_0^s e^{-\lambda t} V_t^{(n)}\, dt\,.
\]
\end{remark}

We now combine the resolvent of $G_n$ with the null and modular smoothings to prove that $\cD_n^{W}$ is dense in $\Dom(G_n)$ in the graph norm.
\begin{proposition}[Ward graph core]
\label{prop:ward-graph-core-fixed}
The subspace $\cD_n^{W}$ is a graph core for $G_n$.  Equivalently, for every $y\in\Dom(G_n)$, there exists a sequence $y_m\in\cD_n^{W}$ such that
\[
  y_m\longrightarrow y,
  \qquad
  G_ny_m\longrightarrow G_ny
  \qquad 
  \text{in }L^n(M).
\]
\end{proposition}

\begin{proof}
The proof has two stages.  First, we use the resolvent of $G_n$ to show that $j_n(M)\cap\Dom(G_n)$ is a graph core.  Second, we approximate each resolvent-core vector by an element of $\cD_n^{W}$ using null and modular-Gaussian smoothing.

The core $j_n(M)$ is norm dense in $L^n(M)$.  
For $a\in M$ define the ultraweak integral
\begin{equation*}
  a_\lambda\defeq \int_0^\infty e^{-\lambda t}T_t(a)\,dt\in M\,.
\end{equation*}
For every normal functional $\varphi \in M_*$, $t\mapsto e^{-\lambda t}\varphi(T_t(a))$ is integrable. Hence the ultraweak integral exists by norm-dual boundedness.

With $R(\l,G_n)$ as defined in Remark \ref{rem:resolvent}, we claim that
\begin{equation}\label{eq:rlambda-inclusion}
  R(\l,G_n)j_n(a)=j_n(a_\lambda)\,,
\end{equation}
so that $R(\l,G_n)$ maps $j_n(M)$ into itself. This follows from Lemma~\ref{lem:jn-compatible-weak}, with $\mathcal S= (0,\I)$, $d\nu(t) = e^{-\lambda t} \,dt$, and $a_t= T_t(a)$. 
Thus $R(\l,G_n)j_n(M)\subset j_n(M)\cap\Dom(G_n)$ for $\l\in(0,\I)\subset\r(G_n)$, and $j_n(M)\cap\Dom(G_n)$ is a graph core by the resolvent-core criterion Lemma~\ref{lem:resolvent-core}. This follows straightforwardly from the fact that $(\l-G_n)R(\l,G_n)j_n(M)=j_n(M)$ is dense in $L^n(M)$.

It remains to approximate the resolvent-core vectors $R(\l,G_n)j_n(a)$ by vectors in $\cD_n^{W}$.  
Consider
\begin{equation*}
    f_m(t) = g(mt)k(t/m)e^{-\l t}\,,
\end{equation*}
with $g,k\in C^\I([0,\I))$ such that $g$ vanishes in a neighbourhood of $0$, $ g(u)=1$ for $u\geq1$, and $k(u)=1$ for $u\in[0,1]$, $k(u)=0$ for $u\geq2$. Then, $f_m\in C_{c,0}^\infty((0,\infty))$ and $f_m\to e^{-\lambda t}$ in $L^1([0,\infty))$. 
Therefore,
\begin{equation*}
  N_{f_m}j_n(a) = \int_0^\I f_m(t)V_t^{(n)}j_n(a)dt \longrightarrow R(\l,G_n)j_n(a)\,,
\end{equation*}
since
\begin{equation*}
  \left\|\int_0^\infty \bigl(f_m(t)-e^{-\lambda t}\bigr)V_t^{(n)}j_n(a)\,dt\right\|_n
  \leq \|f_m-e^{-\lambda\, \cdot\,}\|_{L^1}\|j_n(a)\|_n \to 0\,.
\end{equation*}
The generator part is
\begin{equation*}
    G_nN_{f_m}j_n(a) = - \int_0^\I f_m'(t)V_t^{(n)}j_n(a)dt = A_m+B_m+C_m\,,
\end{equation*}
where
\begin{align*}
    A_m =&\ -\int_0^\I mg'(mt)k(t/m)e^{-\l t}V_t^{(n)}j_n(a)dt\longrightarrow -j_n(a), \\
    B_m =&\ -\frac1m\int_0^\I g(mt)k'(t/m)e^{-\l t}V_t^{(n)}j_n(a)dt \longrightarrow0, \\
    C_m =&\ \l\int_0^\I g(mt)k(t/m)e^{-\l t}V_t^{(n)}j_n(a)dt \longrightarrow \l R(\l,G_n)j_n(a)\,.
\end{align*}
The limits are computed as follows. For \(A_m\), substitute \(u=mt\) and use the strong continuity of
\(V_t^{(n)}\) at \(t=0\), together with
\(\int_0^\infty g'(u)\,du=1\). The limit for \(B_m\) follows
from the prefactor \(m^{-1}\), contractivity of \(V_t^{(n)}\), and
compact support of \(k'\). The convergence of \(C_m\) follows from
\(f_m\to e^{-\lambda\cdot}\) in \(L^1(0,\infty)\).

Combining,
\begin{equation*}
    G_nN_{f_m}j_n(a) \longrightarrow -j_n(a) + \l R(\l,G_n)j_n(a) = G_n R(\l,G_n)j_n(a)\,.
\end{equation*}
So, $N_{f_m}j_n(a)\to R(\l,G_n)j_n(a)$ in graph norm.

For each fixed $m$, Lemma~\ref{lem:modular-smoothing-fixed} gives
\begin{equation*}
  \mathcal M_\varepsilon^{(n)}N_{f_m}j_n(a)\longrightarrow N_{f_m}j_n(a)
\end{equation*}
 in the graph norm as $\varepsilon\downarrow 0$.  
 Moreover, by Lemma~\ref{lem:modular-gaussian-dense-core}
\begin{equation*}
  \mathcal M_\varepsilon^{(n)}N_{f_m}j_n(a)
  =j_n(a_{f_m,\varepsilon})\in\cD_n^{W}\,.
\end{equation*}
Choosing a diagonal sequence in $m$ and $\varepsilon$ gives a graph-norm approximation of each resolvent-core vector by elements of $\cD_n^{W}$.  Combining this with the resolvent-core approximation proves that $\cD_n^{W}$ is a graph core for $G_n$.
\end{proof}


\subsection{The cyclic \texorpdfstring{$n$}{n}-point Ward identity}
\label{subsec:cyclic-ward-identity-fixed}

With $\cD_n^W$ established as the graph core for $G_n$, we now state and prove the cyclic Ward identity.

For $y_1,\ldots,y_n\in L^n(M)$ define the cyclic $n$-linear form
\begin{equation}
  \Lambda_n(y_1,\ldots,y_n)\defeq \tr(y_1y_2\cdots y_n)\,.
\end{equation}
By H\"older's inequality,
\begin{equation}
  |\Lambda_n(y_1,\ldots,y_n)|\leq \prod_{k=1}^n\|y_k\|_n\,.
\end{equation}

\begin{proposition}[Core cyclic Ward identity]
\label{prop:core-cyclic-ward-fixed}
For $a_1,\ldots,a_n\in M^W \subset \Dom(\delta)$,
\begin{equation}\label{eq:ward-on-core-An}
    \sum_{k=0}^{n-1}\zeta_n^k
  \Lambda_n\bigl(
    j_n(a_1),\ldots,j_n(a_k),j_n(\delta a_{k+1}),j_n(a_{k+2}),\ldots,j_n(a_n)
  \bigr)=0\,,
\end{equation}
where $\zeta_n =e^{-2\pi i/n}$.
Equivalently, for $y_1,\ldots,y_n\in\cD_n^{W}$,
\begin{equation}\label{eq:ward-on-core-Dn}
    \sum_{k=0}^{n-1}\zeta_n^k
  \Lambda_n\bigl(
    y_1,\ldots,y_k,G_ny_{k+1},y_{k+2},\ldots,y_n
  \bigr)=0\,.
\end{equation}
\end{proposition}

\begin{proof}
For entire analytic $a_1,\ldots,a_n$, the KMS placement (Lemma~\ref{lem:KMS-placement}) of the symmetric $L^n$ insertions gives
\begin{equation*}
  \Lambda_n(j_n(a_1),\ldots,j_n(a_n))
  =\omega\bigl(a_1\sigma_{-i/n}(a_2)\sigma_{-2i/n}(a_3)\cdots\sigma_{-(n-1)i/n}(a_n)\bigr)\,.
\end{equation*}
Since $M^W$ consists of modular-entire elements, this applies to all $a_i\in M^W$.

Let
\begin{equation*}
  B\defeq a_1\sigma_{-i/n}(a_2)\sigma_{-2i/n}(a_3)\cdots\sigma_{-(n-1)i/n}(a_n)\,.
\end{equation*}
By Proposition \ref{lem:smooth-core-delta}, the algebras $M^W$ and $\sigma_z(M^W)$ are contained in $\Dom(\delta)$, and by Lemma~\ref{lem:derivation-on-dom} $\delta$ is a derivation on $\Dom(\delta)$.  Hence $B\in\Dom(\delta)$.  Since $\omega\circ T_t=\omega$, the scalar
\begin{equation*}
  \omega(T_t(B))
\end{equation*}
 is independent of $t\geq 0$.  Taking a right derivative $t=0$ in operator norm (in the sense of \eqref{eq:delta-dom-def-main}), we get
\begin{equation*}
  0=\omega(\delta B).
\end{equation*}
Using the derivation rule,
\begin{align*}
  0
  ={}&\sum_{k=0}^{n-1}
  \omega\bigl(
    a_1\sigma_{-i/n}(a_2)\cdots
    \delta(\sigma_{-ki/n}(a_{k+1}))
    \cdots\sigma_{-(n-1)i/n}(a_n)
  \bigr)\,.
\end{align*}
For $k=0$, the differentiated factor is $\delta a_1$.  For $k\geq 1$, Borchers covariance from Proposition~\ref{lem:smooth-core-delta} gives
\begin{equation*}
  \delta(\sigma_{-ki/n}(a))
  =e^{-2\pi ik/n}\sigma_{-ki/n}(\delta a)
  =\zeta_n^k\sigma_{-ki/n}(\delta a)\,.
\end{equation*}
Substituting this into the previous expression and translating back through the KMS placement formula gives the first displayed Ward identity.  The equivalent form follows from
\begin{equation*}
  G_nj_n(a)=j_n(\delta a),\qquad a\in M^W\,.
\end{equation*}
\end{proof}

\begin{theorem}[Ward identity on the full generator domain]
\label{thm:full-domain-ward-fixed}
For all $y_1,\ldots,y_n\in\Dom(G_n)$,
\begin{equation}\label{eq:full-domain-ward-fixed}
    \sum_{k=0}^{n-1}\zeta_n^k
  \Lambda_n\bigl(
    y_1,\ldots,y_k,G_ny_{k+1},y_{k+2},\ldots,y_n
  \bigr)=0\,.
\end{equation}
\end{theorem}

\begin{proof}
Equip $\Dom(G_n)$ with the graph norm
\begin{equation*}
  \|y\|_{G_n}\defeq \|y\|_n+\|G_ny\|_n\,.
\end{equation*}
For each $k$, H\"older's inequality gives
\begin{equation*}
  \bigl|
    \Lambda_n(y_1,\ldots,y_k,G_ny_{k+1},y_{k+2},\ldots,y_n)
  \bigr|
  \leq \|G_ny_{k+1}\|_n\prod_{j\neq k+1}\|y_j\|_n\,.
\end{equation*}
Thus the left-hand side of the claimed identity is continuous on $\Dom(G_n)^n$ with respect to the product graph norm.  Since $\cD_n^{W}$ is a graph core by Proposition~\ref{prop:ward-graph-core-fixed}, the identity extends from Proposition~\ref{prop:core-cyclic-ward-fixed} to all of $\Dom(G_n)^n$.
\end{proof}

\begin{remark}[Replica geometry and modular monodromy]
\label{rem:replica-monodromy}
The cyclic Ward identity in Theorem~\ref{thm:full-domain-ward-fixed} admits a transparent replica-geometric interpretation that makes contact with the replica-trick computation of R\'enyi entropies in QFT~\cite{Calabrese:2004eu}.  The KMS placement formula expresses the integer cyclic trace $\Lambda_n(j_n(a_1),\ldots,j_n(a_n))$ as a single vacuum expectation value with the $n$ insertions placed at imaginary modular times $z_k=-ik/n$ for $k=0,\ldots,n-1$.  These insertion points are evenly spaced around the modular thermal circle.  In this picture, $Q_n=\tr(x^n)$ is an algebraic analogue of the replica partition function on the $n$-fold cover. The phases $\zeta_n^k=e^{-2\pi ik/n}$ in the Ward identity then admit an  interpretation as \emph{modular monodromies}. This parallels the monodromy structure of twist operators in two-dimensional CFT, where the replica geometry produces analogous rotation phases~\cite{Calabrese:2004eu}.
\end{remark}

The only output of this section used below in the log-convexity proof is
Theorem~\ref{thm:full-domain-ward-fixed}, or equivalently its
specialization \eqref{eq:cyclic-identity-logconvexity-form}.  The smoothing algebra
\(M^W\) and the graph core \(D_n^W\) are used only to establish this
full-domain identity.
\begin{corollary}
\label{cor:cyclic-identity-logconvexity-form}
Let \(x\in\Dom(G_n^2)\).  Then
\begin{equation}
\label{eq:cyclic-identity-logconvexity-form}
\operatorname{tr}\bigl((G_n^2x)x^{n-1}\bigr)
+
\sum_{j=0}^{n-2}
\zeta_n^{j+1}
\operatorname{tr}\bigl(
(G_nx)x^j(G_nx)x^{n-2-j}
\bigr)
=0.
\end{equation}
\end{corollary}
\begin{proof}
    Follows by applying Theorem~\ref{thm:full-domain-ward-fixed} with $y_1=G_n x$, and $y_2=y_3=\hdots=y_n=x$.
\end{proof}
%%%%%%%%%%%%%%%%%%%%%%%%%%%%%%%%%%%%%%%%%%%%%%%%%%%%%%%%%%%%%%%%%%%%%%%%%%%%%%%%%%%%%%%%%%%%%%%

\section{Log-convexity of \texorpdfstring{$L^n$}{Ln} norms under HSMI}\label{sec:integer-renyi-qnec}

In this section we use the cyclic Ward identity to prove the integer RQNEC. We first prove log-convexity of $Q_n(\psi_c\|\omega;M)$, defined in Lemma \ref{lem:density-propagation-fixed}, for positive initial data in $\Dom(G_n^2)$, then relax this regularity assumption using Yosida regularization, obtaining the integer RQNEC for general functionals with finite SRD.  

Fix an integer $n\geq2$. Assume a HSMI $M_B\subset M_A=:M$ with a common faithful normal state $\omega$, and the consequences from Theorem~\ref{thm:modulartranslations}.
The main inputs for this section are the definition of the SRD (Definition~\ref{def:SRD}, Lemma~\ref{lem:y-op-defn-SRD}), its fixed algebra formulation (Lemma~\ref{lem:fixedalgebraSRD}), the $L^n$ semigroup $(V_c^{(n)})_{c\geq0}$ (Proposition~\ref{prop:null-translation-Ln-semigroup}) and the dual semigroup $(V_c^{(q_n)})_{c\geq0}$ (Propositions~\ref{prop:null-translation-Lqn-semigroup},\ref{prop:equal-dual-semigroups}), and the specialized Ward identity from Corollary~\ref{cor:cyclic-identity-logconvexity-form}. We also use technical results from Appendix~\ref{Sec:app-haagerupLp}.


\subsection{Null-translated sandwiched R\'enyi density}
\label{subsec:integer-density-propagation}
\begin{definition}[Sandwiched R\'{e}nyi density]\label{defn:renyi-density}
Let $\psi\in M_*^+$ be non-zero and assume
\begin{equation}
  Q_n(\psi\|\omega;M)<\infty\,.
\end{equation}
By the Haagerup--Kosaki characterization of SRD (Definition~\ref{def:SRD} and Lemma~\ref{lem:y-op-defn-SRD}), there exists a unique element $x_0\in L^n(M)_+$, which we call the sandwiched R\'enyi density, 
such that
\begin{equation}
  h_\psi=h^{\eta_n}x_0h^{\eta_n}\,, \quad\text{ and }\quad Q_n(\psi\|\omega;M)=\tr(x_0^n)\,,
\end{equation}
with $\eta_n= \frac{n-1}{2n}=\frac1{2q_n}$. 
For $c\geq 0$, define
\begin{equation}
  x_c \defeq V_c^{(n)}x_0\,.
\end{equation}
Since $V_c^{(n)}$ is positive and contractive,
\begin{equation}
  x_c\in L^n(M)_+\,,
  \qquad
  \|x_c\|_n\leq \|x_0\|_n\,.
\end{equation}
\end{definition}
Below, we also show that every $x_0\in L^n(M)_+$ corresponds to a unique non-zero $\psi\in M_*^+$ with finite SRD with respect to $\omega$.
\begin{lemma}[Positive \(L^n\)-elements as sandwiched R\'enyi densities]
\label{lem:positive-Ln-sandwiched-density}
For every non-zero \(x_0\in L^n(M)_+\), there exists a unique
non-zero normal positive functional \(\psi\in M_*^+\) such that
\begin{equation}
    h_\psi=h^{\eta_n}x_0h^{\eta_n}.
    \label{eq:functional-from-positive-Ln}
\end{equation}
Moreover,
\[
    0<Q_n(\psi\|\omega;M)<\infty,
\]
and \(x_0\) is the unique sandwiched Rényi density of \(\psi\)
relative to \(\omega\). In particular,
\begin{equation}
    Q_n(\psi\|\omega;M)
      =\|x_0\|_n^n
      =\tr(x_0^n).
    \label{eq:Qn-positive-Ln}
\end{equation}
\end{lemma}
\begin{proof}
By Theorem~\ref{thm:haagerup-kosaki},
\[
 L^n(M,\omega)
   =h^{1/(2q_n)}L^n(M)h^{1/(2q_n)}
   =h^{\eta_n}L^n(M)h^{\eta_n}
   \subset L^1(M).
\]
Consequently,
\[
    y:=h^{\eta_n}x_0h^{\eta_n}
\]
is a well-defined element of \(L^1(M)_+\). Moreover, the same
identification is isometric:
\[
    \|y\|_{n,\omega}=\|x_0\|_n.
\]
Since \(x_0\neq0\), it follows that \(\|y\|_{n,\omega}>0\), and
hence \(y\neq0\).

By Proposition~\ref{prop:standard-Haagerup-facts}, the map
\[
    M_*\ni\varphi\longmapsto h_\varphi\in L^1(M)
\]
is an isometric linear bijection that maps \(M_*^+\) onto
\(L^1(M)_+\). Therefore, there exists a unique
\(\psi\in M_*^+\) such that
\[
    h_\psi=y=h^{\eta_n}x_0h^{\eta_n}.
\]
Since \(y\neq0\), this functional is non-zero.

Finally, \(h_\psi\in L^n(M,\omega)\). The characterization of the
sandwiched R\'enyi divergence in Lemma~\ref{lem:y-op-defn-SRD} therefore
implies that
\[
    0<Q_n(\psi\Vert\omega;M)<\infty
\]
and that there is a unique element \(x\in L^n(M)_+\) satisfying
\[
    h_\psi=h^{\eta_n}xh^{\eta_n}.
\]
Equation~\eqref{eq:functional-from-positive-Ln} shows that this
unique element is \(x=x_0\). Thus \(x_0\) is the sandwiched Rényi
density of \(\psi\), and
\[
    Q_n(\psi\Vert\omega;M)
      =\|x_0\|_n^n
      =\tr(x_0^n).
\]
\end{proof}

The null translations and the sandwiching operation are compatible.
\begin{lemma}[Null translation of sandwiched R\'enyi density]
\label{lem:density-propagation-fixed}
Consider a non-zero $\psi\in M_*^+$ such that $Q_n(\psi\|\w;M)<\I$. 
Let
\begin{equation}
    \psi_c \defeq \psi\circ T_c
\end{equation}
be the null-translated functional on the fixed algebra $M$.  Then
\begin{equation}
    h_{\psi_c}=h^{\eta_n}x_ch^{\eta_n},
  \qquad x_c=V_c^{(n)}x_0\,.
\end{equation}
With $Q_n(c)\defeq Q_n(\psi_c\|\omega;M)=\tr(x_c^n)$, we also have
\begin{equation}
    0< Q_n(\psi_c\|\omega;M)<\infty \qquad\A c\geq 0.
\end{equation}
\end{lemma}
\begin{proof}
For $a\in M$,
\begin{align*}
\begin{aligned}
  \tr\bigl(x_cj_{q_n}(a)\bigr)
  &=\tr\bigl(V_c^{(n)}x_0\,j_{q_n}(a)\bigr)=\tr\bigl(x_0V_c^{(q_n)}j_{q_n}(a)\bigr)=\tr\bigl(x_0j_{q_n}(T_c(a))\bigr) \\
  &=\tr\bigl(h_\psi T_c(a)\bigr)=\psi(T_c(a))=\psi_c(a)\,.
  \end{aligned}
\end{align*}
On the other hand, by H\"older multiplication and cyclicity from Lemma~\ref{lem:holder-products-trace-cyclicity},
\begin{equation*}
  \tr\bigl(x_cj_{q_n}(a)\bigr)
  =\tr\bigl(x_ch^{\eta_n}ah^{\eta_n}\bigr)
  =\tr\bigl(h^{\eta_n}x_ch^{\eta_n}a\bigr)\,.
\end{equation*}
Thus, $h^{\eta_n}x_ch^{\eta_n}$ represents the normal positive functional $\psi_c$.  Since the map $M_*\ni\varphi\mapsto h_\varphi\in L^1(M)$ is injective,
\begin{equation*}
  h_{\psi_c}=h^{\eta_n}x_ch^{\eta_n}\,.
\end{equation*}

The formula for $Q_n(c)$ follows from the uniqueness of the integer SRD density (Lemma~\ref{lem:y-op-defn-SRD}).  Finiteness follows from $x_c\in L^n(M)_+$. If \(Q_n(c)=0\), then
\(\|x_c\|_n=0\), hence \(x_c=0\). Therefore \(h_{\psi_c}=h^{\eta_n}x_ch^{\eta_n}=0\),
so \(\psi_c=0\) by the injectivity of \(M_*\ni\phi\mapsto h_\phi\in L^1(M)\).
But
\[
\psi_c(1)=\psi(T_c(1))=\psi(1)>0,
\]
a contradiction. Thus \(Q_n(c)>0\).
\end{proof}

\subsection{Log-convexity for smooth semigroup orbits}
\label{subsec:smooth-orbit-log-convexity-fixed}

Under a regularity assumption on $x_0$, Corollary~\ref{cor:cyclic-identity-logconvexity-form} and the Cauchy-Schwarz inequality in the $L^2(M)$ Hilbert space lead to a straightforward proof of log-convexity of $\tr((V_c^{(n)}x_0)^n)$. The regularity restriction is removed in Sec.~\ref{subsec:integer-renyi-qnec-final-fixed}.
\begin{proposition}[Smooth-orbit log-convexity]
\label{prop:smooth-orbit-log-convexity-fixed}
Let $x_0\in L^n(M)_+\cap\Dom(G_n^2)$, $x_0\neq0$. Set
\[
    x_c=V_c^{(n)}x_0, \quad Q_n(c)=\tr(x_c^n), \quad c\geq0.
\]
Then $Q_n(c)$ is log-convex in $c$ on $[0,\infty)$:
\begin{equation}
    Q_n((1-\theta)c_0+\theta c_1)
  \leq Q_n(c_0)^{1-\theta}Q_n(c_1)^\theta
\end{equation}
for all $c_0,c_1\geq 0$ and $0\leq \theta\leq 1$.
\end{proposition}

\begin{proof}
By Lemma~\ref{lem:positive-Ln-sandwiched-density} applied to $x_c$, it follows that $0<Q_n(c)<\I, \, \forall c\geq0$.

Since $x_0\in\Dom(G_n^2)$, it follows from Definition~\ref{def:Gn-def} that
\begin{equation*}
     x_c\in\Dom(G_n^2),
  \qquad
  \frac{d}{dc}x_c=G_nx_c,
  \qquad
  \frac{d^2}{dc^2}x_c=G_n^2x_c\,, \quad \forall c\geq0
\end{equation*}
Fix $c> 0$ and write
\begin{equation}\label{eq:defn-xuv}
  x \defeq x_c,
  \qquad
  u \defeq G_nx_c,
  \qquad
  v \defeq G_n^2x_c\,.
\end{equation}
Because $x_c=x_c^*$ for all $c$ and the involution is an isometry on $L^n(M)$, differentiation in $L^n(M)$ gives $u=u^*,v=v^*$.

The product map
\begin{equation*}
  L^n(M)\times\cdots\times L^n(M)\longrightarrow L^1(M)
\end{equation*}
with $n$ factors is continuous by H\"older's inequality.  Since $c\mapsto x_c$ is twice continuously differentiable in $L^n(M)$ by assumption, the finite Leibniz rule gives
\begin{equation}\label{eq:Qp-C-defn}
    Q_n'(c)=nC, \quad  C \defeq \tr(ux^{n-1})\,.
\end{equation}
The second derivative is
\begin{equation}\label{eq:Qpp-regular}
  Q_n''(c)=n\left(A+\sum_{j=0}^{n-2}B_j\right)\,,
\end{equation}
where
\begin{equation}
  A \defeq \tr(vx^{n-1}),\quad  B_j \defeq \tr(ux^jux^{n-2-j}),
  \qquad j=0,\ldots,n-2\,.
\end{equation}

Each $B_j$ is non-negative.  To see this, define
\begin{equation}\label{eq:defn-Yj}
  Y_j \defeq x^{j/2}ux^{(n-2-j)/2}\in L^2(M)\,.
\end{equation}
Using Lemma~\ref{lem:holder-products-trace-cyclicity}, H\"older multiplication gives $Y_j\in L^2(M)$, and cyclicity gives
\begin{equation*}
  \|Y_j\|_2^2=\tr(Y_j^*Y_j)=\tr\bigl(x^{(n-2-j)/2}ux^jux^{(n-2-j)/2}\bigr)=B_j \geq0\,.
\end{equation*}
Moreover,
\begin{equation*}
  B_j=B_{n-2-j}\,.
\end{equation*}
This follows from cyclicity and self-adjointness:
\begin{align*}
\begin{aligned}
    B_j &=\tr(ux^jux^{n-2-j})=\tr\bigl((ux^jux^{n-2-j})^*\bigr)=\tr(x^{n-2-j}ux^ju)=\tr(ux^{n-2-j}ux^j)=B_{n-2-j}\,.
\end{aligned}
\end{align*}

Since $x\in\Dom(G_n^2)$ by assumption, and $u=G_nx\in\Dom(G_n)$, we apply Corollary~\ref{cor:cyclic-identity-logconvexity-form} with $x=x_c$ to \eqref{eq:Qpp-regular}, which gives
\begin{equation}\label{eq:Qypp}
  \frac{Q_n''(c)}{n}
  =\sum_{j=0}^{n-2}\bigl(1-\zeta_n^{j+1}\bigr)B_j\,.
\end{equation}
Using $B_j=B_{n-2-j}$, one can write this as 
\begin{equation}\label{eq:Qpp-positive-sum}
    \frac{Q_n''(c)}{n} = \sum_{j=0}^{n-2}w_jB_j, \qquad w_j=2\sin^2\bfrac{\pi(j+1)}{n},
\end{equation}
with
\begin{equation*}
    w_j\geq0,\qquad \sum_{j=0}^{n-2}w_j=n.
\end{equation*}

Let
\begin{equation*}
  Z \defeq x^{n/2}\in L^2(M)\,.
\end{equation*}
Then for $j=0,\ldots,n-2$\,, with $Y_j$ defined in \eqref{eq:defn-Yj}, we have
\begin{equation*}
  \|Y_j\|_2^2=B_j,
  \qquad
  \|Z\|_2^2=Q_n(c)\,.
\end{equation*}
Furthermore, by cyclicity,
\begin{equation*}
  \langle Y_j,Z\rangle_{L^2(M)}=\tr(Y_j^*Z) =\tr\bigl(x^{(n-2-j)/2}ux^{j/2}x^{n/2}\bigr)=C\,.
\end{equation*}
where $C$ was defined in \eqref{eq:Qp-C-defn}. 
Cauchy-Schwarz in $L^2(M)$ gives
\begin{equation}\label{eq:C-inequality}
  |C|^2\leq Q_n(c)B_j,
  \qquad j=0,\ldots,n-2\,.
\end{equation}
Also, $C$ is real: using $u=u^*$, $x=x^*$, and cyclicity,
\begin{equation*}
  \overline C
  =\tr\bigl((ux^{n-1})^*\bigr)
  =\tr(x^{n-1}u)
  =\tr(ux^{n-1})
  =C\,.
\end{equation*}
Thus $|C|^2=C^2$.

Then, using \eqref{eq:Qpp-positive-sum}, \eqref{eq:C-inequality}, and the fact that $C$ is real,
\begin{equation*}
  Q_n(c)Q_n''(c)=nQ_n(c)\sum_{j=0}^{n-2} w_j B_j \geq n\sum_{j=0}^{n-2} w_j C^2 =n^2C^2 = Q_n'(c)^2\,.
\end{equation*}
Since $Q_n(c)>0$,
\begin{equation}
  (\log Q_n)''(c)
  =\frac{Q_n(c)Q_n''(c)-Q_n'(c)^2}{Q_n(c)^2}
  \geq 0\,.
\end{equation}
Hence $\log Q_n$ is convex on $(0,\infty)$, which is equivalent to the claimed log-convexity inequality for $c\in(0,\infty)$.

Since $c\mapsto x_c$ is strongly continuous in $L^n(M)$, from the strong continuity of the semigroup $(V_c^{(n)})_{c\geq0}$ (Proposition~\ref{prop:null-translation-Ln-semigroup}), and since the product map $L^n(M)^n\to L^1(M)$ is continuous by H\"{o}lder, $Q_n(c)=\tr(x_c^n)$ is continuous on $[0,\I)$. Since $Q_n(c)>0$, $\log Q_n$ is continuous on $[0,\I)$. Convexity on $(0,\I)$ then extends to convexity on $[0,\I)$ by taking limits from the right.
\end{proof}


\subsection{Log convexity for general positive \texorpdfstring{$L^n$}{Ln} data}
\label{subsec:integer-renyi-qnec-final-fixed}

We now relax the assumption $x_0\in \Dom(G_n^2)$ in Proposition~\ref{prop:smooth-orbit-log-convexity-fixed} and prove log convexity for arbitrary $x_0\in L^n(M)_+$. This is done by approximating $x_0$ by a sequence in $L^n(M)_+\cap\Dom(G_n^2)$, so that Proposition~\ref{prop:smooth-orbit-log-convexity-fixed} applies as we take the limit.

We first collect the definition and some properties of the Yosida resolvent of the generator $G_n$ of the strongly continuous contractive semigroup $(V_c^{(n)})_{c\geq0}$ on the $L^n(M)$ Banach space. See Appendix \ref{appsubsec:semigrp} for more details.
\begin{definition}[Yosida resolvent]\label{def:yosidaresolvent-main}
    Recall the definition of the resolvent of $G_n$, $R(m,G_n)$, from Remark~\ref{rem:resolvent}. The Yosida resolvent is defined to be
\begin{equation}
    R_m=mR(m,G_n)=m(m-G_n)^{-1}, \quad m>0,
\end{equation}
which admits the Laplace-transform representation
\begin{equation}
    R_m= \lim_{t\to\I} m\int_0^t e^{-m s}V_s^{(n)}\,ds.
\end{equation}
\end{definition}

\begin{lemma}\label{lem:yosidaproperties}
    Let the $R_m, m>0$ be the Yosida resolvent of the generator $G_n$ of the strongly continuous contractive semigroup $(V_c^{(n)})_{c\geq0}$, as in Definition~\ref{def:yosidaresolvent-main}. Then the following properties hold:
    \begin{enumerate}
        \item $R_m:L^n(M) \longmapsto \Dom(G_n)$ is a positive operator.
        \item For $x_0\in L^n(M)_+$, $R_m^2 x_0\in L^n(M)_+\cap \Dom(G_n^2)$.
        \item $R_m V_c^{(n)}=V_c^{(n)} R_m$.
        \item For $x_0\in L^n(M)_+$, $R_m^2x_0 \to x_0 $ as $m\to \infty$.
    \end{enumerate}
\end{lemma}
\begin{proof}
We prove statements (1), (2) and (4). Statement (3) follows from Lemma~\ref{lem:resolvent-convergence}.

\emph{Proof of (1):}
Since \(m\in\rho(G_n)\), the operator
\[
m-G_n:\Dom(G_n)\longrightarrow L^n(M)
\]
is bijective. Consequently,
\[
R_m=m(m-G_n)^{-1}
\]
maps \(L^n(M)\) into \(\Dom(G_n)\).
It remains to prove positivity. For \(x\in L^n(M)_+\), the Laplace-transform
representation gives
\[
R_mx
   =\lim_{t\to\infty}
      m\int_0^t e^{-ms}V_s^{(n)}x\,ds,
\]
where the integrals are Bochner integrals in \(L^n(M)\). Since
\(V_s^{(n)}\) is positive,
\[
m e^{-ms}V_s^{(n)}x\in L^n(M)_+
\qquad (s\geq0).
\]
The positive cone \(L^n(M)_+\) is norm closed and convex. Hence every
finite Bochner integral
\[
m\int_0^t e^{-ms}V_s^{(n)}x\,ds
\]
belongs to \(L^n(M)_+\), and so does its \(L^n(M)\)-norm limit as
\(t\to\infty\). Thus
\[
x\geq0\quad\Longrightarrow\quad R_mx\geq0,
\]
and \(R_m\) is positive. This proves statement (1).

\emph{Proof of (2):}
Since \(R_m\) is positive,
\[
R_mx_0\in L^n(M)_+,
\qquad
R_m^2x_0=R_m(R_mx_0)\in L^n(M)_+.
\]
It remains to prove that \(R_m^2x_0\in\Dom(G_n^2)\).

It follows from statement (1) that
\[
R_m^2x_0= R_m(R_mx_0)\in\Dom(G_n).
\]
We  use the standard resolvent commutation identity (Lemma~\ref{lem:resolvent-convergence})
\[
G_nR_my=R_mG_ny,
\qquad y\in\Dom(G_n).
\]
Applying this identity to \(y=R_mx_0\in\Dom(G_n)\), we obtain
\[
G_nR_m^2x_0
   =R_mG_nR_mx_0 \in \Dom(G_n).
\]
Therefore,
\[
R_m^2x_0\in\Dom(G_n)
\quad\text{and}\quad
G_nR_m^2x_0\in\Dom(G_n),
\]
which, by definition, means that
\[
R_m^2x_0\in\Dom(G_n^2).
\]
Combining this with positivity proves
\[
R_m^2x_0\in L^n(M)_+\cap\Dom(G_n^2).
\]

\emph{Proof of (4).}
Since \(\|R_m\|\leq1\) by Theorem~\ref{thm:hille-yosida}, and \(R_m x_0\to x_0\) by Lemma~\ref{lem:resolvent-convergence} , we have
\[
 \|R_m^2x-x\|_n
 \leq
 \|R_m(R_mx-x)\|_n+\|R_mx-x\|_n
 \leq
 2\|R_mx-x\|_n
 \longrightarrow0.
\]
\end{proof}


\begin{theorem}[Log-convexity of $Q_n(c)$ for general functionals]
\label{thm:general-log-convexity}
Let $x_0\in L^n(M)_+$, $x_0\neq0$.
Set
\[
    x_c=V_c^{(n)}x_0, \quad Q_n(c)=\tr(x_c^n), \quad c\geq0.
\]
Then $Q_n$ is log-convex on $[0,\infty)$:
\begin{equation}
    Q_n((1-\theta)c_0+\theta c_1)
  \leq Q_n(c_0)^{1-\theta}Q_n(c_1)^\theta
\end{equation}
for all $c_0,c_1\geq 0$ and $0\leq \theta\leq 1$.
\end{theorem}
\begin{proof}
For $m>0$, define the Yosida resolvent $R_m$ of $G_n$, as in Definition~\ref{def:yosidaresolvent-main}, and the Yosida approximants 
\begin{equation*}
  x_0^{(m)} \defeq R_m^2x_0\,,
\end{equation*}
From Lemma~\ref{lem:yosidaproperties} it follows that 
\begin{equation*}
  x_0^{(m)}\in L^n(M)_+\cap\Dom(G_n^2),
\end{equation*}
and
\begin{equation}
 \lim_{m\to \I} x_0^{(m)}\to x_0
  \quad
  \text{in }L^n(M).
\end{equation}
Moreover, $x_0^{(m)}\neq 0$, since $x_0\neq0$ by assumption and resolvents are injective.

For $c\geq 0$, set
\begin{equation*}
  x_c^{(m)} \defeq V_c^{(n)}x_0^{(m)}\,.
\end{equation*}
Then by Lemma~\ref{lem:yosidaproperties}
\begin{equation*}
    x_c^{(m)}=V_c^{(n)}x_0^{(m)}=R_m^2 V_c^{(n)}x_0\in L^n(M)_+\cap\Dom(G_n^2).
\end{equation*}
Define
\begin{equation*}
  Q_{n,m}(c) \defeq \tr\bigl((x_c^{(m)})^n\bigr)\,.
\end{equation*}
By Lemmas~\ref{lem:positive-Ln-sandwiched-density} and \ref{lem:density-propagation-fixed} applied to $x_0^{(m)}$, it follows that $ 0<Q_{n,m}(c)<\I$. 
Thus, Proposition~\ref{prop:smooth-orbit-log-convexity-fixed} applies to each $Q_{n,m}(c)$, implying that each $Q_{n,m}$ is positive and log-convex:
\begin{equation*}
  Q_{n,m}((1-\theta)c_0+\theta c_1)
  \leq Q_{n,m}(c_0)^{1-\theta}Q_{n,m}(c_1)^\theta
\end{equation*}
for all $c_0,c_1\geq 0$ and $0\leq\theta\leq 1$.

For each fixed $c\geq 0$,
\begin{equation*}
  x_c^{(m)}=V_c^{(n)}R_m^2x_0 =R_m^2V_c^{(n)}x_0 \longrightarrow V_c^{(n)}x_0=x_c \qquad\text{in }L^n(M).
\end{equation*}
By Lemma \ref{lem:integer-power-continuity}, the convergence
\(x_c^{(m)}\to x_c\) in \(L^n(M)_+\) implies
\[
(x_c^{(m)})^n \to x_c^n
\]
in \(L^1(M)\).
Hence, for every $c\geq 0$,
\begin{equation*}
  Q_{n,m}(c)\longrightarrow Q_n(c) \defeq \tr(x_c^n).
\end{equation*}
The limit $Q_n(c)$ is positive by Lemmas~\ref{lem:positive-Ln-sandwiched-density} and \ref{lem:density-propagation-fixed} applied to $x_0$. 
Passing to the limit in the log-convexity inequality for $Q_{n,m}$ gives
\begin{equation*}
  Q_n((1-\theta)c_0+\theta c_1)
  \leq Q_n(c_0)^{1-\theta}Q_n(c_1)^\theta\,.
\end{equation*}
Thus $Q_n$ is log-convex on $[0,\infty)$.
\end{proof}




The integer R\'enyi QNEC follows as a corollary of the above theorem.
\begin{corollary}[Integer R\'enyi QNEC]
\label{cor:integer-renyi-qnec-fixed} 
Suppose that the HSMI \(M_B\subset M_A\) is realized by null-cut
algebras in a Poincar\'e-invariant quantum field theory, with common faithful normal vacuum state \(\omega\). 
Assume the HSMI structure from Theorem~\ref{thm:modulartranslations}, and the fixed algebra formulation of the SRD from Lemma~\ref{lem:fixedalgebraSRD}.  Let $\psi\in M_*^+$ be non-zero and assume
\[
  D_n(\psi\|\omega;M)<\infty\,.
\]
For $c\geq 0$, define
\[
  \psi_c \defeq \psi\circ T_c\,.
\]
Then
\[
  0< Q_n(\psi_c\|\omega;M)<\infty
\]
for all $c\geq 0$, and the function
\[
  c\longmapsto S_n(c)=D_n(\psi_c\|\omega;M)=D_n(\psi\|\omega;M_c)
\]
 is convex on $[0,\infty)$.  Equivalently,
 \[
    S_n((1-\theta)c_0+\theta c_1)
  \leq (1-\theta)S_n(c_0)
     +\theta S_n(c_1)
 \]
for all $c_0,c_1\geq 0$ and $0\leq\theta\leq 1$.  Consequently,
\[
  \partial_c^2S_n(c)\geq 0
\]
 in the sense of distributions on $(0,\infty)$, and pointwise wherever the ordinary second derivative exists.
\end{corollary}
\begin{proof}
By the finite SRD assumption, there exists a unique
\begin{equation*}
  x_0\in L^n(M)_+
\end{equation*}
with
\begin{equation*}
  h_\psi=h^{\eta_n}x_0h^{\eta_n},
  \qquad
  Q_n(\psi\|\omega;M)=\tr(x_0^n)\,.
\end{equation*}
For $c\geq 0$, Lemma~\ref{lem:density-propagation-fixed} gives
\begin{equation*}
  h_{\psi_c}=h^{\eta_n}x_ch^{\eta_n},
  \qquad
  x_c=V_c^{(n)}x_0\,.
\end{equation*}
and 
\begin{equation*}
  0< Q_n(\psi_c\|\omega;M)=\tr(x_c^n)<\infty\,.
\end{equation*}
By the Yosida regularization argument in Theorem \ref{thm:general-log-convexity}, the function
\begin{equation*}
  Q_n(c) \defeq Q_n(\psi_c\|\omega;M)=\tr(x_c^n)
\end{equation*}
 is log-convex on $[0,\infty)$:
\begin{equation*}
  Q_n((1-\theta)c_0+\theta c_1)
  \leq Q_n(c_0)^{1-\theta}Q_n(c_1)^\theta\,.
\end{equation*}
Taking logarithms gives convexity of $\log Q_n(c)$.  Since
\begin{equation*}
  D_n(\psi_c\|\omega;M)
  =\frac{1}{n-1}\log\left(\frac{Q_n(\psi_c\|\omega;M)}{\psi_c(1)}\right)
  =\frac{1}{n-1}\log\left(\frac{Q_n(c)}{\psi(1)}\right)\,,
\end{equation*}
where $\psi_c(1)=\psi(T_c(1))=\psi(1)$ is independent of $c$, convexity of $D_n$ follows.  Convexity implies non-negativity of the second derivative in the distributional sense, and the pointwise statement holds at every point where the ordinary second derivative exists.
\end{proof}



%%%%%%%%%%%%%%%%%%%%%%%%%%%%%%%%%%%%%%%%%%%%%%%%%%%%%%%%%%%%%%%%%%%%%%%%%%%%%%%%%%%%%%%%%%%%%%%%%%%%

\section{Discussion and future directions}\label{Sec:conclusion}

We have proven the integer RQNEC for arbitrary $\sigma$-finite von Neumann algebras carrying a half-sided modular inclusion structure, under the single additional hypothesis that the SRD with respect to the distinguished faithful state $\omega$ is finite on the larger algebra $M_A$.  
In particular, no form-domain assumption $\Psi\in \Dom(P^{1/2})$ is required, in contrast to the proofs of QNEC \cite{Ceyhan:2018zfg,Hollands:2025glm}.
Whether QNEC itself can be proven without such a form-domain assumption remains an interesting open problem.

An immediate target for future work is to prove RQNEC for general real  R\'enyi parameter $\a>1$. We expect that the proof technique used here can be adapted to prove this statement, and we hope to present the details in a future publication.

In the remainder of this section, we discuss several directions for future research.

\subsection*{RQNEC for semigroups}

The integer R\'enyi QNEC proved here is a log-convexity statement for the semigroup $V_c^{(n)}$ on $L^n(M)$: writing $Q_n(c)=\|x_c\|_n^n$ for the propagated density $x_c=V_c^{(n)}x_0$, our result is the log-convexity inequality
\[
  Q_n((1-\theta)c_0+\theta c_1)\leq Q_n(c_0)^{1-\theta}Q_n(c_1)^\theta\,.
\]
This is a noncommutative functional inequality for the semigroup $V_c^{(n)}$, in the spirit of the log-Sobolev and Poincar\'e inequalities developed in \cite{olkiewicz1999hypercontractivity, junge2015noncommutative, carlen2017gradient} and others.
Hollands and Longo~\cite{Hollands:2025glm} noted the connection between QNEC-type inequalities and noncommutative entropy production under semigroup actions, in the sense of~\cite{Spohn:1978pfz, Wirth:2025goz}. 
It would be interesting to investigate whether QNEC and RQNEC-type inequalities exist for semigroups arising from other geometric flows in quantum field theory. Similar proof techniques to those in the present work may apply.

\subsection*{Further generalizations of R\'enyi QNEC}
A family of divergences known as optimized quantum $f$-divergences has been studied for finite-dimensional systems \cite{Wilde:2017okz}. These divergences satisfy the data processing inequality, and the SRD is one example of this family. It would be interesting to study these divergences in the von Neumann algebra setting for general QFTs, and prove convexity statements similar to RQNEC. 
It is also natural to conjecture the {\az} R\'enyi QNEC for the {\az} R\'enyi divergence \cite{Audenaert:2015npv}, a two-parameter generalization of the relative entropy. It has been formulated in terms of Haagerup $L^p$ spaces \cite{Kato:2023aro,Kato:2023hlj,Hiai:2024qve}, and investigating its log-convexity under null translations will also be interesting.


\subsection*{R\'enyi quantum focusing}

Recently, \cite{Chandrasekaran:2026pnc} proved the quantum focusing conjecture of \cite{Bousso:2015mna} at leading non-trivial order in perturbative quantum gravity on Killing horizon backgrounds, building on the proof of QNEC \cite{Ceyhan:2018zfg}. For a one-parameter family of horizon cuts labelled by affine parameter $u$, the quantum expansion is defined as the derivative of the generalized entropy,
\begin{equation}
    \Theta(u) \defeq \del_u\barSgen(u;\hat\psi)\,,
\end{equation} 
evaluated on a dressed state $\ket{\hat\psi}=\int dx\,f(x)\ket{\Psi}\ket{x}$ that pairs an excited bulk state $\ket{\Psi}$ with eigenstates $\ket{x}$ of the area operator. The quantum focusing conjecture asserts that $\Theta$ is non-increasing in $u$, 
\begin{equation}        \label{eq:expansion-monotonicity}
    \del_u \Theta(u)\leq 0\,.
\end{equation}
Equivalently, it is the statement that the generalized entropy is concave.

Following \cite{Chandrasekaran:2022eqq}, the authors of \cite{Chandrasekaran:2026pnc} show that up to perturbative corrections in $G_N$,
\begin{equation}\label{eq:rel-ent-gravitational}
    \barSgen(u;\hat\psi) \approx -S_{\text{rel}}(u) + \text{const.}, \qquad S_{\text{rel}}(u) = -\langle\hat\psi,\log\D_{\hat\psi,\hat\W;\wh{\cM}_u}\hat\psi\rangle\,,
\end{equation}
where $\hat\W$ is the cyclic and separating Hartle--Hawking state and $\widehat{\cM}_u$ is the crossed product of the bulk QFT algebra by its vacuum modular automorphism group. The algebras $\wh{\cM}_u$ form a gravitational half-sided modular inclusion. The proof of \cite{Chandrasekaran:2026pnc} then follows \cite{Ceyhan:2018zfg} to derive a variational expression for $\Theta$ and deduce its monotonicity from convexity of relative entropy.

It would be interesting to investigate whether the SRD and our proof of RQNEC can be used to establish an analogous R\'enyi quantum focusing theorem.  We define the R\'enyi quantum expansion as
\begin{equation}
    \Theta_\a(u) = -\del_u D_\a(\hat\psi\|\hat\W;\wh{\cM}_u)\,,
\end{equation}
and conjecture that the R\'enyi quantum expansion is non-increasing in $u$
\begin{equation}
    \del_u\Theta_\a(u)\leq 0,\quad \forall \a>1\,.
\end{equation} 
Several alternative definitions are possible, reflecting the various R\'enyi generalizations of relative entropy, such as the Petz-R\'enyi divergence. Our choice of the SRD is well-motivated from the perspective that the R\'enyi quantum expansion (and R\'enyi focusing) should have a well-defined limit as $G_N\to0$.

%%%%%%%%%%%%%%%%%%%%%%%%%%%%%%%%%%%%%%%%%%%%%%%%%%%%%%%%%%%%%%%%%%%%%%%%%%%%%%%%%%%%%%%%%%%%%%%%%%%%

\begin{acknowledgments}
We thank Krishna Jalan for comments on the manuscript.
T.K. is supported by a Simons Foundation fellowship through the Targeted Grant to Instituto Balseiro. The work of P.R. has been supported by the Polish National Science Centre through Sonata grant (2022/47/D/ST2/02058). 
During the preparation of this manuscript, the authors used Anthropic Claude and OpenAI ChatGPT as editorial aids in drafting and revising portions of the text. Every mathematical statement, proof, and citation in the final manuscript was independently checked by the authors, who take full responsibility for its content.
\end{acknowledgments}

\begin{appendix}

\section{Analysis on Banach spaces}\label{appsec:technical}
This appendix records the analytic facts used in Sections~\ref{sec:cyclic-ward}
and~\ref{sec:integer-renyi-qnec}.  We avoid recalling standard Banach-space
definitions, except where a convention is needed for the proof. For background on one-parameter semigroups and vector-valued integration we refer to
\cite{engel.nagel:00:one-parameter,arendt.batty.ea:11:vector-valued}.

\subsection{Ultraweak and Bochner integrals}

Let $M$ be a $\sigma$-finite von Neumann algebra with a faithful normal state $\omega\in M_*^+$, and modular automorphisms $\sigma^{\omega}_t$.
\begin{definition}[Ultraweak integral]\label{def:ultraweak-int}
Let $F: \mathbb R \to M$ be an operator--valued function. The function $F$ is said to be ultraweakly integrable with respect to a measure $\mu(t)$ if the scalar function $t\mapsto \varphi(F(t))$ is Lebesgue integrable with respect to $\mu(t)$ for all $\varphi\in M_*$.
\end{definition}

The Bochner integral provides a theory of Lebesgue integration for vector-valued functions. Our overview closely follows that in \cite{engel.nagel:00:one-parameter}. Consider vector-valued functions $f:I\subset \R\to X$ for $X$ a Banach space.
\begin{definition}[Simple function]
The function $f$ is said to be simple if it can be written as a finite sum of the form
\[
    f(x) = \sum_{k=1}^n\chi_{I_k}x_k
\]
where $\chi_{I_k}$ are characteristic functions for measurable subsets $I_k\subset I.$ The integral of a simple function is defined to be
\[
    \int_I f(s)ds\defeq \sum_{k=1}^nx_k\mu(I_k),
\]
where $\mu$ denotes the Lebesgue measure on $\R$.
\end{definition}
\begin{definition}[Strong measurability]
We say that $f$ is strongly measurable if there exists a sequence $f_n$ of simple functions such that
\[
    \lim_{n\to\I}\|f(s)-f_n(s)\|=0 \qquad\text{a.e.}
\]
\end{definition}
\begin{definition}[Bochner integral]\label{def:Bochner-int}
We say that $f$ is Bochner integrable if $f$ is strongly measurable, and there exists a sequence $f_n$ of simple functions on $I$ such that
\[
    \lim_{n\to\I}\int_I\|f(s)-f_n(s)\|\,ds=0.
\]
For a Bochner integrable function, we define its Bochner integral by
\[
    \int_I f(s)\,ds\defeq \lim_{n\to\I}\int_I f_n(s)\,ds,
\]
which is independent of the choice of the approximating simple functions.
\end{definition}


\begin{definition}[Approximate identity]\label{def:approx-id}
An approximate identity \(\mu_\varepsilon\) is a family of finite positive Borel measures on \(\mathbb R\) such that as $\ve\to0$,
\[
        \mu_\varepsilon(\mathbb R)\to 1,
        \qquad
        \mu_\varepsilon(\{|s|\geq \delta\})\to 0
        \quad\text{for every }\delta>0 .
\]
\end{definition}
\begin{lemma}\label{lem:approx-id}
Let \(X\) be a Banach space and let $\mu_\ve$ be an approximate identity. Then, for every bounded function \(F:\mathbb R\to X\) that is norm-continuous
at \(0\),
\[
        \int_{\mathbb R}F(s)\,d\mu_\varepsilon(s)
        \longrightarrow F(0)
        \qquad\text{in }X .
\]
\end{lemma}
\begin{proof}
We have
\[
\begin{aligned}
\left\|
        \int_{\mathbb R}F(s)\,d\mu_\varepsilon(s)-F(0)
\right\|
&\leq
        \int_{\mathbb R}\|F(s)-F(0)\|\,d\mu_\varepsilon(s)
        +
        |\mu_\varepsilon(\mathbb R)-1|\,\|F(0)\| .
\end{aligned}
\]
Given \(\eta>0\), choose \(\delta>0\) such that
\(\|F(s)-F(0)\|<\eta\) for \(|s|<\delta\).  If
\(\|F\|_\infty\leq C\), then
\[
        \int_{\mathbb R}\|F(s)-F(0)\|\,d\mu_\varepsilon(s)
        \leq
        \eta\,\mu_\varepsilon(\mathbb R)
        +
        2C\,\mu_\varepsilon(\{|s|\geq\delta\}),
\]
which tends to at most \(\eta\).  Since \(\eta\) is arbitrary, the claim
follows.
\end{proof}

\subsection{One-parameter semigroups}\label{appsubsec:semigrp}
We provide a concise summary of the relevant aspects of one-parameter semigroups on Banach spaces. We refer the reader to \cite{engel.nagel:00:one-parameter,arendt.batty.ea:11:vector-valued,hille1981functional} for more details. Our overview is based on \cite{engel.nagel:00:one-parameter}.

\begin{definition}[Strongly continuous semigroup]
A one-parameter family $(T(t))_{t\geq0}$ of bounded linear operators on a Banach space $X$ is called a strongly continuous semigroup if
\begin{equation}
    T(0) = I, \qquad \text{and}\qquad T(t+s)=T(t)T(s)\quad\A t,s\geq0,
\end{equation}
and the map $\R_+\ni t\mapsto T(t)\in B(X)$
is continuous w.r.t.~the strong operator topology on $X$. 
Here $B(X)$ is the Banach algebra of all bounded linear operators on $X$ with the operator norm $\|\cdot\|$. Define the orbit maps of the semigroup to be $\xi_x:t\mapsto T(t)x\in X$.
\end{definition}
\begin{proposition}\label{prop:semigroup}
For a semigroup $(T(t))_{t\geq0}$ on a Banach space $X$, the following are equivalent.
\begin{enumerate}
    \item $(T(t))_{t\geq0}$ is strongly continuous.
    \item $\lim_{t\downarrow0}T(t)x=x$ for all $x\in X$.
    \item There exist $\d>0,M\geq 1$, and a dense subset $D\subset X$ such that
    \begin{enumerate}
        \item $\|T(t)\|\leq M$ for all $t\in[0,\d]$,
        \item $\lim_{t\downarrow0}T(t)x=x$ for all $x\in D$.
    \end{enumerate}
\end{enumerate}
\end{proposition}


Strongly continuous semigroups are uniformly bounded on each compact interval.
\begin{definition}[Strongly continuous contractive semigroup]
A strongly continuous semigroup $(T(t))_{t\geq0}$ is said to be bounded if there exists $M\geq 1$ such that $\|T(t)\|\leq M$ for all $t\geq 0$. A bounded semigroup is called contractive if $M=1$.
\end{definition}
\begin{definition}[Generator of semigroup]\label{def:generator}
The generator $A:\Dom(A)\subseteq X\to X$ of a strongly continuous semigroup $(T(t))_{t\geq0}$ on a Banach space $X$ is the operator
\begin{equation}\label{eq:defn-semigp-gen}
    Ax \defeq \lim_{h\downarrow0}\frac{T(h)x-x}{h},
\end{equation}
defined for every $x$ in its domain
\begin{equation}\label{eq:defn-semigp-gen-dom}
    \Dom(A)\defeq \lbr x\in X:\lim_{h\downarrow0}\frac{T(h)x-x}{h} \text{ exists} \rbr.
\end{equation}
\end{definition}
The domain $\Dom(A)$, a linear subspace of $X$, is an essential part of the definition of the generator $A$. All generators $A$ should be understood to implicitly refer to the pair $(A,\Dom(A))$.
\begin{theorem}\label{thm:generator-is-closed}
The generator of a strongly continuous semigroup is a closed densely defined linear operator that determines the semigroup uniquely.
\end{theorem}
\begin{definition}[Operator graph and graph norm]\label{def:graph}
The graph of a linear operator $A$ with domain $\Dom(A)$ in a Banach space $X$, is defined to be the subset $\cG(A)\defeq\{(x,Ax):x\in \Dom(A)\}$. The graph norm associated with $A$ is defined to be
\begin{equation}\label{eq:graph-norm}
    \|x\|_{\cG(A)}\defeq \|x\|+\|Ax\|, \qquad x\in \Dom(A).
\end{equation}
\end{definition}\label{def:core}
\begin{definition}[Graph core of an operator]
A subspace $\cD$ of $\Dom(A)$ that is dense in $\Dom(A)$ for the graph norm is called a core of $A$.
\end{definition}

\begin{lemma}[Resolvent core criterion]\label{lem:resolvent-core}
    $\cD \subset\Dom(A)$ is a graph core of $A$ if and only if $(\l-A)\cD$ is dense in $X$ for some $\l\in\r(A)$.
\end{lemma}


\begin{definition}\label{def:resolvent-appendix}
For a (not necessarily densely defined) closed linear operator $A:\Dom(A)\subset X\to X$, the resolvent set of $A$ is
\begin{equation}
    \r(A)\defeq \{\l\in\C:\l-A \text{ is bijective as a map }\Dom(A)\to X \}.
\end{equation}
The complement of the resolvent set, $\s(A)\defeq\C\backslash\r(A)$ is called the spectrum of $A$. For $\l\in\r(A)$, the inverse
\begin{equation}
    R(\l,A)\defeq (\l-A)^{-1},
\end{equation}
is a bounded operator on $X$ called the resolvent of $A$ (at point $\l$).
\end{definition}
The resolvent of the generator of a strongly continuous contractive semigroup has the following integral representation,
\begin{equation}
    R(\l,A)x = \lim_{t\to\I}\int_0^t e^{-\l s}T(s)x\,ds, \qquad \A x\in X,
\end{equation}
which we write as
\begin{equation}
    R(\l,A) = \int_0^\I e^{-\l s}T(s)\,ds.
\end{equation}
This relation can be heuristically understood as the statement that the resolvent is the Laplace transform of $T(t)$. Schematically, $T(t)x=e^{tA}x$, and the resolvent definition $R(\l,A)x= (\l-A)^{-1}x$ can be thought of as the Laplace transform $(\l-\a)^{-1}= \int_0^\I e^{-\l s}e^{\a s}ds$. \cite{hille1981functional}
\begin{theorem}[Hille-Yosida Generation Theorem]\label{thm:hille-yosida}
For a linear operator $(A,\Dom(A))$ on a Banach space $X$, the following are equivalent.
\begin{enumerate}
    \item $(A,\Dom(A))$ generates a strongly continuous contraction semigroup.
    \item $(A,\Dom(A))$ is closed, densely defined, and for all $\l>0$, one has $\l\in\r(A)$ and $\|\l R(\l,A)\|\leq 1.$
    \item $(A,\Dom(A))$ is closed, densely defined, and for all $\l\in\C$ with $\mathrm{Re}\,\l>0$, one has $\l\in\r(A)$ and $\|R(\l,A)\|\leq (\mathrm{Re}\,\l)^{-1}$.
\end{enumerate}
\end{theorem}
\begin{lemma}\label{lem:resolvent-convergence}
Let $(A,\Dom(A))$ be the generator of a strongly continuous contraction semigroup $(T(s))_{s\geq0}$ on a Banach space $X$. Then, the following statements hold.
\begin{enumerate}
    \item $\lim_{\lambda\to\I}\l R(\l,A)x\to x $ for all $x\in X$.
    \item $T(s)R(\lambda,A)=R(\lambda,A)T(s)$ for all $s\geq0$, and $\lambda\in \rho(A)$. In particular, by Theorem~\ref{thm:hille-yosida}(2), this holds for all $\lambda>0$
    \item $\l AR(\l,A)x=\l R(\l,A)Ax\to Ax $ for all $x\in \Dom(A)$, and $\lambda\in \rho(A)$.
\end{enumerate}
\end{lemma}
\begin{definition}[Yosida resolvents and approximants]
Let $(A,\Dom(A))$ be the generator of a strongly continuous contraction semigroup on a Banach space $X$. We define the Yosida resolvents of $A$ to be
\begin{equation}
    R_m \defeq m(m-A)^{-1} = mR(m,A).
\end{equation}
The Yosida approximants of the generator $A$ are defined to be 
\begin{equation}
    A_m = R_mA.
\end{equation}
For general $x\in X$, we define Yosida approximants $x_m$ to be
\begin{equation}
    x_m \defeq R_mx.
\end{equation}
Note that by Lemma \ref{lem:resolvent-convergence}, we have
\begin{equation}
    A_mx\to Ax,\qquad\text{for }x\in \Dom(A).
\end{equation}
\end{definition}


\section{Analytical results for Haagerup \texorpdfstring{$L^p$}{Lp} spaces}\label{Sec:app-haagerupLp}

Throughout this subsection we assume the same notations and definitions as in Sec.~\ref{Sec:prelim}. \(L^p(M)\) denotes Haagerup's noncommutative
\(L^p\)-space and \(\operatorname{tr}:L^1(M)\to\mathbb C\) denotes the canonical functional on Haagerup's $L^1$ space.  Products are products of affiliated measurable operators.
We use the convention \(1/\infty=0\).

\begin{lemma}[H\"older products and cyclicity]\label{lem:holder-products-trace-cyclicity}
Let \(1\le p_1,\ldots,p_k\le\infty\) and suppose
\[
        \frac1r=\frac1{p_1}+\cdots+\frac1{p_k}\le1 .
\]
If \(a_j\in L^{p_j}(M)\), then
\[
        a_1a_2\cdots a_k\in L^r(M),
        \qquad
        \|a_1\cdots a_k\|_r
        \le
        \prod_{j=1}^k\|a_j\|_{p_j}.
\]
In particular, if \(\sum_j p_j^{-1}=1\), then \(a_1\cdots a_k\in L^1(M)\).
In this case cyclicity holds:
\[
        \operatorname{tr}(a_1a_2\cdots a_k)
        =
        \operatorname{tr}(a_2\cdots a_ka_1).
\]
\end{lemma}

\begin{proof}
The norm estimate is the generalized H\"older inequality for Haagerup
\(L^p\)-spaces.  For cyclicity, put \(b=a_2\cdots a_k\).  By H\"older,
\(b\in L^q(M)\), where \(1/p_1+1/q=1\).  The two-factor trace pairing gives
\(\operatorname{tr}(a_1b)=\operatorname{tr}(ba_1)\).  This is the desired
cyclic permutation.
\end{proof}

\begin{lemma}[Integer power continuity]
\label{lem:integer-power-continuity}
Let \(N\geq 1\) be an integer. If \(x_m,x\in L^N(M)_+\) and
\begin{equation*}
    \|x_m-x\|_N\longrightarrow 0,
\end{equation*}
then
\begin{equation*}
    \|x_m^N-x^N\|_1\longrightarrow 0.
\end{equation*}
\end{lemma}

\begin{proof}
Since \(x_m\to x\) in \(L^N(M)\), the sequence
\(\bigl(\|x_m\|_N\bigr)_m\) is bounded. Set
\begin{equation*}
    C
    :=
    \max\left\{
        1,\|x\|_N,\sup_m\|x_m\|_N
    \right\}
    <\infty.
\end{equation*}

All products below are taken in the algebra
\(L^0(\cR,\tau)\) of measurable operators affiliated with the crossed
product. The algebraic telescoping identity follows by direct
expansion:
\begin{align*}
    \sum_{k=0}^{N-1}
        x_m^{N-1-k}(x_m-x)x^k
    &=
    \sum_{k=0}^{N-1}
        \left(
            x_m^{N-k}x^k
            -
            x_m^{N-1-k}x^{k+1}
        \right)
    \\
    &=
    x_m^N
    +
    \sum_{k=1}^{N-1}x_m^{N-k}x^k
    -
    \sum_{k=1}^{N-1}x_m^{N-k}x^k
    -
    x^N
    \\
    &=x_m^N-x^N.
\end{align*}

For each \(k=0,\ldots,N-1\), the summand
\begin{equation*}
    x_m^{N-1-k}(x_m-x)x^k
\end{equation*}
is a product of exactly \(N\) elements of \(L^N(M)\): \(N-1-k\)
copies of \(x_m\), followed by \(x_m-x\), followed by \(k\) copies of
\(x\). 
Applying Lemma~\ref{lem:holder-products-trace-cyclicity} with
\begin{equation*}
    p_1=\cdots=p_N=N,
    \qquad
    \frac{1}{N}+\cdots+\frac{1}{N}=1,
\end{equation*}
shows that this product belongs to \(L^1(M)\) and gives
\begin{align*}
    \left\|
        x_m^{N-1-k}(x_m-x)x^k
    \right\|_1
    &\leq
    \|x_m\|_N^{N-1-k}
    \|x_m-x\|_N
    \|x\|_N^k
    \\
    &\leq
    C^{N-1}\|x_m-x\|_N.
\end{align*}
Consequently,
\begin{align*}
    \|x_m^N-x^N\|_1
    &\leq
    \sum_{k=0}^{N-1}
    \left\|
        x_m^{N-1-k}(x_m-x)x^k
    \right\|_1
    \\
    &\leq
    N C^{N-1}\|x_m-x\|_N
    \longrightarrow 0.
\end{align*}
\end{proof}

\begin{lemma}[Sandwich continuity]\label{lem:sandwich-continuity}
Let \(h\in L^1(M)_+\), let \(1\le p<\infty\), and let
\((a_i)_i\subset M\) be uniformly bounded with \(a_i\to a\) strongly.  Then
\[
        h^{1/(2p)}a_i h^{1/(2p)}
        \longrightarrow
        h^{1/(2p)}a h^{1/(2p)}
        \qquad\text{in }L^p(M).
\]
\end{lemma}

\begin{proof}
By \cite[Lemma 2.3]{Junge:2002osf},
\[
        (a_i-a)h^{1/(2p)}
        \longrightarrow0
        \qquad\text{in }L^{2p}(M).
\]
Hence H\"older's inequality gives
\[
\begin{aligned}
\left\|
        h^{1/(2p)}(a_i-a)h^{1/(2p)}
\right\|_p
&\le
        \|h^{1/(2p)}\|_{2p}
        \,
        \|(a_i-a)h^{1/(2p)}\|_{2p}
        \longrightarrow0 .
\end{aligned}
\]
\end{proof}


\section{Analytical results used for semigroup and modular smoothing}\label{Sec:misc-results-Sec45}

In this section, we collect some technical results used in Sec.~\ref{sec:cyclic-ward} to construct the semigroup and null smoothing. We adopt the notations and definitions from Sec.~\ref{Sec:prelim} and the HSMI structure from Sec.~\ref{Sec:HSMI}. Throughout this section $M_B\subset M_A=:M$ is a HSMI with common faithful normal state $\omega$, and the null translation endomorphism \[T_c: M_A\longmapsto M_c\subset M_A, \quad T_c(a)=U_{-c}aU_c, \quad a\in M_A.\] We denote the modular flow with respect to $\omega$ by $\sigma^\omega_t$. 

Let $L^n(M)$ denote the Haagerup $L^n$ space. We also recall the dense embedding of $M$ into $L^n(M)$ denoted by $j_n(M)$ in Definition~\ref{def:symmetricembedding}.

\subsection{Generator of the null translation endomorphism}
 
Recall the definition of the generator, $\delta$, of the null translation endomorphism $T_r$ from \eqref{eq:delta-dom-def-main}:
\begin{equation}\label{eq:delta-dom-def}
    \d a\defeq \lim_{r\downarrow0}\frac{T_r(a)-a}{r}, \quad \Dom(\d)\defeq\lbr a\in M:\lim_{r\downarrow0}\frac{T_r(a)-a}{r} \ \text{exists in operator norm} \rbr.
\end{equation}

The generator $\delta$ acts as a derivation on its domain.
\begin{lemma}[Derivation on the domain]\label{lem:derivation-on-dom}
$\Dom(\delta)$ is a unital $*$-subalgebra of $M$, and
    \begin{equation}\label{eq:delta-derivation-domain}
        \delta(ab)= \delta(a)b+a\delta(b),\quad a,b\in \Dom(\delta).
    \end{equation}
\end{lemma}
\begin{proof}
    To establish the derivation property, let $a,b\in M$ both lie in $\Dom(\delta)$ with operator-norm-convergent difference quotients.  Multiplicativity of $T_r$ gives the algebraic identity
\begin{equation*}
  \frac{T_r(ab)-ab}{r}
  =\frac{T_r(a)-a}{r}\,T_r(b)+a\,\frac{T_r(b)-b}{r}\,.
\end{equation*}
The factor $T_r(b)$ converges to $b$ in operator norm and the difference quotient on $a$ converges to $\delta a$ in operator norm, so the first summand converges to $\delta(a)b$.  The second summand converges to $a\,\delta(b)$ by left multiplication by the fixed element $a\in M$.  Hence $ab\in\Dom(\delta)$ with $\delta(ab)=\delta(a)b+a\,\delta(b)$.

Using the fact that $T_r(a)^*=T_r(a^*)$ and the definition of the domain, it is straightforward to see that
\begin{equation*}
    a\in \Dom(\delta) \Rightarrow a^*\in \Dom(\delta), \quad \delta(a^*)=(\delta a)^*\,.
\end{equation*}
Finally, $\Dom(\delta)$ is unital since $T_r$ preserves the identity operator. All of the above combined with the linearity of $T_r$ establishes the claim that $\Dom(\delta)$ is a unital $*$-subalgebra of $M$.
\end{proof}

The following Lipschitz estimate is useful for dominated convergence arguments involving $T_r(a_f)$ and $\delta$.
\begin{lemma}[Lipschitz bound for null-smoothed elements]
\label{lem:lipschitz-af}
Let $a\in M$ and $f\in C_{c,0}^\infty((0,\I))$, and let
\[
  a_f=\int_0^\I f(t)\,T_t(a)\,dt,\qquad
  \delta a_f=-\int_0^\I f'(t)\,T_t(a)\,dt
\]
be the ultraweak integrals defined in \eqref{eq:def-null-smearing} and \eqref{eq:delta-af-defn}.  Then the
integrated identity
\begin{equation}
\label{eq:integrated-difference-af}
  T_u(a_f)-a_f=\int_0^u T_r(\delta a_f)\,dr\,,\qquad u\ge 0,
\end{equation}
holds in $M$, where the right-hand side is an ultraweak integral.  In particular, $a_f$
satisfies the uniform Lipschitz estimate
\begin{equation}
\label{eq:lipschitz-af-quotient}
  \left\|\frac{T_u(a_f)-a_f}{u}\right\|\le\|\delta a_f\|,\qquad u>0\,.
\end{equation}
\end{lemma}

\begin{proof}
Extend $f$ to all of $\mathbb R$ by zero; the extension is smooth because $f$ vanishes in a neighbourhood of zero.  
The $*$-endomorphism $T_u$ is normal on $M$, so it commutes with the ultraweak integral defining $a_f$:
\begin{equation*}
  T_u(a_f)=\int_0^\I f(t)\,T_{t+u}(a)\,dt=\int_u^\I f(v-u)\,T_{v}(a)\,dv
  =\int_0^\I f(v-u)\,T_{v}(a)\,dv.
\end{equation*}
This gives us
\begin{equation*}
  T_u(a_f)-a_f
  =\int_0^\I (f(v-u)-f(v))\,T_{v}(a)\,dv
  =-\int_0^\I\!\! \lb\int_0^u f'(v-r)\,dr \rb\,T_{v}(a)\,dv\,.
\end{equation*}
The integrand on the right is jointly measurable and majorized by
$|f'(v-r)|\,\|a\|$, with
\[
    \int_0^u \int_0^\I|f'(v-r)|\,dv\,dr=u\,\|f'\|_{L^1}<\I.
\]
So, by Fubini's theorem, the order of integration may be exchanged.  Substituting $w=v-r$ in the inner integral, and using $f'(w)=0$ for $w<0$ from the smooth zero-extension of $f$,
\begin{equation*}
  \int_0^\I f'(v-r)\,T_{v}(a)\,dv
  =\int_0^\I f'(w)\,T_{w+r}(a)\,dw
  =T_r\!\biggl(\int_0^\I f'(w)\,T_w(a)\,dw\biggr)
  =-T_r(\delta a_f),
\end{equation*}
where the second equality again uses normality of $T_r$ to commute it with the ultraweak
integral.  This proves~\eqref{eq:integrated-difference-af}.

Now~\eqref{eq:lipschitz-af-quotient} follows from
contractivity of $T_r$: $\|T_r(\delta a_f)\|\le\|\delta a_f\|$
for all $r\ge 0$, and the triangle inequality applied to the ultraweak integral
in~\eqref{eq:integrated-difference-af} proves the claim:
\begin{equation*}
  \|T_u(a_f)-a_f\|\le\int_0^u\|T_r(\delta a_f)\|\,dr\le u\,\|\delta a_f\|\,.
\end{equation*}
\end{proof}

\subsection{Compatibility of ultraweak and Bochner integrals with the dense embedding}

The following result shows that ultraweak and Bochner integrals are compatible with the embedding.
\begin{lemma}[Compatibility of $j_n$ with ultraweak and Bochner integrals]\label{lem:jn-compatible-weak}
Let $(\mathcal S,\nu)$ be a finite complex measure space and let $s\mapsto a_s \in M$ be bounded and ultraweakly measurable. Assume that $j_n(a_s)$ is Bochner integrable in $L^n(M)$, and define
\begin{equation}
    a \defeq \int_{\mathcal S} a_s \,d\nu(s)
\end{equation}
as an ultraweak integral. Then
\begin{equation} \label{eq:jn-compatible-int}
    j_n(a) = \int_{\mathcal S} j_n(a_s) \, d\nu(s) \qquad\text{in }L^n(M).
\end{equation}
\end{lemma}
\begin{proof}
For $b\in M$, set 
\begin{equation}
\label{eq:pairing-functional}
    \varphi_b:M\to\mathbb C,\qquad
    \varphi_b(c)\defeq \tr\bigl(j_n(c)\,j_{q_n}(b)\bigr).
\end{equation}
By the pairing formula \eqref{eq:core-pairing},
\begin{equation*}
    \varphi_b(c)=\langle c^*\Omega,J_\omega b^*\Omega\rangle=\langle\Omega,c\,J_\omega b^*\Omega\rangle\,,
\end{equation*}
which is the vector functional $\omega_{\Omega,J_\omega b^*\Omega}\in M_*$, hence normal.

Pairing the left hand side of \eqref{eq:jn-compatible-int} with $j_{q_n}(b)$, gives by normality of $\varphi_b$
\begin{equation*}
    \varphi_b \lb \int_{\mathcal S} a_s\, d\nu(s)\rb=\int_{\mathcal S} \varphi_b(a_s)\, d\nu(s). 
\end{equation*}
Pairing the right hand side gives the same result, since the trace pairing is a continuous linear functional on $L^n(M)$ and thus commutes with the Bochner integral \cite[Proposition C.4]{engel.nagel:00:one-parameter}.
Since $j_{q_n}(M)$ is dense in $L^{q_n}(M)$, and the pairing separates points of the dual $L^n(M)$, the two $L^n(M)$ elements agree,  establishing \eqref{eq:jn-compatible-int}. 
\end{proof}

\subsection{Modular smoothing}

For $\varepsilon>0$, define
\begin{equation}\label{eq:gamma-appendix}
  \gamma_\varepsilon(s)\defeq \frac{1}{\sqrt{2\pi}\,\varepsilon}
  \exp\left(-\frac{s^2}{2\varepsilon^2}\right),
  \qquad s\in\mathbb R\,.
\end{equation}
Thus $\int_{\mathbb R}\gamma_\varepsilon(s)\,ds=1$, and the Gaussian measures $\g_\ve(s)\,ds$ form an approximate identity (see Definition~\ref{def:approx-id}). We use this approximate identity together with the ultraweak integral to construct modular analytic approximations for any operator in the algebra as follows.

\begin{lemma}[Modular-Gaussian smearing produces entire elements]
\label{lem:gaussian-modular-entire}
Let $M$ be a von Neumann algebra with faithful normal state $\w$ and modular automorphism group $\s^\w$.  For $\ve>0$, let $\g_\ve$ be the Gaussian kernel in \eqref{eq:gamma-appendix}.
For $b\in M$, define the ultraweak integral (Definition~\ref{def:ultraweak-int})
\begin{equation}
\label{eq:b-eps-def}
  b_\ve\defeq \int_{\mathbb R}\g_\ve(s)\,\s^\w_s(b)\,ds\in M\,.
\end{equation}
Then $b_\ve$ is entire analytic for $\s^\w$, and its complex modular extension is given by the ultraweak integral
\begin{equation}
\label{eq:sigma-z-b-eps}
  \s^\w_z(b_\ve)=\int_{\mathbb R}\g_\ve(s-z)\,\s^\w_s(b)\,ds,\qquad z\in\mathbb C\,.
\end{equation}
Further, as $\ve\to0$, $b_\ve\to b$ in the strong-$*$ operator topology.
\end{lemma}

\begin{proof}
For each $z\in\mathbb C$, the kernel bound
\begin{equation}
\label{eq:kernel-bound}
  |\g_\ve(s-z)|\le e^{(\Im z)^2/(2\ve^2)}\,\g_\ve(s-\Re z)
\end{equation}
shows that $s\mapsto\g_\ve(s-z)$ is integrable, with the bound locally uniform in $z$.  Hence the right-hand side of \eqref{eq:sigma-z-b-eps} is a well-defined ultraweak integral; denote it by $\widetilde b(z)\in M$.

For each $\varphi\in M_*$, the scalar function
\begin{equation}
  z\mapsto\varphi(\widetilde b(z))=\int_{\mathbb R}\g_\ve(s-z)\,\varphi(\s^\w_s(b))\,ds
\end{equation}
is entire: $\g_\ve(s-z)$ is entire in $z$ for each $s$, 
and \eqref{eq:kernel-bound} together with $\|\varphi\circ\s^\w_s\|\le\|\varphi\|$ provides locally uniform domination.  Thus $\widetilde b:\mathbb C\to M$ is weakly entire.

For real $z=t\in\mathbb R$, normality of $\s^\w_t$ gives $\varphi\circ\s^\w_t\in M_*$ for every $\varphi\in M_*$, so applying $\varphi\circ\s^\w_t$ to \eqref{eq:b-eps-def} and substituting $s\mapsto s-t$ in the kernel yields
\begin{equation}
  \varphi(\s^\w_t(b_\ve))=\int_{\mathbb R}\g_\ve(s)\,\varphi(\s^\w_{s+t}(b))\,ds=\int_{\mathbb R}\g_\ve(s-t)\,\varphi(\s^\w_s(b))\,ds=\varphi(\widetilde b(t))\,.
\end{equation}
Since this holds for all $\varphi\in M_*$, $\s^\w_t(b_\ve)=\widetilde b(t)$ for every $t\in\mathbb R$.

Hence $\widetilde b$ is a weakly entire $M$-valued extension of the modular orbit $t\mapsto\s^\w_t(b_\ve)$.  By definition, $b_\ve$ is entire analytic for $\s^\w$ with $\s^\w_z(b_\ve)=\widetilde b(z)$ for all $z\in\mathbb C$, proving \eqref{eq:sigma-z-b-eps}.

The convergence is easily seen using the fact that $\g_\ve(s)$ is an approximate identity.
\end{proof}

\begin{lemma}[Strongly continuous modular isometry group on \(L^n(M)\)]
\label{lem:modular-Ln-group}
Define $\sigma_s:= \sigma_s^\omega$.
For every \(s\in\mathbb R\), the map
\[
        j_n(a)\longmapsto j_n(\sigma_s(a)),
        \qquad a\in M,
\]
extends uniquely to a positive surjective isometry
\[
        \Sigma_s^{(n)}:L^n(M)\longrightarrow L^n(M).
\]
Moreover,
\begin{equation}
\label{eq:Sigma-group-property}
        \Sigma_{s+t}^{(n)}
        =
        \Sigma_s^{(n)}\Sigma_t^{(n)},
        \qquad
        \Sigma_0^{(n)}=I,
        \qquad s,t\in\mathbb R,
\end{equation}
and the map
\[
        \mathbb R\ni s\longmapsto \Sigma_s^{(n)}x\in L^n(M)
\]
is norm-continuous for every \(x\in L^n(M)\).  Thus
\((\Sigma_s^{(n)})_{s\in\mathbb R}\) is a strongly continuous group of
positive isometries on \(L^n(M)\).
\end{lemma}

\begin{proof}
For each \(s\in\mathbb R\), the modular automorphism
\(\sigma_s:M\to M\) is normal, unital, positive, and
\(\omega\)-preserving.  Hence Theorem~\ref{thm:extension-Lp} shows that the
map
\[
        j_n(a)\longmapsto j_n(\sigma_s(a))
\]
extends uniquely to a positive contraction
\[
        \Sigma_s^{(n)}:L^n(M)\longrightarrow L^n(M).
\]

We first verify the group property.  For \(a\in M\) and \(s,t\in\mathbb R\),
\[
\begin{aligned}
        \Sigma_s^{(n)}\Sigma_t^{(n)}j_n(a)=
        j_n(\sigma_s(\sigma_t(a)))  =
        j_n(\sigma_{s+t}(a)) =
        \Sigma_{s+t}^{(n)}j_n(a).
\end{aligned}
\]
Since \(j_n(M)\) is norm-dense in \(L^n(M)\) and both sides are bounded,
the equality extends to all of \(L^n(M)\).  Similarly,
\(\Sigma_0^{(n)}=I\).  In particular,
\[
        \Sigma_{-s}^{(n)}\Sigma_s^{(n)}
        =
        \Sigma_s^{(n)}\Sigma_{-s}^{(n)}
        =
        I,
\]
so \(\Sigma_s^{(n)}\) is bijective with inverse
\(\Sigma_{-s}^{(n)}\).

Since both \(\Sigma_s^{(n)}\) and \(\Sigma_{-s}^{(n)}\) are contractions,
for every \(x\in L^n(M)\),
\[
        \|x\|_n
        =
        \|\Sigma_{-s}^{(n)}\Sigma_s^{(n)}x\|_n
        \leq
        \|\Sigma_s^{(n)}x\|_n
        \leq
        \|x\|_n.
\]
Consequently,
\[
        \|\Sigma_s^{(n)}x\|_n=\|x\|_n,
\]
and \(\Sigma_s^{(n)}\) is an isometry.

It remains to prove strong continuity.  We first prove continuity at
\(s=0\) on the dense subspace \(j_n(M)\).  Let \(a\in M\).  The modular
automorphism group is point strong-\(*\) continuous, and
\[
        \|\sigma_s(a)\|=\|a\|,
        \qquad s\in\mathbb R.
\]
Therefore,
\[
        \sigma_s(a)\longrightarrow a
        \qquad\text{strongly as }s\to0,
\]
with the family \((\sigma_s(a))_{s\in\mathbb R}\) uniformly bounded.
Lemma~\ref{lem:sandwich-continuity} then gives
\begin{equation}
\label{eq:Sigma-cont-core}
\begin{aligned}
        \|\Sigma_s^{(n)}j_n(a)-j_n(a)\|_n
        &=
        \|j_n(\sigma_s(a))-j_n(a)\|_n
        \longrightarrow0
        \qquad (s\to0).
\end{aligned}
\end{equation}

Now let \(x\in L^n(M)\) and let \(\varepsilon>0\).  Choose
\(y\in j_n(M)\) such that
\[
        \|x-y\|_n<\varepsilon.
\]
Since \(\Sigma_s^{(n)}\) is an isometry,
\[
\begin{aligned}
        \|\Sigma_s^{(n)}x-x\|_n
        &\leq
        \|\Sigma_s^{(n)}(x-y)\|_n
        +
        \|\Sigma_s^{(n)}y-y\|_n
        +
        \|y-x\|_n                                    \\
        &\leq
        2\varepsilon
        +
        \|\Sigma_s^{(n)}y-y\|_n.
\end{aligned}
\]
By \eqref{eq:Sigma-cont-core}, the last term tends to zero as \(s\to0\).
Since \(\varepsilon>0\) is arbitrary,
\[
        \Sigma_s^{(n)}x\longrightarrow x
        \qquad\text{in }L^n(M)
\]
as \(s\to0\).

Finally, for any \(s_0\in\mathbb R\), the group property and isometry give
\[
\begin{aligned}
        \|\Sigma_s^{(n)}x-\Sigma_{s_0}^{(n)}x\|_n
        &=
        \left\|
        \Sigma_{s_0}^{(n)}
        \bigl(\Sigma_{s-s_0}^{(n)}x-x\bigr)
        \right\|_n                                   \\
        &=
        \|\Sigma_{s-s_0}^{(n)}x-x\|_n
        \longrightarrow0
\end{aligned}
\]
as \(s\to s_0\).  This proves strong continuity on \(\mathbb R\).
\end{proof}


\section{R\'enyi QNEC for \texorpdfstring{$n=2$}{n=2}}
\label{subsec:alpha-two-rigorous}

We have proven RQNEC for all integers $n\geq2$ in Corollary~\ref{cor:integer-renyi-qnec-fixed}. In this section we present a slightly different proof of RQNEC for $n=2$ that is easier to interpret directly on $\mathcal H$.
The R\'enyi QNEC for $n=2$ is special since $L^2(M)$ is a Hilbert space. 

The main inputs for this section are the HSMI $M_B\subset M_A=: M$, with common faithful $\omega\in (M_A)_*^+$ and associated properties from Theorem~\ref{thm:modulartranslations}, the null translation semigroup from Proposition~\ref{prop:null-translation-Ln-semigroup}, and the sandwiched densities from Sec~\ref{subsec:integer-density-propagation}.

Using the isometric equivalence of standard forms, we can identify the action of the null translation semigroup $V_c^{(2)}$ on $L^2(M)$ with left multiplication by $e^{-cP}$ on the vector representative $\xi_0 \in \mathcal H$ of the $L^2(M)$ sandwiched density $x_0$. This lends itself to a more transparent physical interpretation of the SRD and RQNEC in terms of the averaged null energy operator $P$.

\subsection{Preliminaries on standard forms and modular operators}

We first recall some useful preliminaries. We refer to \cite{Hiai:2021qfv} for this section.
In order to understand how modular data, such as the modular operators, relative modular operators, etc., are represented on the Haagerup $L^p$ space it is useful to define, for each $x\in M$, the left and right actions $\mathfrak L(x), \mathfrak R(x)$ on the $L^2(M)$ Hilbert space.
\begin{definition}[Left and right action]
    \[
        \mathfrak L(x) a \defeq x a, \quad \mathfrak R(x) a \defeq ax, \, a\in L^2(M)\,.
    \]
    The involution $J$ is represented on $L^2(M)$ as: 
    \[
        Ja=a^*\,.
    \]
\end{definition}

\begin{theorem}\label{thm:std-form-Haagerup}
    For each $x\in M$ define the left, right actions $\mathfrak L(x),\mathfrak R(x)$ and the involution $J$. Then:
    \begin{enumerate}
        \item $\mathfrak L(M)$($\mathfrak R(M)$) is a normal faithful (anti-)representation of $M$ on $L^2(M)$.
        \item The von Neumann algebras $\mathfrak L(M), \mathfrak R(M)$ are commutants of each other with $\mathfrak R(M) = J \mathfrak L(M)J$.
        \item $(\mathfrak L(M),L^2(M),J,L^2(M)_+)$ is a standard form of $M$ (see \cite{Hiai:2021cks,Hiai:2021qfv}).
    \end{enumerate}
\end{theorem}


\begin{remark}[Modular operator on $L^2(M)$]\label{rem:Modular-op-L2}
    The previous lemma shows that if $\vf$ is faithful, then
    \begin{equation}
        \Delta^{it}_\varphi \xi = h_\varphi^{it} \xi h_\varphi^{-it}, \quad \forall \xi\in L^2(M),\,t\in \mathbb R\,.
    \end{equation}
    Furthermore, since the analytic continuation is unique, we have:
    \begin{equation}\label{eq:delta-Lp-action}
        \Delta^{\frac p2}_\varphi (x h_\varphi^{1/2}) =  h_\varphi^{\frac p2} x  h_\varphi^{\frac{1-p}{2}}, \quad x \in M, \, 0\leq p\leq 1\,.
    \end{equation}
    For non-faithul $\vf$, the formulas are to be understood with usual support projections on $s(\vf)Ms(\vf)$.
\end{remark}
Similarly, we have the representation of the relative modular operator on $L^2(M)$.
\begin{remark}[Relative modular operator on $L^2(M)$]\label{rem:Relative-Modular-op-L2}
    For any $\psi,\varphi \in M_*^+$, we have the relative modular operator \cite{Hiai:2021qfv}
    \[
        \Delta_{\psi,\varphi}^{it}\xi = h_\psi^{it}\xi h_\varphi^{-it}, \quad \forall \xi\in s(\psi)L^2(M)s(\vf)\,\,, t\in \mathbb R\,,
    \]
    \[
        \Delta_{\psi,\varphi}^{\frac p2}(x h_\varphi^{1/2}) = h_\psi^{\frac{p}{2}}x h_\varphi^{\frac{1-p}{2}}, \quad x \in M, \, 0\leq p\leq 1\,,
    \]
    with the convention
    \[
        h_\psi^0=s(\psi), \quad h_\varphi^0=s(\varphi),\quad \Delta_{\psi,\varphi}^0=s(\psi)Js(\varphi)J.
    \]
\end{remark}

\subsection{Proof of RQNEC for \texorpdfstring{$n=2$}{n=2}}
In this subsection we work with Haagerup's $L^2(M)$ space and the pairing $\braket{\cdot,\cdot}_{L^2(M)}$ refers to the inner product on $L^2(M)$ defined as
\[
    \braket{a,b}\defeq \tr(a^*b), \quad a,b\in L^2(M). 
\]

Recall the $L^2(M)$ semigroup $V_c^{(2)}$ defined in Sec.~\ref{subsec:test-and-dual-semigroups}. The following proposition provides an explicit characterization of $V_c^{(2)}$ in terms of the null translation generator $P$ associated with HSMI. Let $\omega\in M_*^+$ be faithful and denote its vector representative by $\Omega\in \mathcal H$.

\begin{lemma}[\(L^2\)-implementation of the null flow]
\label{prop:standard-L2-implementation}
Define \(W : L^2(M) \to \mathcal H\) on the dense subspace \(M h^{1/2} \subset L^2(M)\) by
\[
W(a h^{1/2})  \defeq  a \Omega, \qquad a \in M \,.
\]
Then \(W\) extends to a unitary. Let
\[
\widehat{P}  \defeq  W^{-1} P W,
\qquad
\widehat{V}_c^{(2)}  \defeq  e^{-c \widehat{P}}, \qquad c \ge 0 \,.
\]
For every \(a \in M^{\mathrm{an}}\),
\[
\widehat{V}_c^{(2)}\big(h^{1/4} a h^{1/4}\big)
=
h^{1/4} T_c(a) h^{1/4}\,,
\]
which allows the identification
\[
    V_c^{(2)} = \widehat{V}_c^{(2)}.
\]
In particular, \((V_c^{(2)})_{c \ge 0}\) is a strongly continuous semigroup of self-adjoint contractions on
\(L^2(M)\).
\end{lemma}

\begin{proof}
For \(a \in M\),
\[
\|a h^{1/2}\|_2^2
=
\tr(h^{1/2} a^* a h^{1/2})
=
\omega(a^* a)
=
\|a \Omega\|^2\,.
\]
Since \(M h^{1/2}\) is dense in \(L^2(M)\) and \(M \Omega\) is dense in \(\mathcal H\), \(W\)
extends to a unitary. Moreover, \(W\) intertwines the left action of \(M\) on \(L^2(M)\) with
the given representation of \(M\) on \(\mathcal H\). By uniqueness of the standard form \cite{Haagerup1975},
\[
W \Delta_{L^2(M)}^{it} = \Delta_{\Omega;M}^{it} W, \qquad t \in \mathbb{R}\,,
\]
where \(\Delta_{L^2(M)}\) denotes the modular operator of the Haagerup standard form defined in Theorem~\ref{thm:std-form-Haagerup}. Hence, for
\(a \in M^{\mathrm{an}}\),
\[
W\big(h^{1/4} a h^{1/4}\big)
=
W\big(\Delta_{L^2(M)}^{1/4}(a h^{1/2})\big)
=
\Delta_{\Omega;M}^{1/4} a \Omega\,.
\]

Now let \(a \in M^{\mathrm{an}}\). Using \(T_c(a)\Omega = U_{-c} a \Omega\), we get
\[
W\big(h^{1/4} T_c(a) h^{1/4}\big)
=
\Delta_{\Omega;M}^{1/4} U_{-c} a \Omega\,.
\]
By the imaginary-time continuation of the Borchers covariance relation (Theorem~\ref{thm:modulartranslations}(4)),
\[
\Delta_{\Omega;M}^{1/4} U_{-c} a \Omega
=
U_{-ic} \Delta_{\Omega;M}^{1/4} a \Omega
=
e^{-cP} \Delta_{\Omega;M}^{1/4} a \Omega\,.
\]
Therefore
\[
W\big(h^{1/4} T_c(a) h^{1/4}\big)
=
e^{-cP} W\big(h^{1/4} a h^{1/4}\big)
=
W \widehat{V}_c^{(2)}\big(h^{1/4} a h^{1/4}\big)\,.
\]
Since \(W\) is unitary, this proves the claimed action formula on the dense core $h^{1/4}M^{\text{an}}h^{1/4}$ of $L^2(M)$. The equality extends by boundedness to all of $L^2(M)$. The remaining statements follow
from the spectral calculus for the positive self-adjoint operator \(P\).
\end{proof}

\begin{proposition}[\(n=2\) SRD as expectation value]
\label{prop:xc-semigroup-orbit}
Let $\psi,x_0,x_c, Q_2(c)$ be as in Lemma \ref{lem:density-propagation-fixed} with $n=2$, so that $0<Q_2(c)<\I$ for $c\geq0$. 
Then
\[
x_c = V_c^{(2)} x_0 = W^{-1} e^{-cP} W x_0\,.
\]
If \(\xi_0  \defeq  W x_0 \in \mathcal H\), then
\[
Q_2(c)  \defeq  Q_2(\psi_c\|\omega;M)
=
\|x_c\|_2^2
=
\|e^{-cP}\xi_0\|^2
=
\langle \xi_0, e^{-2cP}\xi_0\rangle_{\mathcal H}\,.
\]
\end{proposition}

\begin{proof}
By Lemma \ref{prop:standard-L2-implementation},
\[
    V_c^{(2)}=W^{-1}e^{-cP}W.
\]
Therefore, with \(\xi_0=Wx_0\),
\[
    x_c=W^{-1}e^{-cP}\xi_0,
\]
and
\[
    Q_2(c)=\|x_c\|_2^2 = \|e^{-cP}\xi_0\|^2 = \langle \xi_0,e^{-2cP}\xi_0\rangle.
\]
The formula for \(Q_2(c)\) is immediate.
\end{proof}

Thus, we have translated the null-flowed SRD to the logarithm of the expectation value of $e^{-2cP}$ in the $\mathcal H$ representative $\xi_0$ of the $L^2(M)$ sandwiched density $x_0$, where $P$ is the averaged null energy operator for the null cut algebras. 
RQNEC for $n=2$ now follows simply from spectral calculus.
\begin{theorem}[R\'enyi QNEC at \(n=2\)]
\label{thm:RQNEC-alpha-two}
Assume \(Q_2(\psi\|\omega;M) < \infty\). Then the function
\[
S_2(c)  \defeq  D_2(\psi_c\|\omega;M)
\]
is \(C^\infty\) and convex on \((0,\infty)\). Equivalently,
\[
S_2''(c) \ge 0, \qquad c>0\,.
\]
\end{theorem}

\begin{proof}
By Proposition~\ref{prop:xc-semigroup-orbit},
\[
Q_2(c) = \langle \xi_0, e^{-2cP}\xi_0\rangle_{\mathcal H},
\qquad
\xi_0  \defeq  W x_0 \in \mathcal H\,. 
\]
Let \(E_P\) be the spectral measure of \(P\), and define a finite positive measure on
\([0,\infty)\) by
\[
\mu(B)  \defeq  \|E_P(B)\xi_0\|^2\,.
\]
Then
\[
Q_2(c)=\int_{[0,\infty)} e^{-2c\lambda}\,d\mu(\lambda)\,.
\]
For \(c>0\), differentiation under the integral gives
\[
Q_2'(c) = -2\int \lambda e^{-2c\lambda}\,d\mu(\lambda),
\qquad
Q_2''(c) = 4\int \lambda^2 e^{-2c\lambda}\,d\mu(\lambda)\,.
\]
Set \(d\nu_c(\lambda)  \defeq  e^{-2c\lambda}\,d\mu(\lambda)\). By Cauchy-Schwarz in
\(L^2([0,\infty),d\nu_c)\),
\[
\big(Q_2'(c)\big)^2
=
4\Big(\int \lambda\,d\nu_c\Big)^2
\le
4\Big(\int 1\,d\nu_c\Big)\Big(\int \lambda^2\,d\nu_c\Big)
=
Q_2(c)Q_2''(c)\,.
\]
Since \(Q_2(c)>0\), it follows that
\[
\frac{d^2}{dc^2}\log Q_2(c)
=
\frac{Q_2(c)Q_2''(c)-Q_2'(c)^2}{Q_2(c)^2}
\ge 0\,.
\]
This proves the convexity of \(S_2(c)\).

For normalized states, \(S_2(c) = \log Q_2(c)\); in general one has
\(S_2(c) = \log Q_2(c) - \log \psi(1)\), so the additive constant is irrelevant for convexity.
\end{proof}


\end{appendix}


\providecommand{\href}[2]{#2}\begingroup\raggedright\begin{thebibliography}{10}

\bibitem{Bousso:2015mna}
R.~Bousso, Z.~Fisher, S.~Leichenauer and A.C.~Wall, \emph{{Quantum focusing
  conjecture}}, \href{https://doi.org/10.1103/PhysRevD.93.064044}{\emph{Phys.
  Rev. D} {\bfseries 93} (2016) 064044}
  [\href{https://arxiv.org/abs/1506.02669}{{\ttfamily 1506.02669}}].

\bibitem{Bousso:2015wca}
R.~Bousso, Z.~Fisher, J.~Koeller, S.~Leichenauer and A.C.~Wall, \emph{{Proof of
  the Quantum Null Energy Condition}},
  \href{https://doi.org/10.1103/PhysRevD.93.024017}{\emph{Phys. Rev. D}
  {\bfseries 93} (2016) 024017}
  [\href{https://arxiv.org/abs/1509.02542}{{\ttfamily 1509.02542}}].

\bibitem{Malik:2019dpg}
T.A.~Malik and R.~Lopez-Mobilia, \emph{{Proof of the quantum null energy
  condition for free fermionic field theories}},
  \href{https://doi.org/10.1103/PhysRevD.101.066028}{\emph{Phys. Rev. D}
  {\bfseries 101} (2020) 066028}
  [\href{https://arxiv.org/abs/1910.07594}{{\ttfamily 1910.07594}}].

\bibitem{Koeller:2015qmn}
J.~Koeller and S.~Leichenauer, \emph{{Holographic Proof of the Quantum Null
  Energy Condition}},
  \href{https://doi.org/10.1103/PhysRevD.94.024026}{\emph{Phys. Rev. D}
  {\bfseries 94} (2016) 024026}
  [\href{https://arxiv.org/abs/1512.06109}{{\ttfamily 1512.06109}}].

\bibitem{Balakrishnan:2017bjg}
S.~Balakrishnan, T.~Faulkner, Z.U.~Khandker and H.~Wang, \emph{{A General Proof
  of the Quantum Null Energy Condition}},
  \href{https://doi.org/10.1007/JHEP09(2019)020}{\emph{JHEP} {\bfseries 09}
  (2019) 020} [\href{https://arxiv.org/abs/1706.09432}{{\ttfamily
  1706.09432}}].

\bibitem{Ceyhan:2018zfg}
F.~Ceyhan and T.~Faulkner, \emph{{Recovering the QNEC from the ANEC}},
  \href{https://doi.org/10.1007/s00220-020-03751-y}{\emph{Commun. Math. Phys.}
  {\bfseries 377} (2020) 999}
  [\href{https://arxiv.org/abs/1812.04683}{{\ttfamily 1812.04683}}].

\bibitem{Hollands:2025glm}
S.~Hollands and R.~Longo, \emph{{A New Proof of the QNEC}},
  \href{https://doi.org/10.1007/s00220-025-05450-y}{\emph{Commun. Math. Phys.}
  {\bfseries 406} (2025) 269}
  [\href{https://arxiv.org/abs/2503.04651}{{\ttfamily 2503.04651}}].

\bibitem{Wall:2017blw}
A.C.~Wall, \emph{{Lower Bound on the Energy Density in Classical and Quantum
  Field Theories}},
  \href{https://doi.org/10.1103/PhysRevLett.118.151601}{\emph{Phys. Rev. Lett.}
  {\bfseries 118} (2017) 151601}
  [\href{https://arxiv.org/abs/1701.03196}{{\ttfamily 1701.03196}}].

\bibitem{Borchers:1991xk}
H.J.~Borchers, \emph{{The CPT theorem in two-dimensional theories of local
  observables}}, \href{https://doi.org/10.1007/BF02099011}{\emph{Commun. Math.
  Phys.} {\bfseries 143} (1992) 315}.

\bibitem{Wiesbrock:1992mg}
H.W.~Wiesbrock, \emph{{Half sided modular inclusions of von Neumann algebras}},
  \href{https://doi.org/10.1007/BF02098019}{\emph{Commun. Math. Phys.}
  {\bfseries 157} (1993) 83}.

\bibitem{Borchers:1995zg}
H.J.~Borchers, \emph{{On the use of modular groups in quantum field theory}},
  {\emph{Ann. Inst. H. Poincare Phys. Theor.} {\bfseries 63} (1995) 331}.

\bibitem{Borchers1996}
H.-J.~Borchers, \emph{Half-sided modular inclusion and the construction of the
  poincaré group},
  \href{https://doi.org/10.1007/BF02100104}{\emph{Communications in
  Mathematical Physics} {\bfseries 179} (1996) 703}.

\bibitem{Araki:2005we}
H.~Araki and L.~Zsido, \emph{{Extension of the structure theorem of Borchers
  and its application to half-sided modular inclusions}},
  \href{https://doi.org/10.1142/S0129055X05002388}{\emph{Rev. Math. Phys.}
  {\bfseries 17} (2005) 491}
  [\href{https://arxiv.org/abs/math/0412061}{{\ttfamily math/0412061}}].

\bibitem{Kibe:2021qjy}
T.~Kibe, A.~Mukhopadhyay and P.~Roy, \emph{{Quantum Thermodynamics of
  Holographic Quenches and Bounds on the Growth of Entanglement from the
  Quantum Null Energy Condition}},
  \href{https://doi.org/10.1103/PhysRevLett.128.191602}{\emph{Phys. Rev. Lett.}
  {\bfseries 128} (2022) 191602}
  [\href{https://arxiv.org/abs/2109.09914}{{\ttfamily 2109.09914}}].

\bibitem{Banerjee:2022dgv}
A.~Banerjee, T.~Kibe, N.~Mittal, A.~Mukhopadhyay and P.~Roy, \emph{{Erasure
  Tolerant Quantum Memory and the Quantum Null Energy Condition in Holographic
  Systems}}, \href{https://doi.org/10.1103/PhysRevLett.129.191601}{\emph{Phys.
  Rev. Lett.} {\bfseries 129} (2022) 191601}
  [\href{https://arxiv.org/abs/2202.00022}{{\ttfamily 2202.00022}}].

\bibitem{Kibe:2024icu}
T.~Kibe, A.~Mukhopadhyay and P.~Roy, \emph{{Generalized Clausius inequalities
  and entanglement production in holographic two-dimensional CFTs}},
  \href{https://doi.org/10.1007/JHEP04(2025)096}{\emph{JHEP} {\bfseries 04}
  (2025) 096} [\href{https://arxiv.org/abs/2412.13256}{{\ttfamily
  2412.13256}}].

\bibitem{Kibe:2025cqc}
T.~Kibe and P.~Roy, \emph{{Quantum null energy condition in quenched 2D CFTs}},
  \href{https://doi.org/10.1103/2p7v-f28f}{\emph{Phys. Rev. D} {\bfseries 111}
  (2025) 126009} [\href{https://arxiv.org/abs/2503.17448}{{\ttfamily
  2503.17448}}].

\bibitem{Mezei:2019sla}
M.~Mezei and J.~Virrueta, \emph{{The Quantum Null Energy Condition and
  Entanglement Entropy in Quenches}},
  \href{https://arxiv.org/abs/1909.00919}{{\ttfamily 1909.00919}}.

\bibitem{Casini:2023kyj}
H.~Casini, I.~Salazar~Landea and G.~Torroba, \emph{{Irreversibility, QNEC, and
  defects}}, \href{https://doi.org/10.1007/JHEP07(2023)004}{\emph{JHEP}
  {\bfseries 07} (2023) 004}
  [\href{https://arxiv.org/abs/2303.16935}{{\ttfamily 2303.16935}}].

\bibitem{Muller-Lennert:2013liu}
M.~M{\"u}ller-Lennert, F.~Dupuis, O.~Szehr, S.~Fehr and M.~Tomamichel,
  \emph{{On quantum R{\'e}nyi entropies: A new generalization and some
  properties}}, \href{https://doi.org/10.1063/1.4838856}{\emph{J. Math. Phys.}
  {\bfseries 54} (2013) 122203}
  [\href{https://arxiv.org/abs/1306.3142}{{\ttfamily 1306.3142}}].

\bibitem{Wilde:2013bdg}
M.M.~Wilde, A.~Winter and D.~Yang, \emph{{Strong Converse for the Classical
  Capacity of Entanglement-Breaking and Hadamard Channels via a Sandwiched
  Renyi Relative Entropy}},
  \href{https://doi.org/10.1007/s00220-014-2122-x}{\emph{Commun. Math. Phys.}
  {\bfseries 331} (2014) 593}
  [\href{https://arxiv.org/abs/1306.1586}{{\ttfamily 1306.1586}}].

\bibitem{Berta:2016vnw}
M.~Berta, V.B.~Scholz and M.~Tomamichel, \emph{{R\'enyi Divergences as Weighted
  Non-commutative Vector-Valued $L_p$ -Spaces}},
  \href{https://doi.org/10.1007/s00023-018-0670-x}{\emph{Annales Henri
  Poincare} {\bfseries 19} (2018) 1843}
  [\href{https://arxiv.org/abs/1608.05317}{{\ttfamily 1608.05317}}].

\bibitem{Jencova:2016tqz}
A.~Jen\v{c}ov\'a, \emph{{R{\'e}nyi relative entropies and noncommutative
  $L_p$-spaces}},
  \href{https://doi.org/10.1007/s00023-018-0683-5}{\emph{Annales Henri
  Poincare} {\bfseries 19} (2018) 2513}
  [\href{https://arxiv.org/abs/1609.08462}{{\ttfamily 1609.08462}}].

\bibitem{Jencova:2017txf}
A.~Jen\v{c}ov\'a, \emph{{R\'enyi relative entropies and noncommutative
  $L_p$-spaces II}},
  \href{https://doi.org/10.1007/s00023-021-01074-9}{\emph{Annales Henri
  Poincare} {\bfseries 22} (2021) 3235}
  [\href{https://arxiv.org/abs/1707.00047}{{\ttfamily 1707.00047}}].

\bibitem{Lashkari:2018nsl}
N.~Lashkari, \emph{{Constraining Quantum Fields using Modular Theory}},
  \href{https://doi.org/10.1007/JHEP01(2019)059}{\emph{JHEP} {\bfseries 01}
  (2019) 059} [\href{https://arxiv.org/abs/1810.09306}{{\ttfamily
  1810.09306}}].

\bibitem{Moosa:2020jwt}
M.~Moosa, P.~Rath and V.P.~Su, \emph{{A R\'enyi quantum null energy condition:
  proof for free field theories}},
  \href{https://doi.org/10.1007/JHEP01(2021)064}{\emph{JHEP} {\bfseries 01}
  (2021) 064} [\href{https://arxiv.org/abs/2007.15025}{{\ttfamily
  2007.15025}}].

\bibitem{Roy:2022yzm}
P.~Roy, \emph{{Proof of the R\'enyi quantum null energy condition for free
  fermions}}, \href{https://doi.org/10.1103/PhysRevD.108.045010}{\emph{Phys.
  Rev. D} {\bfseries 108} (2023) 045010}
  [\href{https://arxiv.org/abs/2212.02331}{{\ttfamily 2212.02331}}].

\bibitem{Casini:2017roe}
H.~Casini, E.~Teste and G.~Torroba, \emph{{Modular Hamiltonians on the null
  plane and the Markov property of the vacuum state}},
  \href{https://doi.org/10.1088/1751-8121/aa7eaa}{\emph{J. Phys. A} {\bfseries
  50} (2017) 364001} [\href{https://arxiv.org/abs/1703.10656}{{\ttfamily
  1703.10656}}].

\bibitem{Terp:1981lp}
M.~Terp, \emph{{$L^p$} spaces associated with von Neumann algebras},
  Math.~Inst.~Copenhagen University, Copenhagen (1981).

\bibitem{Hiai:2021qfv}
F.~Hiai, \emph{Quantum f-Divergences in von Neumann Algebras: Reversibility of
  Quantum Operations}, Mathematical Physics Studies, Springer, Singapore
  (2021),
  \href{https://doi.org/10.1007/978-981-33-4199-9}{10.1007/978-981-33-4199-9}.

\bibitem{Hiai:2021cks}
F.~Hiai, \emph{Lectures on Selected Topics in von Neumann Algebras}, EMS Press,
  Berlin (2021), \href{https://doi.org/10.4171/elm/32}{10.4171/elm/32}.

\bibitem{KOSAKI198429}
H.~Kosaki, \emph{Applications of the complex interpolation method to a von
  neumann algebra: Non-commutative lp-spaces},
  \href{https://doi.org/https://doi.org/10.1016/0022-1236(84)90025-9}{\emph{Journal
  of Functional Analysis} {\bfseries 56} (1984) 29}.

\bibitem{Kato:2023aro}
S.~Kato and Y.~Ueda, \emph{A remark on non-commutative {$L^p$}-spaces},
  \href{https://doi.org/10.4064/sm230724-11-10}{\emph{Studia Mathematica}
  {\bfseries 275} (2023) 235}
  [\href{https://arxiv.org/abs/2307.01790}{{\ttfamily 2307.01790}}].

\bibitem{Haagerup:2008arm}
U.~Haagerup, M.~Junge and Q.~Xu, \emph{A reduction method for noncommutative
  {$L_p$}-spaces and applications},
  \href{https://doi.org/10.1090/S0002-9947-09-04935-6}{\emph{Transactions of
  the American Mathematical Society} {\bfseries 362} (2008) 2125}
  [\href{https://arxiv.org/abs/0806.3635}{{\ttfamily 0806.3635}}].

\bibitem{Pisier:2003dsk}
G.~Pisier and Q.~Xu, \emph{Non-commutative lp-spaces},  in \emph{Handbook of
  the Geometry of Banach Spaces}, 2003,
  \href{https://doi.org/10.1016/S1874-5849(03)80041-4}{DOI}.

\bibitem{phillips1955adjoint}
R.S.~Phillips, \emph{The adjoint semigroup}, {\emph{Pacific Journal of
  Mathematics} {\bfseries 5} (1955) 269}.

\bibitem{engel.nagel:00:one-parameter}
K.-J.~Engel and R.~Nagel, \emph{One-parameter semigroups for linear evolution
  equations}, vol.~194 of \emph{Graduate Texts in Mathematics},
  Springer-Verlag, New York, New York (2000).

\bibitem{arendt.batty.ea:11:vector-valued}
W.~Arendt, C.J.K.~Batty, M.~Hieber and F.~Neubrander, \emph{Vector-valued
  {L}aplace Transforms and {C}auchy Problems}, vol.~96 of \emph{Monographs in
  Mathematics}, Birkh{\"a}user, Basel, 2nd~ed. (2011).

\bibitem{Calabrese:2004eu}
P.~Calabrese and J.L.~Cardy, \emph{{Entanglement entropy and quantum field
  theory}}, \href{https://doi.org/10.1088/1742-5468/2004/06/P06002}{\emph{J.
  Stat. Mech.} {\bfseries 0406} (2004) P06002}
  [\href{https://arxiv.org/abs/hep-th/0405152}{{\ttfamily hep-th/0405152}}].

\bibitem{olkiewicz1999hypercontractivity}
R.~Olkiewicz and B.~Zegarlinski, \emph{Hypercontractivity in noncommutative
  {$L_p$} spaces}, \href{https://doi.org/10.1006/jfan.1998.3342}{\emph{Journal
  of Functional Analysis} {\bfseries 161} (1999) 246}.

\bibitem{junge2015noncommutative}
M.~Junge and Q.~Zeng, \emph{Noncommutative martingale deviation and
  {Poincar\'e} type inequalities with applications},
  \href{https://doi.org/10.1007/s00440-014-0552-1}{\emph{Probability Theory and
  Related Fields} {\bfseries 161} (2015) 449}.

\bibitem{carlen2017gradient}
E.A.~Carlen and J.~Maas, \emph{Gradient flow and entropy inequalities for
  quantum {M}arkov semigroups with detailed balance},
  \href{https://doi.org/10.1016/j.jfa.2017.05.003}{\emph{Journal of Functional
  Analysis} {\bfseries 273} (2017) 1810}.

\bibitem{Spohn:1978pfz}
H.~Spohn, \emph{{Entropy production for quantum dynamical semigroups}},
  \href{https://doi.org/10.1063/1.523789}{\emph{J. Math. Phys.} {\bfseries 19}
  (1978) 1227}.

\bibitem{Wirth:2025goz}
M.~Wirth, \emph{{Exponential Relative Entropy Decay Along Quantum Markov
  Semigroups}},  \href{https://arxiv.org/abs/2505.07549}{{\ttfamily
  2505.07549}}.

\bibitem{Wilde:2017okz}
M.M.~Wilde, \emph{{Optimized quantum f-divergences and data processing}},
  \href{https://doi.org/10.1088/1751-8121/aad5a1}{\emph{J. Phys. A} {\bfseries
  51} (2018) 374002} [\href{https://arxiv.org/abs/1710.10252}{{\ttfamily
  1710.10252}}].

\bibitem{Audenaert:2015npv}
K.M.R.~Audenaert and N.~Datta, \emph{{alpha-z-relative Renyi entropies}},
  \href{https://doi.org/10.1063/1.4906367}{\emph{J. Math. Phys.} {\bfseries 56}
  (2015) 022202} [\href{https://arxiv.org/abs/1310.7178}{{\ttfamily
  1310.7178}}].

\bibitem{Kato:2023hlj}
S.~Kato, \emph{{On {\ensuremath{\alpha}}-z-R{\'e}nyi divergence in the von
  Neumann algebra setting}}, \href{https://doi.org/10.1063/5.0186552}{\emph{J.
  Math. Phys.} {\bfseries 65} (2024) 042202}
  [\href{https://arxiv.org/abs/2311.01748}{{\ttfamily 2311.01748}}].

\bibitem{Hiai:2024qve}
F.~Hiai and A.~Jen{\v{c}}ov{\'a}, \emph{{$\alpha $-z-R{\'e}nyi Divergences in
  von Neumann Algebras: Data Processing Inequality, Reversibility, and
  Monotonicity Properties in $\alpha ,z$}},
  \href{https://doi.org/10.1007/s00220-024-05124-1}{\emph{Commun. Math. Phys.}
  {\bfseries 405} (2024) 271}
  [\href{https://arxiv.org/abs/2404.07617}{{\ttfamily 2404.07617}}].

\bibitem{Chandrasekaran:2026pnc}
V.~Chandrasekaran and {\'E}.{\'E}.~Flanagan, \emph{{Subregion algebras in
  classical and quantum gravity}},
  \href{https://arxiv.org/abs/2601.07915}{{\ttfamily 2601.07915}}.

\bibitem{Chandrasekaran:2022eqq}
V.~Chandrasekaran, G.~Penington and E.~Witten, \emph{{Large N algebras and
  generalized entropy}},
  \href{https://doi.org/10.1007/JHEP04(2023)009}{\emph{JHEP} {\bfseries 04}
  (2023) 009} [\href{https://arxiv.org/abs/2209.10454}{{\ttfamily
  2209.10454}}].

\bibitem{hille1981functional}
E.~Hille and R.~Phillips, \emph{Functional Analysis and Semi-groups}, American
  Mathematical Society: Colloquium publications, AMS, New York City (1981).

\bibitem{Junge:2002osf}
M.~Junge, \emph{Doob's inequality for non-commutative martingales},
  \href{https://doi.org/10.1515/crll.2002.061}{\emph{J. Reine Angew. Math.}
  (2002) 149} [\href{https://arxiv.org/abs/math/0206062}{{\ttfamily
  math/0206062}}].

\bibitem{Haagerup1975}
U.~Haagerup, \emph{The standard form of von neumann algebras},
  \href{https://doi.org/10.7146/math.scand.a-11606}{\emph{Mathematica
  Scandinavica} {\bfseries 37} (1975) 271}.

\end{thebibliography}\endgroup


\end{document}
